% Les formes de violences — ECM 1re Tech — Cours signé OnBuch+
% Programme officiel MINESEC. Compiler avec : tectonic N04-formes-violences.tex
\newcommand{\DOCMATIERE}{ECM}
\newcommand{\DOCNIVEAU}{1re Tech}
\newcommand{\DOCMODULE}{ECM – classe de Première technique}
\newcommand{\DOCLECON}{N04}
\newcommand{\DOCTITRE}{Les formes de violences}
% OnBuch+ — préambule « cours signés » ENRICHI (compilé avec tectonic / XeLaTeX)
% Système visuel « Pop » d'OnBuch+ : fond crème (jamais blanc), bordures encre
% épaisses, ombres « dures » décalées, titres Archivo Black en capitales.
% Version MATHÉMATIQUES : graphes (pgfplots/tikz), boîtes pédagogiques
% (définition, propriété, méthode, attention, le savais-tu, exercice résolu,
% exercice, TP, bilan), page de garde et signature OnBuch+ ancrée en bas.
%
% Macros attendues dans main.tex AVANT \input{preamble} :
%   \DOCMATIERE  \DOCNIVEAU  \DOCMODULE  \DOCLECON (numéro)  \DOCTITRE
\documentclass[11pt]{article}

\usepackage{fontspec}
\usepackage[french]{babel}
\usepackage[a4paper,margin=2cm,top=3.4cm,bottom=2.8cm,headheight=1.4cm,footskip=1.3cm]{geometry}
\usepackage{amsmath,amssymb,amsfonts}
\usepackage{mathtools} % \xmapsto, \coloneqq... souvent utilisés par les modèles
\usepackage{cancel} % simplifications barrées (bilans de forces)
% \abs{z} (module, valeur absolue) et \norm{u} (norme), utilisés par les modèles
\providecommand{\abs}{}\let\abs\relax\DeclarePairedDelimiter{\abs}{\lvert}{\rvert}
\providecommand{\norm}{}\let\norm\relax\DeclarePairedDelimiter{\norm}{\lVert}{\rVert}
\usepackage[table]{xcolor} % [table] : fournit aussi \rowcolors (alternance de lignes)
\usepackage{pagecolor}
\usepackage{tikz}
\usetikzlibrary{arrows.meta,positioning,calc,shapes.geometric,shapes.misc,shapes.arrows,decorations.pathmorphing,decorations.pathreplacing,decorations.markings,patterns,shadows,mindmap,trees,fit,backgrounds,matrix,intersections,angles,quotes}
\usepackage{pgfplots}
\usepackage[version=4]{mhchem} % équations nucléaires \ce{^{235}_{92}U}
\usepackage[european,siunitx]{circuitikz} % schémas électriques
\pgfplotsset{compat=1.18}
\usepgfplotslibrary{fillbetween,groupplots} % aires entre courbes (calcul intégral)
\usepackage{enumitem}
\usepackage{fancyhdr}
% [most] charge toutes les bibliothèques tcolorbox (listings, minted, raster,
% documentation...) et gonfle inutilement la mémoire TeX sur un document long
% (~300 boîtes) : on ne charge que ce qui est réellement utilisé.
\usepackage[skins,breakable]{tcolorbox}
\usepackage{graphicx}
\usepackage{colortbl}
\usepackage{booktabs}
\usepackage{tabularx}
\usepackage{multirow}
\usepackage{diagbox} % case d'en-tête diagonale des tableaux à double entrée
% Colonnes extensibles centrée / gauche / droite (C, L, R), souvent utilisées par les modèles
\newcolumntype{C}{>{\centering\arraybackslash}X}
\newcolumntype{L}{>{\raggedright\arraybackslash}X}
\newcolumntype{R}{>{\raggedleft\arraybackslash}X}
\usepackage{array}
\usepackage{multicol}
\usepackage{siunitx}
% parse-numbers=false : accepte aussi \SI{2,5 \times 10^{-3}}{...}
\sisetup{locale=FR,output-decimal-marker={,},group-minimum-digits=4,parse-numbers=false}
\DeclareSIUnit\ohm{\ensuremath{\symup{\Omega}}} % U+2126 absent de la police
\DeclareSIUnit\division{div} % réglages d'oscilloscope (ms/div, V/div)
\DeclareSIUnit\tr{tr} % tours (vitesse de rotation, ex. \SI{1500}{\tr\per\minute})
\DeclareSIUnit\molar{M} % concentration molaire (ex. \SI{3}{\molar}, \SI{1}{\micro\molar})
\DeclareSIUnit\an{an} % année (le modèle confond parfois avec \year, un compteur TeX) — ex. \SI{5}{\kilo\meter\per\an}
\DeclareSIUnit\FCFA{FCFA} % franc CFA (ex. \SI{500}{\FCFA\per\kilo\gram})
\DeclareSIUnit\degreeBrix{\textdegree Brix} % teneur en sucre d'un sirop/jus (réfractométrie)
\DeclareSIUnit\month{mois} % le modèle utilise parfois \month, un compteur TeX, comme unité de durée
\DeclareSIUnit\microliter{\micro\liter} % le modèle écrit parfois \microliter en un seul token
\DeclareSIUnit\million{millions} % dénombrement (ex. \SI{15}{\million\per\mL} spermatozoïdes)
\usepackage{qrcode}
% \degree : commande fréquemment utilisée par les modèles pour un angle ou une
% température en dehors de \SI{}{} (non fournie par les paquets chargés).
\providecommand{\degree}{^{\circ}}
% \qed (fin de démonstration), utilisé par les modèles sans amsthm
\providecommand{\qed}{\hfill$\square$}
% \barre : double barre des valeurs interdites dans un tableau de variations (tkz-tab)
\providecommand{\barre}{\ensuremath{\|}}
% environnement proof (amsthm non chargé) utilisé par les modèles
\ifdefined\proof\else\newenvironment{proof}[1][Démonstration]{\par\noindent\textit{#1.}\ }{\qed\par}\fi
\usepackage{lastpage}
\usepackage{hyperref}
\hypersetup{hidelinks}

% ============================================================
% Palette « Pop » OnBuch+ — mêmes hex que la classe P (plus_ui.dart)
% ============================================================
\definecolor{poporange}{HTML}{F59321}
\definecolor{popdark}{HTML}{E07A0C}
\definecolor{popcream}{HTML}{FFF6E8}
\definecolor{popink}{HTML}{1C1714}
\definecolor{poppurple}{HTML}{7A5AE0}
\definecolor{popgreen}{HTML}{2E9E6A}
\definecolor{poppink}{HTML}{E0457B}
\definecolor{popblue}{HTML}{2F7DE1}
\definecolor{popgold}{HTML}{C98A12}
\definecolor{popmuted}{HTML}{8A7F76}
\definecolor{popline}{HTML}{EADFD2}
\definecolor{popgray}{HTML}{8A7F76}
\colorlet{popgrey}{popgray}
% Alias tolérants (le modèle nomme parfois les couleurs différemment)
\colorlet{popwhite}{white}
\colorlet{popblack}{popink}
\colorlet{popred}{poppink}
\colorlet{popyellow}{popgold}
% Teintes claires (fonds de tableaux/encadrés) — le modèle nomme parfois
% les couleurs avec un suffixe L (ex. popblueL) sans les avoir définies.
\colorlet{poporangeL}{poporange!20}
\colorlet{popdarkL}{popdark!20}
\colorlet{poppurpleL}{poppurple!20}
\colorlet{popgreenL}{popgreen!20}
\colorlet{poppinkL}{poppink!20}
\colorlet{popblueL}{popblue!20}
\colorlet{popgoldL}{popgold!20}
\colorlet{popredL}{popred!20}
\colorlet{popgrayL}{popgray!20}
% Teintes claires pour fonds de tableaux / zones
\colorlet{popblueL}{popblue!10!white}
\colorlet{poporangeL}{poporange!14!white}
\colorlet{popgreenL}{popgreen!12!white}
\colorlet{poppurpleL}{poppurple!12!white}
\colorlet{poppinkL}{poppink!12!white}
\colorlet{popgoldL}{popgold!14!white}

\pagecolor{popcream}

% ============================================================
% Typographie
% ============================================================
\defaultfontfeatures{Ligatures=TeX}
\setmainfont{PlusJakartaSans-Medium.ttf}[Path=../../../fonts/,BoldFont=PlusJakartaSans-Bold.ttf]
% Maths en sans-serif (Fira Math) avec chiffres et lettres droites de Plus
% Jakarta, pour que les formules se fondent dans le texte.
\usepackage[math-style=ISO,bold-style=ISO]{unicode-math}
\setmathfont{FiraMath-Regular.otf}[Path=../../../fonts/]
\setmathfont{PlusJakartaSans-Medium.ttf}[Path=../../../fonts/,range={up/num,up/{Latin,latin},\mathup}]
\usepackage{icomma} % virgule décimale sans espace en mode math
% Caractères mathématiques Unicode absents des polices de texte (Plus Jakarta,
% Archivo Black) : rendus par leur équivalent mathématique, en texte comme en
% formule — sinon ils s'affichent en carré vide (ex. « Arithmétique dans ℤ »).
\usepackage{newunicodechar}
\newunicodechar{ℝ}{\ensuremath{\mathbb{R}}}
\newunicodechar{ℤ}{\ensuremath{\mathbb{Z}}}
\newunicodechar{ℂ}{\ensuremath{\mathbb{C}}}
\newunicodechar{ℕ}{\ensuremath{\mathbb{N}}}
\newunicodechar{ℚ}{\ensuremath{\mathbb{Q}}}
\newunicodechar{𝔻}{\ensuremath{\mathbb{D}}}
\newunicodechar{α}{\ensuremath{\alpha}}
\newunicodechar{β}{\ensuremath{\beta}}
\newunicodechar{γ}{\ensuremath{\gamma}}
\newunicodechar{δ}{\ensuremath{\delta}}
\newunicodechar{ε}{\ensuremath{\varepsilon}}
\newunicodechar{θ}{\ensuremath{\theta}}
\newunicodechar{λ}{\ensuremath{\lambda}}
\newunicodechar{μ}{\ensuremath{\mu}}
\newunicodechar{σ}{\ensuremath{\sigma}}
\newunicodechar{τ}{\ensuremath{\tau}}
\newunicodechar{φ}{\ensuremath{\varphi}}
\newunicodechar{ψ}{\ensuremath{\psi}}
\newunicodechar{ω}{\ensuremath{\omega}}
\newunicodechar{Σ}{\ensuremath{\Sigma}}
% majuscules produites par \MakeUppercase dans les titres (α-aminés -> Α)
\newunicodechar{Α}{\ensuremath{\alpha}}
\newunicodechar{Β}{\ensuremath{\beta}}
\newunicodechar{Γ}{\ensuremath{\Gamma}}
\newunicodechar{Δ}{\ensuremath{\Delta}}
\newunicodechar{Ε}{\ensuremath{\varepsilon}}
\newunicodechar{Θ}{\ensuremath{\Theta}}
\newunicodechar{Λ}{\ensuremath{\Lambda}}
\newunicodechar{Μ}{\ensuremath{\mu}}
\newunicodechar{Τ}{\ensuremath{\tau}}
\newunicodechar{Φ}{\ensuremath{\Phi}}
\newunicodechar{Ψ}{\ensuremath{\Psi}}
\newunicodechar{Ω}{\ensuremath{\Omega}}
\newunicodechar{↦}{\ensuremath{\mapsto}}
\newunicodechar{∈}{\ensuremath{\in}}
\newunicodechar{∉}{\ensuremath{\notin}}
\newunicodechar{∧}{\ensuremath{\wedge}}
\newunicodechar{⇔}{\ensuremath{\Leftrightarrow}}
\newunicodechar{⟺}{\ensuremath{\Longleftrightarrow}}
\newunicodechar{⇒}{\ensuremath{\Rightarrow}}
\newunicodechar{⟹}{\ensuremath{\Longrightarrow}}
\newunicodechar{∘}{\ensuremath{\circ}}
\newunicodechar{∪}{\ensuremath{\cup}}
\newunicodechar{∩}{\ensuremath{\cap}}
\newunicodechar{≡}{\ensuremath{\equiv}}
\newunicodechar{⊂}{\ensuremath{\subset}}
\newunicodechar{⊥}{\ensuremath{\perp}}
\newunicodechar{∃}{\ensuremath{\exists}}
\newunicodechar{∀}{\ensuremath{\forall}}
\newunicodechar{∛}{\ensuremath{\sqrt[3]{\,}}}
\newunicodechar{∜}{\ensuremath{\sqrt[4]{\,}}}
\newunicodechar{⃗}{\ensuremath{{}^{\to}}}
\newunicodechar{ⁿ}{\ifmmode^{n}\else\textsuperscript{\ensuremath{n}}\fi}
\newunicodechar{ˣ}{\ifmmode^{x}\else\textsuperscript{\ensuremath{x}}\fi}
\newunicodechar{ᵇ}{\ifmmode^{b}\else\textsuperscript{\ensuremath{b}}\fi}
\newunicodechar{ʳ}{\ifmmode^{r}\else\textsuperscript{\ensuremath{r}}\fi}
\newunicodechar{ᵘ}{\ifmmode^{u}\else\textsuperscript{\ensuremath{u}}\fi}
\newunicodechar{ʸ}{\ifmmode^{y}\else\textsuperscript{\ensuremath{y}}\fi}
\newunicodechar{ᵃ}{\ifmmode^{a}\else\textsuperscript{\ensuremath{a}}\fi}
\newunicodechar{ᵉ}{\ifmmode^{e}\else\textsuperscript{\ensuremath{e}}\fi}
\newunicodechar{ᵏ}{\ifmmode^{k}\else\textsuperscript{\ensuremath{k}}\fi}
\newunicodechar{ᵖ}{\ifmmode^{p}\else\textsuperscript{\ensuremath{p}}\fi}
\newunicodechar{ᴺ}{\ifmmode^{N}\else\textsuperscript{\ensuremath{N}}\fi}
\newunicodechar{ᶜ}{\ifmmode^{c}\else\textsuperscript{\ensuremath{c}}\fi}
\newunicodechar{ʰ}{\ifmmode^{h}\else\textsuperscript{\ensuremath{h}}\fi}
\newunicodechar{ᵠ}{\ifmmode^{q}\else\textsuperscript{\ensuremath{q}}\fi}
\newunicodechar{ₙ}{\ifmmode_{n}\else\textsubscript{\ensuremath{n}}\fi}
\newunicodechar{ₐ}{\ifmmode_{a}\else\textsubscript{\ensuremath{a}}\fi}
\newunicodechar{ᵢ}{\ifmmode_{i}\else\textsubscript{\ensuremath{i}}\fi}
\newunicodechar{ᵣ}{\ifmmode_{r}\else\textsubscript{\ensuremath{r}}\fi}
\newunicodechar{ₖ}{\ifmmode_{k}\else\textsubscript{\ensuremath{k}}\fi}
\newunicodechar{ᵦ}{\ifmmode_{\beta}\else\textsubscript{\ensuremath{\beta}}\fi}
\newunicodechar{ₓ}{\ifmmode_{x}\else\textsubscript{\ensuremath{x}}\fi}
\newfontfamily\popdisplay{ArchivoBlack-Regular.ttf}[Path=../../../fonts/]
\newfontfamily\poplogo{SpaceGrotesk-Bold.ttf}[Path=../../../fonts/]
\newfontfamily\popsemi{PlusJakartaSans-SemiBold.ttf}[Path=../../../fonts/]
\newfontfamily\popxbold{PlusJakartaSans-ExtraBold.ttf}[Path=../../../fonts/]
\newfontfamily\popxboldls{PlusJakartaSans-ExtraBold.ttf}[Path=../../../fonts/,LetterSpace=3.0]

\setlist[enumerate]{leftmargin=1.7em,itemsep=5pt,topsep=4pt}
\setlist[itemize]{leftmargin=1.7em,itemsep=4pt,topsep=4pt,label={\color{poporange}\textbullet}}
\setlength{\parindent}{0pt}
\setlength{\parskip}{5pt}
\renewcommand{\arraystretch}{1.25}

% pgfplots : style Pop par défaut
\pgfplotsset{
  every axis/.append style={axis line style={popink,line width=1pt},tick style={popink},
    grid style={popline,line width=0.5pt},label style={font=\small},tick label style={font=\footnotesize}},
  popcurve/.style={poporange,line width=1.8pt,no markers,smooth},
}
% Même style disponible dans un \draw TikZ simple (hors pgfplots)
\tikzset{popcurve/.style={poporange,line width=1.8pt}}
\tikzset{popgrid/.style={popline,very thin}}
\colorlet{popcurve}{poporange} % le nom du style est parfois utilisé comme couleur

% ============================================================
% Pill / squiggle
% ============================================================
\newcommand{\poppill}[3]{%
  \tikz[baseline=-0.55ex]\node[draw=popink,line width=1.2pt,rounded corners=2.6pt,%
    fill=#2,inner xsep=6.5pt,inner ysep=3pt]{%
    \popxboldls\fontsize{7}{7}\selectfont\color{#3}\MakeUppercase{#1}};%
}
\newcommand{\popsquiggle}[1]{%
  \tikz[baseline=-0.4ex]{
    \draw[#1,line width=2.6pt,line cap=round]
      plot [smooth] coordinates {(0,0.09) (0.34,0.01) (0.68,0.21) (1.02,0.01) (1.36,0.10)};
  }%
}
% Mot-clé surligné (marqueur orange sous le texte)
\newcommand{\cle}[1]{{\popxbold\color{popdark}#1}}

% ============================================================
% En-tête / pied
% ============================================================
\pagestyle{fancy}
\fancyhf{}
\renewcommand{\headrulewidth}{0pt}
\renewcommand{\footrulewidth}{0pt}
\lhead{\raisebox{-0.15ex}{\poplogo\fontsize{13}{13}\selectfont\color{popink}OnBuch\color{poporange}+}}
\rhead{\poppill{\DOCMATIERE\ \textbullet\ \DOCNIVEAU\ \textbullet\ Leçon \DOCLECON}{white}{popink}}
\lfoot{\popxbold\fontsize{7.6}{7.6}\selectfont\color{popmuted}\MakeUppercase{Cours signé OnBuch+}\ \textbullet\ {\popsemi\DOCTITRE}}
\rfoot{\tikz[baseline=-0.6ex]\node[rounded corners=8.5pt,fill=popink,minimum height=17pt,minimum width=17pt,inner xsep=5pt,inner sep=0]{\popxbold\fontsize{8}{8}\selectfont\color{popcream}\thepage\,/\,\pageref*{LastPage}};}

% ============================================================
% Titres
% ============================================================
\newcommand{\chaptitre}[2]{%
  \begin{tcolorbox}[enhanced,colback=poporange,colframe=popink,boxrule=2.6pt,arc=11pt,
      drop shadow=popink,left=15pt,right=15pt,top=12pt,bottom=11pt,before skip=3pt,after skip=18pt]
    \poppill{#2}{popink}{popcream}\par\vspace{9pt}
    {\raggedright\hyphenpenalty=10000\exhyphenpenalty=10000\popdisplay\fontsize{22}{24}\selectfont\color{popcream}\MakeUppercase{#1}\par}
  \end{tcolorbox}
}
% Section numérotée : pastille orange avec le numéro + titre Archivo
\newcounter{popsec}
\newcommand{\coursec}[1]{%
  \stepcounter{popsec}\setcounter{popsub}{0}%
  \par\vspace{16pt}\noindent
  \begin{minipage}{\linewidth}
  \tikz[baseline=-0.7ex]\node[draw=popink,line width=1.6pt,fill=poporange,rounded corners=4pt,minimum size=22pt,inner sep=0]{\popdisplay\fontsize{12}{12}\selectfont\color{popcream}\thepopsec};%
  \hspace{8pt}{\popdisplay\fontsize{15}{17}\selectfont\color{popink}\MakeUppercase{#1}}\par
  \vspace{4pt}\noindent{\color{poporange}\rule{36pt}{3.6pt}}
  \end{minipage}\par\nopagebreak\vspace{8pt}
}
\newcounter{popsub}
\newcommand{\courssub}[1]{%
  \stepcounter{popsub}%
  \par\vspace{9pt}\noindent{\popxbold\fontsize{11.5}{13}\selectfont\color{popink}{\color{poporange}\thepopsec.\thepopsub}\hspace{6pt}#1}\par\nopagebreak\vspace{4pt}
}

% ============================================================
% Boîtes pédagogiques — même recette : fond blanc, bordure épaisse colorée,
% ombre dure encre, badge plein. Titre optionnel [#1].
% ============================================================
\tcbset{popbase/.style={breakable,enhanced,colback=white,boxrule=2.3pt,arc=9pt,
  shadow={3pt}{-3pt}{0pt}{popink},left=14pt,right=14pt,top=11pt,bottom=11pt,before skip=11pt,after skip=15pt,
  fontupper=\popsemi\fontsize{10.6}{14.5}\selectfont\color{popink},
  fontlower=\popsemi\fontsize{10.6}{14.5}\selectfont\color{popink}}}
% Coche dessinée (le glyphe ✓ n'existe pas dans les polices du document)
\newcommand{\popcheck}[1]{\tikz[baseline=-0.5ex]\draw[#1,line width=1.6pt,line cap=round,line join=round](-0.13,0.0)--(-0.04,-0.09)--(0.13,0.1);}
\providecommand{\coche}{\popcheck{popgreen}}
\providecommand{\resultat}[1]{\par\smallskip\noindent\fcolorbox{popgreen}{popgreenL}{\parbox{\dimexpr\linewidth-2\fboxsep-2\fboxrule\relax}{\textbf{Résultat.} #1}}\par\smallskip}
\newcommand{\popboxhead}[3]{\poppill{#1}{#2}{white}\ifx\relax#3\relax\else\hspace{6pt}{\popxbold\fontsize{10.6}{12}\selectfont\color{popink}#3}\fi\par\vspace{7pt}}

\newtcolorbox{aretenir}[1][]{popbase,colframe=popgreen,before upper={\popboxhead{À retenir}{popgreen}{#1}}}
\newtcolorbox{exemplebox}[1][]{popbase,colframe=poporange,before upper={\popboxhead{Exemple}{poporange}{#1}}}
\newtcolorbox{definition}[1][]{popbase,colframe=popblue,before upper={\popboxhead{Définition}{popblue}{#1}}}
\newtcolorbox{propriete}[1][]{popbase,colframe=poppurple,before upper={\popboxhead{Propriété}{poppurple}{#1}}}
\newtcolorbox{methode}[1][]{popbase,colframe=popgold,before upper={\popboxhead{Méthode}{popgold}{#1}}}
\newtcolorbox{attention}[1][]{popbase,colframe=poppink,before upper={\popboxhead{Attention}{poppink}{#1}}}
\newtcolorbox{experience}[1][]{popbase,colframe=popblue,colback=popblueL,before upper={\popboxhead{Expérience}{popblue}{#1}}}
% « Le savais-tu ? » : carte violette pleine (contexte, culture, Cameroun)
\newtcolorbox{savaistu}[1][]{popbase,colback=poppurple,colframe=popink,
  fontupper=\popsemi\fontsize{10.4}{14.2}\selectfont\color{white},
  before upper={\poppill{Le savais-tu\,?}{white}{poppurple}\ifx\relax#1\relax\else\hspace{6pt}{\popxbold\color{white}#1}\fi\par\vspace{7pt}}}
% Exercice résolu : énoncé en haut, \tcblower puis la solution sur fond crème
\newcounter{popexo}\newcounter{popexr}
\newtcolorbox{exoresolu}[1][]{popbase,colframe=poporange,colbacklower=popcream,
  segmentation style={popink,line width=1.2pt,dashed},
  before upper={\stepcounter{popexr}\popboxhead{Exercice résolu \thepopexr}{poporange}{#1}},
  before lower={\poppill{Solution}{popink}{popcream}\par\vspace{6pt}}}
% Exercice (non résolu en place) — la correction est dans la partie Corrigés
\newtcolorbox{exercice}[1][]{popbase,colframe=popink,
  before upper={\stepcounter{popexo}\popboxhead{Exercice \thepopexo}{popink}{#1}}}
% Corrigé
\newtcolorbox{corrige}[1][]{popbase,colframe=popgreen,colback=popgreenL,
  before upper={\popboxhead{Corrigé}{popgreen}{#1}}}
% Objectifs / prérequis / bilan
\newtcolorbox{objectifs}{popbase,colframe=popink,colback=poporangeL,
  before upper={\popboxhead{Objectifs de la leçon}{popink}{}}}
\newtcolorbox{prerequis}{popbase,colframe=popink,colback=popblueL,
  before upper={\popboxhead{Prérequis}{popblue}{}}}
\newtcolorbox{bilan}{popbase,colframe=popink,colback=white,boxrule=2.8pt,
  before upper={\popboxhead{Fiche bilan}{poporange}{L'essentiel à connaître pour l'examen}}}
% Figure encadrée avec légende
\newtcolorbox{popfigure}[1][]{enhanced,breakable=false,colback=white,colframe=popline,boxrule=1.4pt,arc=8pt,
  left=10pt,right=10pt,top=10pt,bottom=8pt,before skip=10pt,after skip=12pt,halign=center,
  lower separated=false,halign lower=center,
  fontlower=\popsemi\fontsize{9}{11.5}\selectfont\color{popmuted},
  #1}
\newcommand{\legende}[1]{\par\vspace{6pt}{\popsemi\fontsize{9}{11.5}\selectfont\color{popmuted}#1}}

% ============================================================
% Signature OnBuch+
% ============================================================
\newcommand{\poplogoinline}[1]{{\poplogo\fontsize{#1}{#1}\selectfont\color{popink}OnBuch\color{poporange}+}}
\newcommand{\signatureonbuch}{%
  \begin{tcolorbox}[enhanced,colback=popink,colframe=popink,boxrule=2.2pt,arc=10pt,
      drop shadow=poporange,left=14pt,right=14pt,top=10pt,bottom=10pt,before skip=0pt,after skip=0pt]
    \tikz[baseline=-0.7ex]\node[circle,fill=popgreen,draw=popcream,line width=1.3pt,minimum size=22pt,inner sep=0]{\popcheck{white}};%
    \hspace{9pt}\begin{minipage}[c]{0.62\linewidth}
      {\popxbold\fontsize{12}{13}\selectfont\color{popcream}Signé\ \textbullet\ {\poplogo OnBuch\color{poporange}+}}\par\vspace{2pt}
      {\raggedright\popsemi\fontsize{8}{10.5}\selectfont\color{popline}\MakeUppercase{Équipe pédagogique OnBuch+}\\\MakeUppercase{Conforme au programme officiel MINESEC}\par}
    \end{minipage}\hfill
    {\poplogo\fontsize{22}{22}\selectfont\color{popcream}OnBuch\color{poporange}+}
  \end{tcolorbox}%
}

% ============================================================
% Page de garde
% #1 titre · #2 matière · #3 niveau · #4 module · #5 n° leçon · #6 durée ·
% #7 description / objectifs (texte ou itemize)
% ============================================================
\newcommand{\pagegarde}[7]{%
  \thispagestyle{empty}
  \newgeometry{a4paper,margin=1.7cm,top=1.6cm,bottom=1.4cm}%
  \begin{tikzpicture}[remember picture,overlay]
    \fill[poporange!18] ([xshift=-3.2cm,yshift=-3.6cm]current page.north east) circle (4.6cm);
    \fill[poppurple!14] ([xshift=2.4cm,yshift=7.6cm]current page.south west) circle (3.8cm);
  \end{tikzpicture}%
  {\poplogo\fontsize{30}{30}\selectfont\color{popink}OnBuch\color{poporange}+}\hfill
  \raisebox{0.6ex}{\poppill{Cours signé \textbullet\ Premium}{popink}{popcream}}\par
  \vspace{1pt}\popsquiggle{poppurple}\par
  \vspace{8pt}
  \begin{tcolorbox}[enhanced,colback=poporange,colframe=popink,boxrule=2.8pt,arc=13pt,
      drop shadow=popink,left=18pt,right=18pt,top=12pt,bottom=13pt,after skip=10pt]
    \poppill{#2 \textbullet\ #3}{popink}{popcream}\hspace{5pt}\poppill{#4}{white}{popink}\par\vspace{8pt}
    {\popdisplay\fontsize{14}{14}\selectfont\color{popink}LEÇON #5}\par\vspace{4pt}
    {\raggedright\hyphenpenalty=10000\exhyphenpenalty=10000\popdisplay\fontsize{25}{27}\selectfont\color{popcream}\MakeUppercase{#1}\par}
    \vspace{8pt}
    {\popxbold\fontsize{9.5}{9.5}\selectfont\color{popink}\MakeUppercase{Durée conseillée : #6}}
  \end{tcolorbox}
  \begin{tcolorbox}[enhanced,colback=white,colframe=popink,boxrule=2.2pt,arc=10pt,
      drop shadow=popink,left=16pt,right=16pt,top=10pt,bottom=10pt,after skip=8pt]
    \poppill{Dans cette leçon}{poppurple}{white}\par\vspace{5pt}
    \popsemi\fontsize{9.3}{12.6}\selectfont\color{popink}#7
  \end{tcolorbox}
  \noindent\begin{minipage}[t]{0.487\linewidth}
    \begin{tcolorbox}[enhanced,colback=popgreen,colframe=popink,boxrule=2.2pt,arc=10pt,
        drop shadow=popink,left=11pt,right=11pt,top=10pt,bottom=10pt,height=3.1cm,valign=center]
      \tcbox[colback=white,colframe=popink,boxrule=1.3pt,arc=5pt,boxsep=0pt,nobeforeafter,
        left=3pt,right=3pt,top=3pt,bottom=3pt,valign=center]{\qrcode[height=1.9cm]{https://play.google.com/store/apps/details?id=com.luvvix.onbuch&hl=fr}}%
      \hspace{8pt}\begin{minipage}[c]{0.5\linewidth}
        {\popdisplay\fontsize{11}{12.5}\selectfont\color{white}\MakeUppercase{Télécharge OnBuch}}\par\vspace{4pt}
        {\raggedright\popsemi\fontsize{8.4}{11}\selectfont\color{white}Gratuit \textbullet\ Google Play\\Tuteur IA, annales, concours, résultats\par}
      \end{minipage}
    \end{tcolorbox}
  \end{minipage}\hfill
  \begin{minipage}[t]{0.487\linewidth}
    \begin{tcolorbox}[enhanced,colback=poppurple,colframe=popink,boxrule=2.2pt,arc=10pt,
        drop shadow=popink,left=11pt,right=11pt,top=10pt,bottom=10pt,height=3.1cm,valign=center]
      \tikz[baseline=-0.7ex]\node[circle,fill=white,draw=popink,line width=1.3pt,minimum size=1.6cm,inner sep=0]{\popdisplay\fontsize{24}{24}\selectfont\color{poppurple}+};%
      \hspace{8pt}\begin{minipage}[c]{0.6\linewidth}
        {\popdisplay\fontsize{11}{12.5}\selectfont\color{white}\MakeUppercase{Passe à OnBuch+}}\par\vspace{4pt}
        {\raggedright\popsemi\fontsize{8.4}{11}\selectfont\color{white}Tous les cours signés, Tuteur IA illimité, fiches PDF — onglet OnBuch+ de l'app\par}
      \end{minipage}
    \end{tcolorbox}
  \end{minipage}
  \vspace{8pt}
  \popcontenu
  \vfill
  \signatureonbuch
  \restoregeometry
  \newpage
}

% Rangée « Ce cours contient » (page de garde)
\newcommand{\poptile}[3]{% #1 couleur #2 gros texte #3 légende
  \begin{tcolorbox}[enhanced,colback=#1,colframe=popink,boxrule=2pt,arc=9pt,shadow={2.5pt}{-2.5pt}{0pt}{popink},
      left=6pt,right=6pt,top=9pt,bottom=9pt,halign=center,valign=center,height=1.75cm,width=\linewidth,nobeforeafter]
    {\popdisplay\fontsize{12.5}{13}\selectfont\color{white}\MakeUppercase{#2}}\par\vspace{3pt}
    {\centering\popsemi\fontsize{7.8}{9.5}\selectfont\color{white}#3\par}
  \end{tcolorbox}}
\newcommand{\popcontenu}{%
  \par\noindent{\popxboldls\fontsize{7.5}{7.5}\selectfont\color{popmuted}CE COURS SIGNÉ CONTIENT}\par\vspace{7pt}
  \noindent\begin{minipage}[t]{0.235\linewidth}\poptile{popblue}{Cours}{complet, illustré, pas à pas}\end{minipage}\hfill
  \begin{minipage}[t]{0.235\linewidth}\poptile{poporange}{Méthodes}{et exercices résolus}\end{minipage}\hfill
  \begin{minipage}[t]{0.235\linewidth}\poptile{poppink}{Exercices}{type examen + corrigés}\end{minipage}\hfill
  \begin{minipage}[t]{0.235\linewidth}\poptile{popgreen}{Fiche bilan}{l'essentiel à retenir}\end{minipage}\par}

% Sommaire
\newtcolorbox{sommaire}{enhanced,colback=white,colframe=popink,boxrule=2.2pt,arc=10pt,
  drop shadow=popink,left=18pt,right=18pt,top=13pt,bottom=13pt,before skip=6pt,after skip=20pt,
  before upper={\poppill{Sommaire}{poppurple}{white}\par\vspace{9pt}\popsemi\fontsize{10.6}{14.5}\selectfont\color{popink}}}

% Signature de fin de cours — poussée tout en bas de la dernière page
\newcommand{\findecours}{%
  \par\vfill\nopagebreak
  \begin{center}\popsquiggle{poporange}\end{center}\vspace{4pt}
  \signatureonbuch
}
\AtBeginDocument{\def\llbracket{\mathopen{[\mkern-3.5mu[}}\def\rrbracket{\mathclose{]\mkern-3.5mu]}}} % intervalles d'entiers (glyphe ⟦ absent de la police)
% Glyphes absents de Fira Math : redéfinis à partir de glyphes disponibles
\AtBeginDocument{%
  \def\setminus{\mathbin{\backslash}}\let\smallsetminus\setminus
  \def\vdots{\mathord{\vbox{\baselineskip3.2pt\lineskiplimit0pt\kern4pt\hbox{.}\hbox{.}\hbox{.}}}}%
  \def\ddots{\mathinner{\mkern1mu\raise6.5pt\hbox{.}\mkern2mu\raise3.6pt\hbox{.}\mkern2mu\raise0.7pt\hbox{.}\mkern1mu}}%
  \def\triangle{\mathord{\scalebox{0.9}{$\symup{\Delta}$}}}%
  \def\nabla{\mathord{\raisebox{\depth}{\scalebox{1}[-1]{$\symup{\Delta}$}}}}%
  \def\checkmark{\popcheck{popdark}}%
  \def\star{\mathbin{\ast}}\def\bigstar{\mathord{\ast}}%
  \def\diamond{\mathbin{\scalebox{0.6}{$\Diamond$}}}%
  \def\bigcup{\mathop{\mathchoice{\raisebox{-0.3em}{\scalebox{1.7}{$\cup$}}}{\scalebox{1.25}{$\cup$}}{\cup}{\cup}}\displaylimits}%
  \def\bigcap{\mathop{\mathchoice{\raisebox{-0.3em}{\scalebox{1.7}{$\cap$}}}{\scalebox{1.25}{$\cap$}}{\cap}{\cap}}\displaylimits}%
  \def\lesseqqgtr{\mathrel{\lessgtr}}\def\bigoplus{\mathop{\oplus}\displaylimits}%
}
\providecommand{\ssi}{si et seulement si} % abréviation fréquente
\AtBeginDocument{\providecommand{\wideparen}{\overparen}} % arc de cercle

%% FIGFIX-BEGIN v4 — correctifs globaux des figures (docs/onbuch-generation/tools/figfix)
\usepackage{adjustbox}\usepackage{etoolbox}
% toute figure TikZ/pgfplots plus large que la ligne est réduite à la largeur disponible
\BeforeBeginEnvironment{tikzpicture}{\begin{adjustbox}{max width=\linewidth}}
\AfterEndEnvironment{tikzpicture}{\end{adjustbox}}
% légendes pgfplots lisibles par-dessus les courbes
\pgfplotsset{every axis/.append style={legend style={fill=white,fill opacity=0.92,text opacity=1,draw=black!15,inner sep=3pt}}}
% étiquettes d'arêtes (syntaxe edge["..."]) avec fond blanc
\tikzset{every edge quotes/.append style={fill=white,inner sep=1.2pt}}
%% FIGFIX-END
\begin{document}

\pagegarde{Les formes de violences}{ECM}{1re Tech}{ECM – classe de Première technique}{N04}{2 séances (fiche de progression harmonisée)}{Cette leçon présente les différentes formes de violences qui menacent la personne humaine et la société : les atteintes à l'intégrité physique, morale et au patrimoine, ainsi que l'extrémisme violent. À partir de documents et de situations concrètes du Cameroun, elle outille l'élève pour reconnaître, analyser et proposer des réponses citoyennes à ces fléaux.
\vspace{4pt}
\begin{multicols}{2}\small
\begin{itemize}
\item À la fin de cette leçon, je sais définir la violence et distinguer ses principales formes
\item identifier et qualifier les atteintes à l'intégrité physique (coups et blessures, meurtre, violences faites aux femmes)
\item identifier et qualifier les atteintes à l'intégrité morale (diffamation, injures, harcèlement, discrimination)
\item identifier et qualifier les atteintes au patrimoine (vol, escroquerie, dégradation, spoliation foncière)
\item expliquer les causes, les manifestations et les conséquences de l'extrémisme violent au Cameroun (Boko Haram, crise anglophone)
\item analyser des documents (photographies, données chiffrées, articles de presse, textes de loi) relatifs aux violences
\end{itemize}\end{multicols}}

\begin{sommaire}
\begin{multicols}{2}
\begin{enumerate}
\item Généralités sur les violences
\item Les atteintes à l'intégrité physique
\item Les atteintes à l'intégrité morale
\item Les atteintes au patrimoine
\item L'extrémisme violent
\item Prévenir et combattre les violences : le rôle du citoyen
\item Activité d'intégration
\item Méthodes et exercices résolus
\item Exercices
\item Corrigés des exercices
\item Fiche bilan
\end{enumerate}
\end{multicols}
\end{sommaire}

\begin{objectifs}
\begin{itemize}
\item À la fin de cette leçon, je sais définir la violence et distinguer ses principales formes
\item identifier et qualifier les atteintes à l'intégrité physique (coups et blessures, meurtre, violences faites aux femmes)
\item identifier et qualifier les atteintes à l'intégrité morale (diffamation, injures, harcèlement, discrimination)
\item identifier et qualifier les atteintes au patrimoine (vol, escroquerie, dégradation, spoliation foncière)
\item expliquer les causes, les manifestations et les conséquences de l'extrémisme violent au Cameroun (Boko Haram, crise anglophone)
\item analyser des documents (photographies, données chiffrées, articles de presse, textes de loi) relatifs aux violences
\item appliquer les notions de la leçon à des situations concrètes de la vie courante et proposer des conduites citoyennes adaptées
\item rédiger et justifier une démarche, une réponse ou une conclusion argumentée sur un cas de violence
\end{itemize}
\end{objectifs}

\begin{prerequis}
\begin{itemize}
\item La dignité de la personne humaine et les droits de l'homme (notions de Seconde et de début d'année)
\item Les valeurs de la citoyenneté : respect, tolérance, solidarité (classes de 4e-3e)
\item Les institutions de la République et le rôle des forces de maintien de l'ordre
\item La notion de conflit et de règlement pacifique des différends
\end{itemize}
\end{prerequis}

\newpage

\coursec{Généralités sur les violences}

\courssub{Situation d'accroche et problématique}

En février 2024, à Yaoundé, une élève de Première d'un lycée de la capitale est agressée à la sortie des cours. Frappée par deux individus, elle est hospitalisée ; son téléphone portable et son argent sont emportés. La même semaine, à Maroua, un jeune de dix-sept ans raconte à ses camarades comment un groupe radical a tenté de le recruter via les réseaux sociaux, en lui promettant de l'argent et un \og avenir glorieux \fg{}. Deux histoires, deux villes, deux réalités différentes : et pourtant, dans les deux cas, des personnes cherchent à imposer leur volonté à autrui par la force, la menace ou la manipulation.

Face à ces situations, les élèves se posent des questions légitimes : qu'est-ce qu'une violence ? Toutes les violences se ressemblent-elles ? Comment l'État et les citoyens peuvent-ils prévenir et réprimer ces atteintes ? C'est à ces questions que répond cette leçon, qui fournit les outils pour \cle{identifier}, \cle{classer} et \cle{combattre} les formes de violences dans la société camerounaise.

\courssub{Définition de la violence et de la victime}

Dans la vie quotidienne, le mot \og violence \fg{} est employé à tort et à travers : on parle d'un match \og violent \fg{}, d'une pluie \og violente \fg{}. En Éducation à la Citoyenneté et à la Morale, le terme a un sens précis qu'il faut maîtriser.

\begin{definition}[Violence et victime]
La \cle{violence} est l'\textbf{usage de la force physique ou psychologique} pour \textbf{contraindre, dominer, blesser ou détruire} une personne ou un bien. Elle se caractérise par trois éléments : une \emph{intention} (vouloir nuire ou contraindre), un \emph{moyen} (force physique, menace, humiliation, manipulation) et un \emph{dommage} (blessure, souffrance, perte matérielle).\par
La \cle{victime} est la personne qui subit directement ou indirectement le dommage causé par la violence. On distingue la victime \emph{directe} (celle qui est frappée, volée, humiliée) et les victimes \emph{indirectes} (la famille, les témoins, la communauté).
\end{definition}

Reprenons la situation d'accroche : l'élève de Yaoundé a subi une violence car il y a eu intention de nuire (dérober les biens), usage de la force (coups portés) et dommage (blessures, perte du téléphone). Le jeune de Maroua, lui, a été confronté à une forme plus insidieuse de violence : la \emph{manipulation psychologique} visant à le recruter dans un groupe radical.

\begin{attention}[Violence ou autorité légitime ?]
Il ne faut \textbf{pas confondre} la violence, qui est \emph{illégale et condamnable}, avec l'\cle{autorité légitime}, qui est l'encadrement exercé dans le respect des règles : un enseignant qui sanctionne un élève selon le règlement intérieur, un parent qui interdit une sortie dangereuse, un policier qui contrôle une identité exercent une autorité légitime, à condition que la sanction soit \emph{proportionnée} et \emph{respectueuse de la dignité}. En revanche, des coups, des insultes, des humiliations publiques ou des sanctions arbitraires constituent des violences, même commises par une personne investie d'autorité.
\end{attention}

\courssub{Les grandes catégories de violences}

Toutes les violences ne se ressemblent pas. Pour les comprendre et les combattre, il faut d'abord savoir les \textbf{classer}. On distingue d'abord trois grandes catégories selon \emph{ce qui est atteint} chez la victime.

\begin{definition}[Les trois catégories d'atteintes]
\begin{itemize}
\item La \cle{violence physique} atteint le \textbf{corps} de la personne : coups et blessures, agressions, sévices, meurtres.
\item La \cle{violence morale} (ou psychologique) atteint la \textbf{dignité et l'équilibre psychique} : insultes, menaces, humiliations, harcèlement, diffamation, manipulation.
\item La \cle{violence patrimoniale} atteint les \textbf{biens} de la personne : vols, escroqueries, dégradations, spoliations de terres ou d'héritages.
\end{itemize}
\end{definition}

\begin{popfigure}
\begin{center}
\begin{tikzpicture}[
  font=\small,
  racine/.style={rectangle, rounded corners, draw=popink, fill=popink!15, thick, text width=3.4cm, align=center},
  branche/.style={rectangle, rounded corners, draw=popblue, fill=popblueL, thick, text width=3.6cm, align=center},
  ex/.style={rectangle, rounded corners, draw=popmuted, fill=popcream, text width=3.6cm, align=center, font=\footnotesize},
  lien/.style={-{Stealth[length=2.5mm]}, thick, popdark}
]
\node[racine, align=center] (V) at (0,0){\textbf{LES VIOLENCES}\\ (atteintes à la personne et aux biens)};
\node[branche, align=center] (P) at (-6,-2.6){\textbf{Violence physique}\\ atteinte au corps};
\node[branche, align=center] (M) at (0,-2.6){\textbf{Violence morale}\\ atteinte à la dignité};
\node[branche, align=center] (Pa) at (6,-2.6){\textbf{Violence patrimoniale}\\ atteinte aux biens};
\node[ex] (Pe) at (-6,-5) {coups et blessures, agression, sévices corporels, meurtre};
\node[ex] (Me) at (0,-5) {insultes, menaces, harcèlement, diffamation, humiliation};
\node[ex] (Pae) at (6,-5) {vol, escroquerie, dégradation, spoliation d'héritage};
\draw[lien] (V) -- (P);
\draw[lien] (V) -- (M);
\draw[lien] (V) -- (Pa);
\draw[lien, popmuted] (P) -- (Pe);
\draw[lien, popmuted] (M) -- (Me);
\draw[lien, popmuted] (Pa) -- (Pae);
\end{tikzpicture}
\end{center}
\legende{Schéma 1 : les trois grandes catégories d'atteintes. (1) Violence physique : le corps est visé. (2) Violence morale : la dignité et le psychisme sont visés. (3) Violence patrimoniale : les biens sont visés. Un même acte peut relever de plusieurs catégories : l'agression de l'élève de Yaoundé est à la fois physique (coups) et patrimoniale (vol du téléphone).}
\end{popfigure}

À cette classification selon \emph{l'objet atteint}, les spécialistes ajoutent une classification selon \emph{la manière dont la violence s'exerce}. C'est un complément utile pour le Baccalauréat technique.

\begin{definition}[Violence directe, structurelle et culturelle (typologie de J. Galtung)]
\begin{itemize}
\item La \cle{violence directe} est visible et identifiable : elle a un auteur et une victime clairs (une agression, un vol, une insulte).
\item La \cle{violence structurelle} est cachée dans l'organisation de la société : ce sont les \textbf{inégalités et exclusions} qui empêchent certaines personnes d'accéder à l'éducation, à la santé, à l'emploi ou à la justice, sans qu'on puisse désigner un seul agresseur.
\item La \cle{violence culturelle} désigne les \textbf{idées, préjugés et discours} qui rendent les autres violences \og acceptables \fg{} : mépris d'un groupe ethnique, discours de haine, idéologies qui justifient la domination d'une catégorie de personnes.
\end{itemize}
\end{definition}

\begin{exemplebox}[Illustrer les trois niveaux de violence]
Dans une localité, des jeunes d'un quartier défavorisé n'ont pas accès à l'eau potable ni à un centre de santé (violence \emph{structurelle}). Un discours circule selon lequel ces jeunes seraient \og naturellement paresseux \fg{} (violence \emph{culturelle}). Un jour, une rixe éclate et un jeune est blessé à coups de machette (violence \emph{directe}). Les trois formes sont liées : les préjugés et les inégalités nourrissent souvent les violences visibles.
\end{exemplebox}

\courssub{Les lieux et contextes des violences}

La violence n'a pas de domicile fixe : elle peut survenir partout où vivent les hommes. Une observation attentive de la vie quotidienne au Cameroun permet de dresser la liste des principaux contextes.

\begin{itemize}
\item \textbf{La famille} : violences entre époux, sévices corporels infligés aux enfants, privation d'héritage des veuves et des orphelins, mariages forcés.
\item \textbf{L'école} : racket entre élèves, brimades, harcèlement, châtiments corporels excessifs, abus de certains encadreurs.
\item \textbf{La rue et les lieux publics} : agressions, vols à l'arraché, rixes entre bandes, comme dans le cas de l'élève de Yaoundé.
\item \textbf{Les réseaux sociaux} : cyberharcèlement, diffamation, diffusion d'images humiliantes, recrutement par des groupes radicaux, comme l'a vécu le jeune de Maroua.
\item \textbf{Les zones de conflits armés} : dans certaines régions du Cameroun, les populations subissent exactions, enlèvements, destructions de villages et d'écoles.
\end{itemize}

\begin{exoresolu}[Classer des faits de violence]
\textbf{Énoncé.} Voici quatre faits observés dans un quartier : (a) un élève est frappé et dépouillé de son argent à la sortie du lycée ; (b) une jeune fille reçoit chaque jour des messages insultants sur son téléphone ; (c) la maison d'un commerçant est vandalisée pendant la nuit ; (d) une veuve est chassée de sa concession par la famille de son défunt mari. Pour chaque fait, indique le lieu de la violence et la ou les catégories d'atteintes.
\tcblower
\textbf{Solution détaillée.}
\begin{itemize}
\item (a) Lieu : la rue (abords de l'école). Catégories : violence \emph{physique} (coups) et \emph{patrimoniale} (vol de l'argent).
\item (b) Lieu : les réseaux sociaux / le téléphone. Catégorie : violence \emph{morale} (cyberharcèlement).
\item (c) Lieu : la rue / le domicile. Catégorie : violence \emph{patrimoniale} (dégradation de biens).
\item (d) Lieu : la famille. Catégories : violence \emph{morale} (expulsion, humiliation) et \emph{patrimoniale} (spoliation de la concession).
\end{itemize}
On remarque qu'un même fait peut relever de plusieurs catégories : c'est fréquent, et c'est pourquoi la qualification juridique demande de la rigueur.
\end{exoresolu}

\courssub{Les facteurs qui favorisent les violences}

Pourquoi la violence éclate-t-elle davantage à certains endroits et à certaines époques ? L'observation sociologique montre que certains \textbf{facteurs favorisants} créent un terrain propice, sans pour autant justifier ni excuser les actes commis.

\begin{itemize}
\item \textbf{La pauvreté et le chômage} : le manque de ressources et l'oisiveté poussent certains jeunes vers le vol, le racket ou les groupes radicaux qui leur promettent de l'argent.
\item \textbf{La consommation de drogues et d'alcool} : elle diminue le contrôle de soi et multiplie les agressions, notamment dans les rues et les fêtes.
\item \textbf{L'intolérance} : le rejet de l'autre (différence ethnique, religieuse, politique ou d'opinion) transforme les désaccords en affrontements.
\item \textbf{L'impunité} : lorsque les auteurs de violences ne sont pas poursuivis ni punis, d'autres sont encouragés à les imiter. La justice qui fonctionne est donc un rempart contre la violence.
\end{itemize}

\begin{attention}[Facteur favorisant n'est pas excuse]
Dire que la pauvreté favorise la délinquance ne signifie pas que le pauvre est un délinquant, ni que le délinquant est innocent. La plupart des personnes pauvres sont honnêtes et dignes. Le facteur favorisant explique un \emph{risque collectif} ; la responsabilité de l'acte reste \emph{individuelle}. Au Baccalauréat technique, évite les raccourcis du type \og le chômage rend violent \fg{} : rédige plutôt \og le chômage est un facteur qui peut favoriser le passage à l'acte violent \fg{}.
\end{attention}

\courssub{Les conséquences des violences}

La violence ne s'arrête jamais au moment où le coup est porté : ses effets se propagent comme des cercles dans l'eau.

\begin{itemize}
\item \textbf{Sur la victime} : blessures physiques parfois définitives (handicap, cicatrices), traumatismes psychologiques (peur, angoisse, dépression, perte d'estime de soi), pertes matérielles, abandon de l'école ou du travail. L'élève de Yaoundé, au-delà de son séjour à l'hôpital, peut développer une peur durable de sortir seule.
\item \textbf{Sur la famille} : frais médicaux, appauvrissement, tensions et ruptures, deuil. La violence faite à un membre fragilise tout le foyer.
\item \textbf{Sur la cohésion sociale} : la violence installe la méfiance entre voisins, entre groupes, entre citoyens et institutions. Elle freine le développement économique (qui investit dans un quartier dangereux ?), détruit le vivre-ensemble et peut dégénérer en cycles de vengeance.
\end{itemize}

\begin{popfigure}
\begin{center}
\begin{tikzpicture}
\begin{axis}[
  width=0.72\linewidth, height=6.2cm,
  ybar, bar width=14pt,
  ymin=0, ymax=40,
  ylabel={Part des plaintes (\%)},
  symbolic x coords={Physiques, Morales, Patrimoniales, Mixtes},
  xtick=data,
  x tick label style={font=\footnotesize},
  nodes near coords, nodes near coords style={font=\footnotesize},
  every axis plot/.append style={fill=popblue, draw=popdark},
  axis lines=left,
  ylabel style={font=\small}
]
\addplot coordinates {(Physiques,32) (Morales,18) (Patrimoniales,35) (Mixtes,15)};
\end{axis}
\end{tikzpicture}
\end{center}
\legende{Diagramme 1 : répartition indicative des plaintes pour violences enregistrées dans une région du Cameroun au cours d'une année (données fictives réalistes, à but pédagogique). Lecture : les atteintes patrimoniales (35 \%) et physiques (32 \%) dominent, mais les violences morales (18 \%) sont probablement sous-déclarées, car beaucoup de victimes n'osent pas porter plainte pour des insultes ou du harcèlement. Les affaires \og mixtes \fg{} (15 \%) combinent plusieurs catégories.}
\end{popfigure}

\courssub{Le cadre juridique de protection}

Face aux violences, le citoyen camerounais n'est pas démuni : un ensemble de textes protège l'intégrité des personnes et des biens.

\begin{aretenir}[Les grands textes protecteurs]
\begin{itemize}
\item \textbf{La Constitution du 18 janvier 1996} : son préambule affirme l'attachement du peuple camerounais aux libertés fondamentales et garantit notamment le droit à la vie, à l'intégrité physique et morale, et le droit de propriété.
\item \textbf{Le Code pénal camerounais} : il définit et punit les coups et blessures, les homicides, les menaces, les diffamations, les vols, les escroqueries et les dégradations.
\item \textbf{Les conventions internationales ratifiées par le Cameroun} : la Déclaration universelle des droits de l'homme (1948), les pactes internationaux relatifs aux droits de l'homme, la Convention sur l'élimination de toutes les formes de discrimination à l'égard des femmes, la Convention relative aux droits de l'enfant.
\end{itemize}
\end{aretenir}

\begin{savaistu}[La Constitution de 1996 et les droits de l'homme]
La Constitution camerounaise du 18 janvier 1996 affirme dans son préambule l'attachement du peuple camerounais aux libertés fondamentales inscrites dans la \textbf{Déclaration universelle des droits de l'homme}, dans la Charte des Nations unies et dans la Charte africaine des droits de l'homme et des peuples. Mieux : le préambule fait partie intégrante de la Constitution, ce qui signifie que ces droits ont une \emph{valeur constitutionnelle} au Cameroun. Protéger les personnes contre les violences n'est donc pas une simple politesse : c'est une obligation inscrite au sommet de la hiérarchie des normes du pays.
\end{savaistu}

Savoir citer un texte ne suffit pas : il faut savoir \emph{qualifier juridiquement} un fait concret, c'est-à-dire lui donner son nom exact en droit. C'est une compétence exigible au Baccalauréat technique.

\begin{methode}[Qualifier juridiquement un fait de violence]
\begin{enumerate}
\item \textbf{Dégager les faits} : relire le récit et relever objectivement ce qui s'est passé (qui ? quoi ? où ? quel dommage ?), sans interprétation.
\item \textbf{Identifier la catégorie d'atteinte} : le corps (physique), la dignité (morale) ou les biens (patrimoniale) sont-ils touchés ? Plusieurs catégories peuvent se cumuler.
\item \textbf{Donner la qualification} : nommer l'infraction correspondante (coups et blessures, vol, menaces, diffamation, dégradation, escroquerie...).
\item \textbf{Citer le texte applicable} : indiquer le texte qui réprime ce fait (Code pénal, Constitution, convention internationale).
\item \textbf{Conclure} : préciser la voie de recours pour la victime (plainte à la gendarmerie ou à la police, saisine du procureur, associations d'aide aux victimes).
\end{enumerate}
\end{methode}

\begin{exoresolu}[Appliquer la méthode de qualification]
\textbf{Énoncé.} À la sortie du lycée, Amina est frappée par deux individus qui emportent son téléphone et son argent ; elle est hospitalisée trois jours. Qualifie juridiquement ces faits et indique le texte applicable ainsi que la démarche à suivre.
\tcblower
\textbf{Solution détaillée.}
\begin{enumerate}
\item \emph{Faits} : coups portés à Amina, blessures ayant entraîné une hospitalisation, soustraction de son téléphone et de son argent.
\item \emph{Catégories d'atteintes} : atteinte à l'intégrité \emph{physique} (coups et blessures) et atteinte \emph{patrimoniale} (soustraction des biens).
\item \emph{Qualification} : coups et blessures volontaires d'une part ; vol avec violence (aggravé par la réunion de plusieurs auteurs) d'autre part.
\item \emph{Texte applicable} : le Code pénal camerounais, qui réprime les coups et blessures volontaires ainsi que le vol, la violence constituant une circonstance aggravante ; la Constitution de 1996 garantit le droit à l'intégrité physique et à la propriété.
\item \emph{Conclusion} : Amina (ou ses parents) doit porter plainte auprès de la gendarmerie ou du commissariat de police, se faire établir un certificat médical constatant les blessures, afin que les auteurs soient recherchés, poursuivis et punis, et que les biens volés soient restitués ou indemnisés.
\end{enumerate}
\end{exoresolu}

\begin{exercice}[À toi de jouer]
Dans un quartier de Douala, un commerçant reçoit chaque semaine des messages l'accusant faussement de vendre des produits périmés ; sa clientèle diminue. Un soir, sa devanture est brisée. (1) Dégage les faits. (2) Qualifie juridiquement chaque fait. (3) Cite le texte applicable. (4) Indique la démarche que doit suivre le commerçant.
\end{exercice}

\begin{corrige}[Exercice : le commerçant de Douala]
(1) \emph{Faits} : diffusion répétée de fausses accusations portant atteinte à la réputation du commerçant ; destruction de sa devanture ; préjudice économique (perte de clientèle).\par
(2) \emph{Qualification} : les messages mensongers constituent une \textbf{diffamation} (atteinte à l'honneur, donc violence morale) ; la devanture brisée constitue une \textbf{dégradation de bien} (violence patrimoniale).\par
(3) \emph{Texte applicable} : le Code pénal camerounais, qui réprime la diffamation et la destruction ou dégradation du bien d'autrui ; la Constitution de 1996 protège la dignité et la propriété.\par
(4) \emph{Démarche} : le commerçant doit conserver les preuves (captures des messages, photos des dégâts, témoignages), porter plainte auprès de la police ou de la gendarmerie, et peut demander réparation du préjudice subi devant les juridictions compétentes.
\end{corrige}

\coursec{L'extrémisme violent}

\courssub{Définition et caractéristiques}

Au-delà des violences \og ordinaires \fg{} (vols, coups, harcèlement), une forme particulière menace gravement la sécurité des États et des populations : l'extrémisme violent.

\begin{definition}[Extrémisme violent et radicalisation]
L'\cle{extrémisme violent} est l'adhésion à une idéologie radicale (politique, religieuse, identitaire) qui \textbf{justifie l'usage de la violence} pour imposer ses vues, renverser l'ordre établi ou détruire un groupe désigné comme ennemi.\par
La \cle{radicalisation} est le \textbf{processus progressif} par lequel un individu adopte cette idéologie et accepte de passer à l'acte violent. Ce n'est pas un événement brutal, mais une trajectoire en plusieurs étapes.
\end{definition}

\begin{popfigure}
\begin{center}
\begin{tikzpicture}[
  font=\small,
  etape/.style={rectangle, rounded corners, draw=poporange, fill=poporangeL, thick, text width=3.2cm, align=center, minimum height=1.8cm},
  fleche/.style={-{Stealth[length=3mm]}, thick, popdark}
]
\node[etape, align=center] (E1) at (0,0){\textbf{1. Pré-radicalisation}\\\emph{Fragilités} : quête de sens, exclusion, échec scolaire/professionnel, rupture familiale.};
\node[etape, align=center] (E2) at (5,0){\textbf{2. Radicalisation}\\\emph{Endoctrinement} : rencontre avec un recruteur, rupture avec l'entourage, adoption d'un discours binaire (nous/eux), consommation de propagande.};
\node[etape, align=center] (E3) at (10,0){\textbf{3. Engagement}\\\emph{Passage à l'acte} : intégration dans un groupe, entraînement, préparation logistique, rupture définitive avec la société.};
\node[etape, align=center] (E4) at (15,0){\textbf{4. Action violente}\\\emph{Commission} : attentat, assassinat, enlèvement, sabotage, ou soutien logistique/financier.};
\draw[fleche] (E1.east) -- (E2.west);
\draw[fleche] (E2.east) -- (E3.west);
\draw[fleche] (E3.east) -- (E4.west);
\node[draw=popgreen, fill=popgreenL, rounded corners, thick, align=center, text width=4cm, font=\footnotesize] at (7.5,-3) {\textbf{Points de bascule / Prévention} :\par
Famille, école, leaders religieux, pairs, services sociaux peuvent intervenir à chaque étape pour \emph{désamorcer} le processus.};
\draw[dashed, popgreen, thick] (7.5,-2.2) -- (7.5,-1.2);
\end{tikzpicture}
\end{center}
\legende{Schéma 2 : le processus de radicalisation menant à l'extrémisme violent (inspiré des modèles de prévention). La flèche verte indique les fenêtres d'opportunité pour la prévention.}
\end{popfigure}

\courssub{Les formes et manifestations au Cameroun}

Le Cameroun fait face à plusieurs expressions de l'extrémisme violent, qui touchent particulièrement les jeunes.

\begin{itemize}
\item \textbf{L'insurrection terroriste dans l'Extrême-Nord} : groupes affiliés à Boko Haram / État islamique en Afrique de l'Ouest (ISWAP). Recrutement par endoctrinement religieux déviant, promesses d'argent, de pouvoir ou de mariage. Attaques contre les forces de défense, les marchés, les écoles, enlèvements (ex. lycéennes de Dabanga, 2020).
\item \textbf{La crise dans les régions du Nord-Ouest et du Sud-Ouest (NOSO)} : groupes séparatistes armés. Recrutement par discours identitaire, haine de l'État, promesses d'indépendance. Violences : attaques d'écoles, enlèvements d'élèves/enseignants, destruction d'infrastructures, imposition de \og ghost towns \fg{}.
\item \textbf{Le cyber-radicalisation} : utilisation des réseaux sociaux (WhatsApp, Facebook, Telegram, TikTok) pour la propagande, le recrutement, la formation à distance (fabrication d'explosifs, maniement d'armes). Le cas du jeune de Maroua (situation d'accroche) illustre cette porte d'entrée numérique.
\end{itemize}

\begin{attention}[Signaux d'alerte (indicateurs de radicalisation)]
Repérer ces signes chez un camarade, un élève ou un proche permet d'alerter à temps : (1) Rupture brutale avec les amis/famille habituels. (2) Changement soudain de vocabulaire (termes guerriers, takfir, \og kouffar \fg{}, \og laïc \fg{} comme insulte). (3) Rejet de l'autorité légitime (école, parents, État) et des symboles nationaux. (4) Isolement, secret sur le téléphone, consultation de sites/chaînes extrémistes. (5) Prosélytisme agressif auprès de l'entourage. \textbf{Attention :} un seul signe ne suffit pas ; c'est la \emph{conjonction} et la \emph{suddeneté} du changement qui doivent alerter.
\end{attention}

\courssub{Conséquences de l'extrémisme violent}

\begin{itemize}
\item \textbf{Humanitaires} : morts, blessés, déplacés internes (plus d'un million au Cameroun), réfugiés, orphelins, traumatismes psychologiques massifs.
\item \textbf{Scolaires} : fermeture d'écoles, perte d'années scolaires, peur d'aller en classe, recrutement d'enfants-soldats.
\item \textbf{Économiques} : destruction de marchés, de champs, d'infrastructures routières ; fuite des investissements ; appauvrissement des zones touchées.
\item \textbf{Sociales} : méfiance intercommunautaire, stigmatisation de régions ou de groupes ethniques/religieux entiers, rupture du lien de confiance avec les forces de l'ordre.
\end{itemize}

\begin{exemplebox}[Témoignage fictif réaliste]
\og J'avais 16 ans quand ils m'ont approché sur WhatsApp. Ils m'ont dit que l'école était \emph{haram}, que l'État nous méprisait, qu'au combat j'aurais de l'argent et le paradis. J'ai failli partir. Mon grand frère a vu les vidéos sur mon téléphone, il m'a parlé, il m'a emmené voir l'imam du quartier. L'imam m'a montré que le Coran interdit de tuer des innocents. J'ai tout effacé. Aujourd'hui, je sensibilise mes camarades. \fg{} — \emph{Abdoul, 18 ans, Maroua.}
\end{exemplebox}

\coursec{Prévenir et combattre les violences : le rôle du citoyen}

\courssub{La prévention : responsabilité partagée}

La lutte contre les violences ne repose pas seulement sur la police ou la justice. Chaque citoyen, à son niveau, est un acteur de la prévention.

\begin{definition}[Prévention situationnelle, sociale et éducative]
\begin{itemize}
\item \textbf{Prévention situationnelle} : réduire les occasions de passer à l'acte (éclairage public, caméras, rondes de quartier, sécurisation des abords d'écoles).
\item \textbf{Prévention sociale} : traiter les causes profondes (emploi jeunes, accès à la santé, lutte contre les drogues, médiation des conflits fonciers/familiaux).
\item \textbf{Prévention éducative} : former à la résolution pacifique des conflits, à l'esprit critique face aux discours de haine, à l'usage responsable des réseaux sociaux (éducation aux médias et à l'information - EMI).
\end{itemize}
\end{definition}

\begin{methode}[Démarche citoyenne face à une situation de violence ou de radicalisation]
\begin{enumerate}
\item \textbf{Observer et évaluer} : suis-je témoin d'une violence (coups, harcèlement, propos radicaux) ? Y a-t-il danger immédiat ?
\item \textbf{Protéger / Alerter} : si danger immédiat $\rightarrow$ appeler le \textbf{112} (urgence) ou le \textbf{113} (police) / \textbf{119} (gendarmerie) ou le numéro vert \textbf{1500} (protection de l'enfance). Ne pas intervenir physiquement si risque pour soi.
\item \textbf{Signaler} : pour une radicalisation suspectée (camarade, voisin), alerter un adulte de confiance (parent, professeur, chef de quartier, imam/prêtre/pasteur) ou les services spécialisés (Comité de vigilance, Bureau de la jeunesse, MINAS, MINPROFF). Le signalement n'est pas une \og dénonciation \fg{}, c'est une \emph{protection}.
\item \textbf{Accompagner} : ne pas isoler la personne en voie de radicalisation ; maintenir le lien, proposer des alternatives (sport, association, formation), orienter vers des structures d'écoute.
\item \textbf{S'engager durablement} : participer aux comités de veille de quartier, aux clubs paix/ECM au lycée, aux campagnes de sensibilisation.
\end{enumerate}
\end{methode}

\courssub{Les institutions d'appui et de protection}

Le citoyen n'agit pas seul : des structures existent pour l'accompagner.

\begin{center}
\begin{tabularx}{\linewidth}{|l|X|}\hline
\textbf{Institution / Structure} & \textbf{Rôle principal pour le citoyen} \\ \hline
Police nationale (Commissariats) & Réception des plaintes, enquêtes, maintien de l'ordre, police de proximité. \\ \hline
Gendarmerie nationale (Brigades) & Même rôle en zone rurale/périurbaine ; police judiciaire militaire. \\ \hline
Tribunaux (TPI, Cour d'Appel) & Jugement des auteurs, réparation pour les victimes (dommages-intérêts). \\ \hline
Commission Nationale des Droits de l'Homme et des Libertés (CNDHL) & Enquêtes sur violations des droits, médiation, rapports, éducation aux droits. \\ \hline
MINAS (Ministère des Affaires Sociales) & Protection de l'enfance, prise en charge victimes (centres d'accueil, assistance psychosociale). \\ \hline
MINPROFF (Ministère de la Promotion de la Femme et de la Famille) & Lutte contre violences basées sur le genre (VBG), centres d'écoute, autonomisation. \\ \hline
Associations / ONG (ex. ALVF, RELUFA, Plan International) & Accompagnement juridique, psychosocial, médical ; plaidoyer ; réinsertion. \\ \hline
Comités de vigilance / Chefferies traditionnelles & Alerte précoce, médiation locale, relais vers forces de l'ordre. \\ \hline
\end{tabularx}
\end{center}

\courssub{La justice transitionnelle et la réconciliation}

Dans les zones sortant de conflit (NOSO, Extrême-Nord), la simple répression ne suffit pas. La \cle{justice transitionnelle} combine :
\begin{itemize}
\item \textbf{Vérité} : commissions d'enquête, reconnaissance des souffrances.
\item \textbf{Justice} : poursuite des crimes graves, mais aussi mécanismes de justice restaurative pour les actes moins graves (réparation, excuses publiques, travaux d'intérêt général).
\item \textbf{Réparation} : indemnisation, réhabilitation psychosociale, reconstruction d'infrastructures.
\item \textbf{Garanties de non-répétition} : réforme du secteur de sécurité, désarmement/démobilisation/réintégration (DDR), éducation à la paix.
\end{itemize}

\begin{savaistu}[Le Plan d'Urgence Humanitaire (PUH) et le DDR au Cameroun]
Depuis 2018, le Cameroun met en œuvre le \textbf{Plan d'Urgence Humanitaire} pour les régions du NOSO et l'Extrême-Nord (reconstruction d'écoles, centres de santé, routes, appui aux déplacés). Parallèlement, les \textbf{Centres de Désarmement, Démobilisation et Réintégration (DDR)} accueillent les ex-combattants qui déposent les armes : ils y reçoivent un accompagnement psychologique, une formation professionnelle (menuiserie, soudure, agriculture, informatique) et un kit de réinsertion. C'est la traduction concrète de la \emph{non-violence} comme politique d'État.
\end{savaistu}

\begin{exoresolu}[Analyser une affiche de prévention]
\textbf{Énoncé.} Une affiche de la CNDHL montre un jeune devant deux portes : l'une marquée \og Radicalisation \fg{} (sombre, fil barbelé), l'autre \og Insertion \fg{} (lumineuse, outils de travail). Slogan : \og Ton avenir se construit, ne le détruis pas. \fg{} Analyse le message et la stratégie de prévention.
\tcblower
\textbf{Solution détaillée.}
\begin{itemize}
\item \emph{Message} : le choix binaire entre destruction de soi (violence) et construction (paix, travail). La radicalisation n'est pas une fatalité, c'un choix qu'on peut refuser.
\item \emph{Stratégie} : \textbf{prévention éducative} (cible les jeunes), \textbf{valorisation de l'alternative positive} (insertion professionnelle), \textbf{responsabilisation individuelle} (\og Ton avenir \fg{}). L'affiche ne fait pas peur, elle \emph{propose une issue}.
\item \emph{Efficacité} : forte si relayée dans les lycées, centres de formation, réseaux sociaux, et si l'offre d'insertion (formation, emploi) est \emph{réelle} et accessible.
\end{itemize}
\end{exoresolu}

\begin{exercice}[Cas pratique : le devoir du citoyen]
Dans ton quartier, tu remarques depuis trois semaines qu'un camarade de promotion, auparavant très sociable, s'isole, refuse de serrer la main aux filles, traite les enseignants de \og mécréants \fg{} et partage sur un groupe WhatsApp de la classe des vidéos de prédicateurs prônant le djihad. (1) Quels signaux d'alerte identifies-tu ? (2) Quelle démarche concrète suis-tu selon la méthode citoyenne ? (3) Pourquoi ne pas garder le secret \og entre nous \fg{} ?
\end{exercice}

\begin{corrige}[Exercice : le camarade en radicalisation]
(1) \emph{Signaux} : rupture sociale (isolement), changement de vocabulaire (termes religieux détournés, hostilité à l'école/laïcité), prosélytisme numérique (partage de propagande), rejet des normes mixtes (refus serrer la main). Conjonction de 4 signaux = alerte sérieuse.\par
(2) \emph{Démarche} : \textbf{Ne pas affronter seul}. En parler immédiatement à un adulte de confiance (professeur principal, CPE, parent). Celui-ci alerte la direction $\rightarrow$ cellule d'écoute/psychologue scolaire $\rightarrow$ si confirmé, signalement au Procureur de la République / Brigade de gendarmerie / Comité de vigilance (procédure \og Alerte Radic \fg{}). Maintenir le lien amical sans valider le discours.\par
(3) \emph{Secret} : le secret protège le \emph{processus violent}, pas la personne. La radicalisation aboutit à l'action (attentat, départ au front) qui tuera ou blessera. Signaler, c'est \emph{sauver} ton camarade et de potentielles victimes. La loi protège le signalement de bonne foi.
\end{corrige}

\begin{aretenir}[L'essentiel de la leçon complète]
\begin{itemize}
\item La \cle{violence} (force physique/psychologique pour nuire) se décline en \textbf{physique}, \textbf{morale}, \textbf{patrimoniale} (objet atteint) et en \textbf{directe}, \textbf{structurelle}, \textbf{culturelle} (mode d'exercice).
\item Elle sévit dans la famille, l'école, la rue, les réseaux sociaux, les zones de conflit. Facteurs favorisants : pauvreté, chômage, drogues, intolérance, impunité (jamais des excuses).
\item Conséquences : traumatismes victimes, fragilisation familles, rupture cohésion sociale, frein au développement.
\item Cadre juridique : Constitution 1996, Code pénal, Conventions internationales. Compétence clé : \textbf{qualifier juridiquement} (faits $\rightarrow$ qualification $\rightarrow$ texte $\rightarrow$ recours).
\item L'\cle{extrémisme violent} : idéologie justifiant la violence pour imposer ses vues. Processus : pré-radicalisation $\rightarrow$ radicalisation $\rightarrow$ engagement $\rightarrow$ action. Formes au Cameroun : terrorisme (Extrême-Nord), séparatisme armé (NOSO), cyber-radicalisation.
\item \textbf{Rôle du citoyen} : prévenir (éducation, vigilance, signalement 112/113/119/1500), protéger, accompagner, s'engager (clubs paix, comités veille). Institutions relais : Police, Gendarmerie, Justice, CNDHL, MINAS, MINPROFF, ONG, Chefferies.
\item Justice transitionnelle (vérité, justice, réparation, non-répétition) et DDR pour sortir durablement des cycles de violence.
\end{itemize}
\end{aretenir}

\coursec{Les atteintes à l'intégrité physique}

Dans la section précédente, nous avons vu que la violence est tout comportement qui porte atteinte à une personne, et qu'elle peut prendre plusieurs formes : physique, morale ou patrimoniale. Nous allons maintenant étudier la forme la plus visible et la plus grave : les atteintes à l'intégrité physique, c'est-à-dire les violences qui touchent directement le corps et la vie de l'être humain. Ces atteintes sont les plus sévèrement punies par la loi camerounaise, car elles touchent au bien le plus précieux : la vie et la santé.

\courssub{Définition et portée de l'atteinte à l'intégrité physique}

Commençons par une situation concrète. À la sortie de l'atelier, deux élèves se disputent pour une place dans le car. L'un pousse l'autre, qui tombe et se fracture le poignet. Le lendemain, dans le quartier, une rixe entre jeunes se termine par un blessé grave transporté à l'hôpital. Dans les deux cas, le corps d'une personne a été atteint : c'est une atteinte à l'intégrité physique.

\begin{definition}[Atteinte à l'intégrité physique]
Une \cle{atteinte à l'intégrité physique} est toute action volontaire qui porte atteinte au \cle{corps} et à la \cle{vie} d'une personne : coups, blessures, mutilations, viols, tortures, homicide. Elle viole le droit à la vie et à l'intégrité corporelle, droits garantis par le préambule de la Constitution camerounaise et par la Déclaration universelle des droits de l'homme (article 3 : « Tout individu a droit à la vie, à la liberté et à la sûreté de sa personne »).
\end{definition}

\begin{aretenir}[L'essentiel]
L'intégrité physique est un droit fondamental. Nul ne peut porter atteinte au corps d'autrui, même pour « punir », « éduquer » ou « venger ». Toute atteinte volontaire au corps ou à la vie est une infraction pénale : elle engage la responsabilité de son auteur devant les tribunaux.
\end{aretenir}

\courssub{Les coups et blessures volontaires}

\begin{definition}[Coups et blessures volontaires]
Les \cle{coups et blessures volontaires} désignent le fait de frapper, blesser ou commettre toute autre violence sur une personne, de manière intentionnelle, sans intention de donner la mort. La gravité de l'infraction est mesurée par les conséquences : simples coups, blessures légères, blessures entraînant une incapacité de travail, mutilation ou infirmité permanente.
\end{definition}

Les exemples sont nombreux dans la vie courante : les \cle{rixes} entre groupes de jeunes, les agressions à la sortie des établissements ou dans les quartiers, les bagarres lors de matchs de football, les coups portés lors de disputes familiales ou de conflits de voisinage.

Le Code pénal camerounais (Loi n°2016/007 du 12 juillet 2016, articles 277 à 280) punit les coups et blessures volontaires. La peine augmente avec la gravité des conséquences :

\begin{center}
\begin{tabularx}{\linewidth}{|l|X|}\hline
\rowcolor{popblueL} \textbf{Gravité des faits} & \textbf{Sanction encourue (ordre de grandeur)} \\ \hline
Coups ou violences simples (sans ITT ou ITT $\le$ 8 jours) & Emprisonnement de 6 jours à 6 mois et amende \\ \hline
Blessures avec ITT $>$ 8 jours & Emprisonnement de 6 mois à 5 ans et amende \\ \hline
Blessures entraînant mutilation ou infirmité permanente & Réclusion criminelle de 10 à 20 ans \\ \hline
Violences ayant entraîné la mort sans intention de la donner (coups mortels) & Réclusion criminelle de 5 à 10 ans \\ \hline
\end{tabularx}
\end{center}

\begin{attention}[Piège fréquent]
Beaucoup croient que « ce n'était qu'une bagarre » ou que « l'autre a commencé ». Erreur : dès qu'il y a violence volontaire, il y a infraction, même sans blessure grave. De plus, la provocation n'efface pas la faute : répondre à une insulte par des coups, c'est commettre un délit. La légitime défense n'est reconnue que si la riposte est immédiate, nécessaire et proportionnée à l'agression.
\end{attention}

\courssub{L'homicide : meurtre, assassinat, infanticide}

L'atteinte la plus grave à l'intégrité physique est la suppression de la vie humaine : l'homicide.

\begin{definition}[Homicide et ses formes]
L'\cle{homicide} est le fait de donner volontairement la mort à autrui. Le Code pénal (article 276) distingue plusieurs formes :
\begin{itemize}
\item le \cle{meurtre} : homicide volontaire commis sans préméditation. Il est puni de la \cle{réclusion criminelle} de dix à vingt ans ;
\item l'\cle{assassinat} : meurtre commis avec \cle{préméditation} (le meurtrier a préparé et planifié son acte) ou avec guet-apens. Il est puni de la \cle{réclusion criminelle à perpétuité} ;
\item l'\cle{infanticide} : meurtre d'un nouveau-né, souvent commis par la mère dans des circonstances particulières (détresse, abandon, pression sociale). Il constitue un meurtre mais peut bénéficier de circonstances atténuantes liées à l'état de la mère.
\end{itemize}
\end{definition}

\begin{exemplebox}[Distinction à retenir]
Deux voisins se disputent violemment ; l'un, emporté par la colère, frappe l'autre qui meurt : c'est un \textbf{meurtre} (réclusion criminelle 10--20 ans). En revanche, si quelqu'un prépare pendant plusieurs jours la mort d'une personne (achat d'une arme, choix du lieu et du moment), il commet un \textbf{assassinat}, puni de la réclusion criminelle à perpétuité car prémédité.
\end{exemplebox}

\courssub{Les violences faites aux femmes et aux enfants}

Certaines catégories de personnes sont particulièrement exposées aux violences physiques : les femmes et les enfants.

\textbf{Les violences faites aux femmes.} On distingue notamment :
\begin{itemize}
\item les \cle{violences conjugales} : coups, blessures, sévices infligés par le conjoint ou le partenaire, souvent de manière répétée ;
\item le \cle{viol} : tout acte sexuel imposé par la force, la contrainte ou la surprise. Le viol est un crime très sévèrement puni par le Code pénal camerounais (article 296) : réclusion criminelle de 10 à 20 ans, portée à la réclusion criminelle à perpétuité lorsque la victime est mineure de 16 ans ou en cas de circonstances aggravantes ;
\item les \cle{mutilations}, notamment les mutilations génitales féminines (excision), pratiques traditionnelles néfastes qui portent une atteinte grave et irréversible au corps des filles et des femmes.
\end{itemize}

\textbf{Les violences faites aux enfants.} Elles comprennent les \cle{châtiments corporels abusifs} (coups excessifs infligés par des parents ou des enseignants), les sévices, l'exploitation et les abus sexuels. Corriger un enfant ne donne à personne le droit de le blesser : la correction éducative doit rester mesurée et ne jamais porter atteinte au corps ou à la dignité de l'enfant.

\begin{exemplebox}[Un exemple chiffré]
Le Code pénal camerounais punit sévèrement le viol et les violences sur mineurs : selon les circonstances (âge de la victime, usage de la force, relation d'autorité), les peines de \textbf{réclusion criminelle} peuvent aller jusqu'à la \textbf{réclusion criminelle à perpétuité}. Les violences sur un enfant de moins de 15 ans constituent une circonstance aggravante : la loi protège tout particulièrement les plus vulnérables.
\end{exemplebox}

\begin{popfigure}
\begin{tikzpicture}
\begin{axis}[
    ybar,
    bar width=14pt,
    width=0.9\linewidth,
    height=6cm,
    ymin=0, ymax=6000,
    symbolic x coords={2018,2019,2020,2021,2022},
    xtick=data,
    ylabel={Cas signalés (ordre de grandeur)},
    xlabel={Année},
    ylabel style={font=\small},
    xlabel style={font=\small},
    tick label style={font=\small},
    nodes near coords,
    every node near coord/.append style={font=\scriptsize},
    axis lines=left,
    enlarge x limits=0.15
]
\addplot[fill=popblue, draw=popdark] coordinates {(2018,2400) (2019,2900) (2020,4100) (2021,4600) (2022,5200)};
\end{axis}
\end{tikzpicture}
\legende{Évolution indicative des cas de violences basées sur le genre (VBG) signalés au Cameroun (ordres de grandeur, d'après les données UNFPA/MINPROFF). La hausse s'explique à la fois par l'aggravation des violences dans les zones de crise et par la libération progressive de la parole des victimes grâce aux campagnes de sensibilisation et aux centres d'écoute.}
\end{popfigure}

L'histogramme ci-dessus montre une augmentation régulière des cas signalés. Deux lectures sont possibles : d'une part, les crises sécuritaires et les déplacements de populations exposent davantage les femmes et les enfants aux violences ; d'autre part, les victimes osent de plus en plus porter plainte, ce qui est un progrès. Un chiffre qui augmente ne signifie donc pas seulement plus de violences : il peut aussi signifier moins de silence.

\courssub{Les violences en milieu scolaire}

L'école et le lycée technique ne sont pas épargnés. On y observe :
\begin{itemize}
\item le \cle{bizutage} : brimades, humiliations ou épreuves imposées aux nouveaux élèves par les anciens. Même présenté comme une « tradition », le bizutage qui comporte des coups, des blessures ou des traitements dégradants est une infraction ;
\item les \cle{brutalités entre élèves} : bagarres, racket accompagné de coups, agressions dans les cours de récréation ou aux abords de l'établissement ;
\item les \cle{châtiments excessifs} : punitions corporelles disproportionnées infligées par des encadreurs, contraires à la réglementation scolaire et à la dignité de l'élève.
\end{itemize}

\begin{attention}[Le silence et la peur]
Le silence et la peur favorisent la répétition des violences : un agresseur non dénoncé recommence, et d'autres élèves deviennent victimes à leur tour. \textbf{Dénoncer une violence est un devoir civique, pas une délation.} La délation vise à nuire par intérêt ou par méchanceté ; la dénonciation protège une victime et fait respecter la loi. Témoin d'une violence, l'élève doit alerter un enseignant, le chef d'établissement, ses parents ou les forces de l'ordre.
\end{attention}

\courssub{Les atteintes liées aux conflits}

Les conflits qui touchent certaines régions du Cameroun aggravent les atteintes à l'intégrité physique. Dans l'Extrême-Nord, les exactions du groupe \cle{Boko Haram} (enlèvements, tueries, mutilations, recrutement forcé) ont frappé des milliers de civils. Dans les régions du Nord-Ouest et du Sud-Ouest, la \cle{crise anglophone} s'accompagne d'assassinats, d'enlèvements, de blessures par armes à feu et d'attaques contre les écoles. Ces violences touchent particulièrement les populations déplacées, les femmes et les enfants, et rappellent que la paix est la première condition du respect de l'intégrité physique de tous.

\begin{savaistu}[Le saviez-vous ?]
Le Cameroun dispose de mécanismes d'assistance aux victimes de violences basées sur le genre : des \textbf{centres d'écoute} et d'accueil mis en place avec l'appui du MINPROFF (Ministère de la Promotion de la Femme et de la Famille) et de partenaires comme l'UNFPA, ainsi que des \textbf{numéros d'appel d'urgence} permettant d'alerter et d'orienter les victimes. Des associations comme l'ALVF (Association de Lutte contre les Violences faites aux Femmes) accompagnent juridiquement et psychologiquement les victimes. Des campagnes nationales de sensibilisation, comme la campagne « 16 jours d'activisme contre les violences faites aux femmes » (du 25 novembre au 10 décembre), mobilisent chaque année les communautés.
\end{savaistu}

\courssub{Les réponses institutionnelles face aux atteintes physiques}

Face à une atteinte à l'intégrité physique, la victime et ses proches ne sont pas démunis. La démarche à suivre est précise.

\begin{methode}[Que faire en cas d'atteinte à l'intégrité physique ?]
\begin{enumerate}
\item \textbf{Se mettre en sécurité} et, si nécessaire, appeler les secours ou alerter un adulte responsable.
\item \textbf{Consulter un médecin} qui établit un \cle{certificat médical} : ce document décrit les blessures, évalue l'incapacité temporaire de travail (ITT) et constitue une preuve essentielle.
\item \textbf{Déposer plainte} auprès de la \cle{police} ou de la \cle{gendarmerie}, qui enquêtent et transmettent le dossier au procureur de la République.
\item Laisser la justice agir : les \cle{tribunaux} jugent l'auteur des faits et prononcent les sanctions pénales ; la victime peut se constituer partie civile pour obtenir réparation.
\item \textbf{Se faire accompagner} par les ONG et associations de défense des droits (ALVF, centres d'écoute, organisations de défense des droits humains) qui offrent soutien juridique, médical et psychologique.
\end{enumerate}
\end{methode}

\begin{exoresolu}[Étude de cas : la réaction de Chantal]
\textbf{Énoncé.} Chantal, élève en classe de Première dans un lycée technique, est régulièrement frappée par son concubin. Un soir, il la blesse au bras avec un objet tranchant. Par peur et par honte, Chantal hésite à réagir. Sa camarade Amina lui conseille de « se taire pour préserver la famille ».
\begin{enumerate}
\item De quelle forme d'atteinte à l'intégrité physique Chantal est-elle victime ?
\item Le conseil d'Amina est-il judicieux ? Justifiez.
\item Décrivez la démarche que Chantal devrait suivre.
\end{enumerate}
\tcblower
\textbf{Solution détaillée.}
\begin{enumerate}
\item Chantal est victime de \textbf{violences conjugales}, qui constituent des \textbf{coups et blessures volontaires} : son concubin lui porte atteinte au corps de manière intentionnelle et répétée. L'usage d'un objet tranchant constitue une circonstance aggravante (arme par destination).
\item Le conseil d'Amina n'est \textbf{pas judicieux}. Le silence et la peur favorisent la répétition des violences : un agresseur non inquiété recommence, souvent avec plus de gravité. Dénoncer n'est pas une délation mais un acte de protection de soi et un devoir civique. Préserver une « famille » ne peut justifier de laisser détruire la santé et la vie d'une personne.
\item Chantal devrait : (a) se mettre en sécurité et se faire soigner ; (b) demander un \textbf{certificat médical} constatant ses blessures et l'ITT ; (c) \textbf{déposer plainte} au commissariat de police ou à la brigade de gendarmerie ; (d) se faire accompagner par une association comme l'ALVF ou un centre d'écoute ; (e) laisser les tribunaux juger les faits, en se constituant éventuellement partie civile pour obtenir réparation.
\end{enumerate}
\end{exoresolu}

\begin{aretenir}[Conclusion de la section]
Les atteintes à l'intégrité physique vont des coups et blessures volontaires jusqu'à l'homicide, en passant par les violences faites aux femmes et aux enfants, les violences scolaires et les exactions liées aux conflits. Toutes sont des infractions graves punies par le Code pénal camerounais, avec des peines pouvant aller jusqu'à la réclusion criminelle à perpétuité. Face à elles, la victime dispose de moyens concrets : certificat médical, dépôt de plainte, action de la police, de la gendarmerie, des tribunaux et des ONG. Rompre le silence, c'est déjà combattre la violence. La section suivante montrera que la violence peut aussi blesser sans toucher le corps : ce sont les atteintes à l'intégrité morale.
\end{aretenir}

\coursec{Les atteintes à l'intégrité morale}

Dans la section précédente, nous avons étudié les atteintes à l'intégrité \cle{physique} : coups et blessures, violences corporelles, sévices. Ces violences laissent des traces visibles sur le corps. Mais il existe d'autres violences, parfois plus sournoises, qui ne laissent aucune cicatrice apparente : ce sont les atteintes à l'intégrité \cle{morale}. Une insulte répétée, une rumeur lancée sur WhatsApp, une moquerie quotidienne dans l'atelier peuvent détruire une personne aussi sûrement qu'un coup de poing. C'est à ces violences invisibles mais bien réelles que nous consacrons cette section.

\courssub{Définition et intuition}

\begin{definition}[Atteinte à l'intégrité morale]
Une \cle{atteinte à l'intégrité morale} est toute action, parole ou écrit qui porte atteinte à l'\textbf{honneur}, à la \textbf{dignité}, à la \textbf{réputation} ou à l'\textbf{équilibre psychologique} d'une personne. Contrairement à la violence physique, elle s'attaque à ce que la personne \emph{est} et à ce que les autres pensent d'elle, et non à son corps. Elle est réprimée par le Code pénal (loi n° 2016/007 du 12 juillet 2016) et la loi n° 2010/012 du 21 décembre 2010 relative à la cybersécurité et à la cybercriminalité.
\end{definition}

Prenons une situation concrète. Marlyse, élève en Première F4 (génie civil), est la seule fille de son groupe d'atelier. Chaque jour, deux camarades répètent qu'« une fille n'a rien à faire sur un chantier », affichent ses erreurs dans le groupe WhatsApp de la classe et inventent des histoires sur sa vie privée. Marlyse n'a reçu aucun coup, pourtant elle ne dort plus, ses notes chutent et elle veut abandonner la filière. Son intégrité physique est intacte, mais son intégrité \emph{morale} est gravement atteinte.

\begin{aretenir}[L'essentiel]
L'intégrité morale d'une personne comprend quatre biens protégés par la loi : son \textbf{honneur} (le respect qui lui est dû), sa \textbf{dignité} (sa valeur d'être humain), sa \textbf{réputation} (l'image qu'elle a dans la société) et son \textbf{équilibre psychologique} (sa santé mentale). Y toucher, même par de simples paroles, est une violence punie par la loi.
\end{aretenir}

\courssub{Les injures, la diffamation et la calomnie}

Ces trois notions sont proches mais juridiquement distinctes. Bien les différencier est un classique du Probatoire technique.

\begin{definition}[Injure, diffamation, calomnie]
\begin{itemize}
\item L'\cle{injure} est une parole, un geste ou un écrit \textbf{offensant} adressé à une personne, sans préciser de fait précis : « voleur », « bonne à rien », insulte ethnique. Elle peut être publique (entendue de tiers) ou non publique.
\item La \cle{diffamation} est l'allégation ou l'imputation d'un \textbf{fait précis} qui porte atteinte à l'honneur ou à la réputation d'une personne, et qui est rendue publique : affirmer devant toute la classe que « Paul a triché à l'examen » alors que c'est faux ou non prouvé.
\item La \cle{calomnie} est une diffamation dont l'auteur \textbf{sait qu'elle est fausse} : c'est une accusation mensongère faite en connaissance de cause, souvent pour nuire. Elle est aggravée par l'intention dolosive.
\end{itemize}
\end{definition}

\begin{center}
\begin{tabularx}{\linewidth}{|l|X|X|X|}
\hline
\rowcolor{popblueL} \textbf{Notion} & \textbf{Définition} & \textbf{Exemple en milieu technique} & \textbf{Bases légales \& sanctions} \\
\hline
Injure & Parole ou écrit offensant, sans fait précis & Traiter un camarade de « nul » ou de « fainéant » devant l'atelier & Code pénal (art. 303 s.) : amende et/ou peine de prison (injures publiques) \\
\hline
Diffamation & Allégation publique d'un fait précis nuisant à l'honneur & Publier sur Facebook qu'un maître de stage détourne du matériel sans preuve & Code pénal (art. 305 s.) + Loi 2010/012 (si en ligne) : amende et emprisonnement \\
\hline
Calomnie & Diffamation que l'auteur sait fausse & Accuser sciemment un collègue de vol pour obtenir son renvoi & Peines aggravées (Code pénal), dommages et intérêts à la victime \\
\hline
\end{tabularx}
\end{center}

\begin{attention}[Piège fréquent]
Beaucoup d'élèves confondent injure et diffamation. Le critère est simple : la diffamation vise un \textbf{fait précis et vérifiable} (« il a volé la caisse le 12 mars »), l'injure est une simple \textbf{qualification insultante} (« voleur ! »). Autre piège : croire que dire la vérité protège toujours. Même un fait vrai, divulgué pour humilier (vie privée, maladie, orientation), peut constituer une atteinte à la dignité ou au respect de la vie privée (Code pénal, art. 307-1 sur le respect de la vie privée).
\end{attention}

\courssub{Le harcèlement et le cyberharcèlement}

\begin{definition}[Harcèlement]
Le \cle{harcèlement} est un ensemble de paroles, gestes ou comportements \textbf{répétés} qui ont pour but ou pour effet de dégrader les conditions de vie d'une personne, de porter atteinte à sa dignité ou de créer un environnement intimidant, hostile ou humiliant. Le mot clé est la \textbf{répétition} : un seul incident peut être une injure ; la répétition systématique caractérise le harcèlement (Code pénal, art. 302-1).
\end{definition}

On distingue plusieurs formes de harcèlement :
\begin{itemize}
\item le \textbf{harcèlement moral} : moqueries, humiliations, isolement, menaces répétées (en classe, à l'atelier, en entreprise) ;
\item le \textbf{harcèlement sexuel} : gestes, propos ou avances à caractère sexuel non désirés et répétés, y compris le « chantage aux notes » ou au stage (Code pénal, art. 299-1) ;
\item le \textbf{harcèlement scolaire} : brimades entre élèves, racket, exclusion systématique d'un camarade ;
\item le \cle{cyberharcèlement} : harcèlement commis par les moyens numériques (réseaux sociaux, SMS, groupes WhatsApp) : diffusion de photos humiliantes, faux comptes, messages insultants jour et nuit. Sa particularité : il ne s'arrête jamais, car il suit la victime jusque chez elle, et il touche un public immense (Loi 2010/012, art. 74).
\end{itemize}

\begin{attention}[« Ce n'était qu'une blague »]
Une « blague » \textbf{répétée} qui humilie un camarade est déjà une violence morale. En droit comme en morale, c'est l'\textbf{effet sur la victime} qui compte, pas seulement l'intention de l'auteur. Dire « je plaisantais » ne suffit pas à excuser : si la personne rit jaune, s'isole ou souffre, la limite est franchie. Le rire du groupe ne transforme pas une humiliation en jeu.
\end{attention}

\begin{methode}[Réagir face au cyberharcèlement : la démarche en 4 étapes]
\begin{enumerate}
\item \textbf{Conserver les preuves} : captures d'écran datées des messages, liens (URL), noms des comptes ou numéros. Ne rien supprimer, même si c'est douloureux à garder.
\item \textbf{Ne pas répondre} aux agresseurs : toute réponse alimente le harcèlement et peut retourner les propos contre la victime.
\item \textbf{Signaler} : à la plateforme (fonction « signaler » de Facebook, WhatsApp, TikTok), puis à un adulte de confiance : parent, enseignant, conseiller d'orientation, chef d'établissement (censeur, proviseur).
\item \textbf{Porter plainte} : auprès de la police ou de la gendarmerie (brigade de gendarmerie compétente en cybercriminalité / antenne ANTIC), en joignant les preuves conservées.
\end{enumerate}
\end{methode}

\begin{savaistu}[La loi camerounaise sur la cybersécurité]
Le Cameroun dispose de la \textbf{loi n° 2010/012 du 21 décembre 2010 relative à la cybersécurité et à la cybercriminalité}. Elle réprime notamment les atteintes à l'honneur et à la réputation commises par les technologies de l'information et de la communication : diffamation en ligne (art. 73), publication de contenus portant atteinte à la dignité d'autrui (art. 74), usurpation d'identité (art. 75). Un message insultant envoyé depuis un téléphone à Douala peut donc conduire son auteur devant un tribunal. Internet n'est pas une zone de non-droit. L'Agence Nationale des Technologies de l'Information et de la Communication (ANTIC) veille à l'application de cette loi.
\end{savaistu}

\courssub{Les discriminations et la stigmatisation}

\begin{definition}[Discrimination et stigmatisation]
La \cle{discrimination} est le fait de traiter une personne de manière défavorable (refus d'emploi, de stage, de logement, mauvais traitement) en raison de son \textbf{appartenance ethnique}, de sa \textbf{religion}, de son \textbf{sexe}, de son \textbf{handicap} ou de son \textbf{orientation}, et non de ses compétences ou de ses actes. La \cle{stigmatisation} en est la face sociale : on colle à une personne ou à un groupe une « étiquette » négative qui l'exclut du regard des autres.
\end{definition}

Exemples dans la vie d'un lycée technique : refuser une candidate à un stage de chaudronnerie « parce que c'est une fille » ; se moquer d'un élève handicapé moteur ; exclure un camarade d'un travail de groupe à cause de son ethnie ou de sa religion. Ces comportements violent le principe d'\textbf{égalité} inscrit dans le préambule de la Constitution camerounaise, qui affirme que « l'être humain, sans distinction de race, de religion, de sexe ou de croyance, possède des droits inaliénables et sacrés ». Ils sont également sanctionnés par le Code pénal (art. 242-1 sur la discrimination).

\courssub{Les rumeurs et les fausses informations : une violence collective}

Une \cle{rumeur} est une information non vérifiée qui circule de bouche à oreille ; une \emph{fake news} (infox) est une fausse information fabriquée pour tromper. Sur les réseaux sociaux, elles se propagent en quelques heures et deviennent des violences morales \textbf{collectives} : des centaines de personnes participent, souvent sans le vouloir, à la destruction de la réputation d'une victime.

\begin{popfigure}
\begin{center}
\begin{tikzpicture}[scale=1, every node/.style={font=\small}, >=Stealth]
% Styles
\tikzset{
  boxA/.style={draw, fill=poporangeL, rounded corners, minimum width=2.2cm, minimum height=0.9cm, thick},
  boxB/.style={draw, fill=popblueL, rounded corners, minimum width=2cm, minimum height=0.8cm},
  boxC/.style={draw, fill=poppurpleL, rounded corners, minimum width=2.6cm, minimum height=0.9cm, thick},
  boxD/.style={draw, fill=poppinkL, rounded corners, minimum width=2.2cm, minimum height=0.9cm, thick},
  boxE/.style={draw, fill=popcream, rounded corners, align=center, font=\footnotesize},
  arrA/.style={->, thick, poporange},
  arrB/.style={->, thick, popblue},
  arrC/.style={->, very thick, poppurple},
  arrD/.style={->, thick, poppink}
}
% Nœuds
\node[boxA] (auteur) at (0,0) {\textbf{Auteur} (rumeur)};
\node[boxB] (r1) at (3.6,1.6) {Relais 1};
\node[boxB] (r2) at (3.6,0) {Relais 2};
\node[boxB] (r3) at (3.6,-1.6) {Relais 3};
\node[boxC] (masse) at (7.2,0) {\textbf{Masse} (partages)};
\node[boxD] (victime) at (10.8,0) {\textbf{Victime}};
\node[boxE, align=center] (effets) at (10.8,-2.6){Humiliation publique \\ Stress, angoisse \\ Déscolarisation \\ Dépression, idées suicidaires};
% Flèches
\draw[arrA] (auteur) -- (r1);
\draw[arrA] (auteur) -- (r2);
\draw[arrA] (auteur) -- (r3);
\draw[arrB] (r1) -- (masse);
\draw[arrB] (r2) -- (masse);
\draw[arrB] (r3) -- (masse);
\draw[arrC] (masse) -- (victime);
\draw[arrD] (victime) -- (effets);
% Commentaire
\node[align=center, font=\footnotesize\itshape, popmuted] at (3.6,2.6) {Chaque partage amplifie la rumeur};
\end{tikzpicture}
\end{center}
\legende{Schéma 1 : circuit d'une rumeur sur les réseaux sociaux. Une seule publication mensongère (auteur) est relayée par quelques contacts (relais 1 à 3), puis amplifiée par les partages de la masse, avant de frapper la victime. Les effets s'enchaînent : humiliation, stress, déscolarisation, dépression.}
\end{popfigure}

La leçon du schéma est claire : \textbf{celui qui partage sans vérifier devient co-auteur de la violence}. Avant de transférer une information, le citoyen numérique se pose trois questions : la source est-elle fiable ? L'information est-elle vérifiable ailleurs (recoupement) ? Suis-je prêt à en assumer les conséquences juridiques et morales ?

\courssub{Conséquences psychologiques et voies de recours}

Les atteintes morales produisent des effets graves et documentés : \textbf{stress} permanent, troubles du sommeil, \textbf{angoisse}, perte d'estime de soi, \textbf{déscolarisation} (refus d'aller en cours ou à l'atelier), \textbf{dépression} et, dans les cas extrêmes, \textbf{suicide}. La victime se tait souvent par honte ou peur des représailles : c'est pourquoi l'entourage et les témoins ont un rôle décisif.

Face à une atteinte morale, plusieurs voies de recours existent :
\begin{itemize}
\item la \textbf{plainte} auprès de la police, de la gendarmerie ou du procureur de la République (action publique) ;
\item la \textbf{médiation} : un tiers neutre (enseignant, chef d'établissement, médiateur communautaire) aide les parties à réparer le tort sans procès, solution souvent adaptée aux conflits entre élèves (justice restaurative) ;
\item la \textbf{protection des témoins} : celui qui signale un harcèlement ne doit subir aucune représailles ; témoigner est un devoir civique (Code de procédure pénale) ;
\item l'\textbf{éducation aux médias et à la citoyenneté numérique} : apprendre à vérifier les sources, à respecter autrui en ligne, à protéger sa vie privée — c'est la prévention la plus efficace.
\end{itemize}

\begin{exoresolu}[Étude de cas : la photo détournée]
Énoncé. Dans un lycée technique de Yaoundé, des élèves détournent une photo de leur camarade Aïcha et la publient dans un groupe WhatsApp avec des commentaires insultants sur son origine ethnique. La publication est reprise sur Facebook. Aïcha, humiliée, cesse de venir en cours. 1) Identifier les atteintes morales commises. 2) Proposer une démarche de résolution conforme à la loi.
\tcblower
Solution détaillée. 1) On relève plusieurs atteintes : des \textbf{injures} (commentaires offensants), une \textbf{discrimination ethnique} (moqueries fondées sur l'origine, art. 242-1 Code pénal), un \textbf{cyberharcèlement} (diffusion répétée et amplifiée en ligne, Loi 2010/012 art. 74) et une atteinte au \textbf{droit à l'image} (art. 76 Loi 2010/012). La déscolarisation d'Aïcha montre l'atteinte à son équilibre psychologique. 2) Démarche : Aïcha conserve les preuves (captures d'écran datées, URLs), ne répond pas aux agresseurs, signale les contenus aux plateformes (WhatsApp, Facebook) et en informe ses parents et le chef d'établissement. Une \textbf{médiation} scolaire peut d'abord être tentée (retrait des contenus, excuses publiques, réparation morale). Si les faits persistent ou sont graves, la famille \textbf{porte plainte} à la gendarmerie (brigade cybercriminalité) ou au procureur. Les témoins qui signalent les faits doivent être protégés. Conclusion : toute publication humiliante en ligne engage la responsabilité de son auteur et de ceux qui la relaient.
\end{exoresolu}

\begin{exercice}[À toi de juger]
Ton voisin d'atelier reçoit chaque jour des SMS anonymes le traitant d'incapable et lui conseillant d'abandonner la filière. Il te montre les messages en pleurant. 1) De quelle forme de violence s'agit-il ? Justifie avec la définition. 2) Rédige en cinq lignes la conduite à tenir, en appliquant la méthode en quatre étapes vue en cours.
\end{exercice}

\begin{corrige}[Exercice « À toi de juger »]
1) Il s'agit de \textbf{harcèlement moral} (par SMS, donc avec une dimension de cyberharcèlement) : des messages insultants \emph{répétés} qui dégradent les conditions de vie de la victime et atteignent son équilibre psychologique (Code pénal art. 302-1, Loi 2010/012 art. 74). 2) Conduite à tenir : conserver les SMS comme preuves (captures, dates, numéro) ; ne pas répondre à l'anonyme ; signaler les faits à un adulte (censeur, enseignant, parents) et à l'opérateur téléphonique ; si les faits continuent, accompagner le camarade pour porter plainte à la gendarmerie. En tant que témoin, je soutiens la victime et je témoigne sans crainte, car la loi protège ceux qui signalent.
\end{corrige}

\begin{aretenir}[Bilan de la section]
Les atteintes à l'intégrité morale touchent l'honneur, la dignité, la réputation et l'équilibre psychologique. Elles prennent la forme d'injures, de diffamations, de calomnies, de harcèlements (dont le cyberharcèlement), de discriminations et de rumeurs. Leurs conséquences peuvent être tragiques (déscolarisation, dépression, suicide). Le citoyen réagit en conservant les preuves, en signalant, en portant plainte et en pratiquant une citoyenneté numérique responsable. La section suivante montrera comment ces violences peuvent aussi viser les biens : les atteintes au patrimoine.
\end{aretenir}

\coursec{Les atteintes au patrimoine}

Dans la section précédente, nous avons étudié les atteintes à l'intégrité morale : la diffamation, l'injure, la violation de la vie privée ou encore le harcèlement, qui blessent une personne dans son honneur et sa dignité sans toucher son corps. Mais la violence ne s'exerce pas seulement contre les personnes. Elle peut aussi s'attaquer à ce qu'elles possèdent : un téléphone arraché dans la rue, un terrain vendu deux fois, une école incendiée, de l'argent public détourné. Ces actes, qui portent atteinte aux \cle{biens} des personnes, constituent les \cle{atteintes au patrimoine}. C'est l'objet de cette section.

\courssub{Définition et cadre général}

\begin{definition}[Atteinte au patrimoine]
Une \cle{atteinte au patrimoine} est toute atteinte portée aux \cle{biens} d'une \cle{personne physique} (un individu) ou d'une \cle{personne morale} (une entreprise, une association, l'État, une commune). Les biens peuvent être \emph{meubles} (argent, téléphone, moto, mobilier) ou \emph{immeubles} (maison, terrain, usine).
\end{definition}

Le \cle{patrimoine} est l'ensemble des biens et des droits qu'une personne possède et qui ont une valeur économique : ses biens meubles, ses biens immeubles, son argent, ses créances. Le droit de propriété est protégé : la Déclaration universelle des droits de l'homme de 1948 affirme en son article 17 que « toute personne a le droit de propriété » et que « nul ne peut être arbitrairement privé de sa propriété ». Au Cameroun, le Code pénal punit sévèrement les atteintes aux biens, car elles détruisent la confiance sans laquelle aucune vie économique et sociale n'est possible.

Pour bien comprendre, distinguons trois grandes familles d'atteintes au patrimoine :
\begin{itemize}
\item les atteintes par \emph{soustraction} : on prend le bien d'autrui (vol, escroquerie, abus de confiance, détournement) ;
\item les atteintes par \emph{destruction} : on détériore ou on détruit le bien d'autrui (dégradations, vandalisme) ;
\item les atteintes par \emph{usurpation} : on s'approprie frauduleusement un droit (spoliation foncière, double vente).
\end{itemize}

\courssub{Le vol : simple et aggravé}

\begin{definition}[Vol]
Le \cle{vol} est la soustraction frauduleuse de la chose d'autrui, c'est-à-dire le fait de prendre un bien qui ne nous appartient pas, sans le consentement de son propriétaire et avec l'intention de se l'approprier.
\end{definition}

Le \emph{vol simple} est le vol commis sans circonstance aggravante : un pickpocket qui dérobe un portefeuille dans un marché de Douala, un élève qui prend discrètement le téléphone de son voisin. Mais le Code pénal camerounais aggrave les peines lorsque le vol est commis avec certaines circonstances :

\begin{itemize}
\item le vol \emph{avec violence} : l'agresseur frappe ou menace la victime pour lui arracher ses biens (agression dans une rue de Yaoundé la nuit) ;
\item le vol \emph{à main armée} : le voleur utilise une arme, vraie ou fausse (braquage d'un « micro », c'est-à-dire d'un taxi ou d'un car de transport urbain, par des bandits armés) ;
\item le vol \emph{en réunion} : le vol est commis par plusieurs personnes associées (une bande organisée) ;
\item le vol \emph{commis la nuit} ou avec \emph{effraction} (briser une porte, escalader une clôture).
\end{itemize}

\begin{exemplebox}[Les braquages de « micros » en milieu urbain]
Dans les grandes villes camerounaises, des bandits montent parfois dans un taxi ou un car comme de simples passagers, puis dégainent une arme et dépouillent tous les occupants de leur argent et de leurs téléphones. Il s'agit d'un vol aggravé (à main armée, en réunion), passible de lourdes peines de prison. La victime doit porter plainte au commissariat ou à la gendarmerie le plus proche.
\end{exemplebox}

\begin{attention}[Le recel est aussi un délit]
Acheter un bien que l'on sait volé (un téléphone « pas cher » sans facture, une moto sans carte grise) constitue le délit de \cle{recel}. Le receleur est puni par la loi \emph{même s'il n'a pas participé au vol}. Refuser d'acheter un objet d'origine douteuse, c'est déjà lutter contre le vol.
\end{attention}

\courssub{L'escroquerie, l'abus de confiance et les arnaques en ligne}

\begin{definition}[Escroquerie et abus de confiance]
L'\cle{escroquerie} est le fait d'obtenir un bien, de l'argent ou un service par des \emph{manœuvres frauduleuses} : mensonges, faux documents, fausse identité. La victime est trompée et remet « volontairement » le bien, mais son consentement est vicié par le mensonge.
L'\cle{abus de confiance} est le fait de détourner un bien qui nous a été \emph{confié} (en dépôt, en location, pour un travail) et que l'on refuse de restituer ou que l'on utilise à des fins personnelles.
\end{definition}

La différence avec le vol est essentielle : dans le vol, on \emph{prend} ; dans l'escroquerie, on \emph{fait donner} par la ruse ; dans l'abus de confiance, on \emph{garde indûment} ce qui nous a été remis.

\begin{exemplebox}[Les arnaques en ligne et par mobile money]
Avec la généralisation des téléphones portables et du mobile money au Cameroun, les escroqueries se sont déplacées sur les réseaux :
\begin{itemize}
\item les \emph{faux concours} : un message annonce que vous avez gagné un lot et vous demande d'envoyer d'abord des frais de « déblocage » ;
\item les \emph{faux agents mobile money} : un inconnu prétend qu'un transfert a été fait par erreur sur votre compte et vous demande de le « retourner », alors qu'aucun transfert n'a eu lieu ;
\item les \emph{fausses annonces de vente} sur les réseaux sociaux : on paie un article qui n'arrive jamais.
\end{itemize}
La règle d'or : ne jamais communiquer son code secret, ne jamais envoyer d'argent à un inconnu, vérifier toute annonce avant de payer.
\end{exemplebox}

\courssub{Les dégradations et destructions de biens}

\begin{definition}[Dégradation]
La \cle{dégradation} est la destruction, la détérioration ou l'endommagement volontaire d'un bien appartenant à autrui, qu'il soit \emph{privé} (maison, véhicule, boutique) ou \emph{public} (école, hôpital, pont, mobilier urbain).
\end{definition}

Contrairement au vol, l'auteur de la dégradation ne cherche pas à s'approprier le bien : il le détruit par colère, par vengeance ou par vandalisme. Les cas fréquents au Cameroun sont :

\begin{itemize}
\item le \emph{vandalisme lors de manifestations} : bris de vitrines, incendie de véhicules, saccage de bâtiments administratifs ;
\item la \emph{destruction d'écoles et de marchés}, notamment dans les régions en crise, ce qui prive des milliers d'enfants de leur droit à l'éducation ;
\item les dégradations du mobilier urbain : lampadaires arrachés, panneaux de signalisation détruits.
\end{itemize}

\begin{attention}[Un bien public appartient à tous]
Dégrader une école, un hôpital ou une borne-fontaine, c'est s'attaquer au patrimoine de toute la communauté, donc aussi au sien. Le contribuable paiera deux fois : une fois pour construire, une fois pour réparer. Le citoyen responsable protège les biens publics et dénonce ceux qui les détruisent.
\end{attention}

\courssub{La spoliation foncière et les conflits fonciers}

Au Cameroun, la terre a une valeur économique, sociale et symbolique considérable. C'est pourquoi les atteintes au patrimoine immobilier sont particulièrement graves et fréquentes.

\begin{definition}[Spoliation foncière]
La \cle{spoliation foncière} est le fait de s'approprier frauduleusement la terre d'autrui, par la force, la ruse ou l'abus de pouvoir : fausses ventes, faux titres fonciers, occupation illégale, expropriation abusive sans indemnisation juste.
\end{definition}

Les formes les plus courantes sont :
\begin{itemize}
\item la \emph{double vente} : un même terrain est vendu à deux acheteurs différents, souvent avec de faux documents ;
\item la vente d'un terrain familial par un seul membre, sans l'accord des autres héritiers ;
\item l'\emph{expropriation abusive} : une personne est chassée de sa terre sans procédure régulière ni indemnisation équitable ;
\item l'occupation illégale des terres d'autrui avec construction rapide pour créer un « fait accompli ».
\end{itemize}

\begin{exoresolu}[La double vente d'un terrain]
\textbf{Énoncé.} À Douala, M. Etoundi vend sa parcelle de 400 m\textsuperscript{2} à Mme Ngo Bell pour 3 000 000 FCFA, avec un reçu signé. Trois mois plus tard, il vend la \emph{même} parcelle à M. Kamga pour 3 500 000 FCFA, en lui remettant un second reçu. M. Kamga commence à construire. Mme Ngo Bell découvre les travaux et saisit le tribunal. Analysez la situation : quelle infraction est commise ? Que risque M. Etoundi ? Quels recours ont les deux acheteurs ?
\tcblower
\textbf{Solution.} M. Etoundi a commis une \cle{escroquerie} (il a obtenu de l'argent de M. Kamga en vendant un bien qu'il ne pouvait plus vendre, par manœuvre frauduleuse) et une \cle{spoliation foncière} à l'égard de Mme Ngo Bell. Il risque des peines de prison et des amendes prévues par le Code pénal. Devant le tribunal, le juge examinera les titres et les dates des ventes : en principe, le premier acheteur de bonne foi (Mme Ngo Bell) voit son droit reconnu. M. Kamga, lésé, peut demander l'\emph{annulation de sa vente}, la \emph{restitution} de ses 3 500 000 FCFA et des \emph{dommages et intérêts} (réparation civile). La leçon à retenir : avant tout achat de terrain, il faut vérifier le titre foncier auprès de la conservation foncière et s'assurer que le vendeur est bien le propriétaire.
\end{exoresolu}

\begin{exemplebox}[Un contentieux massif devant les tribunaux]
Les conflits fonciers représentent une part très importante des affaires portées devant les tribunaux camerounais — selon les acteurs de la justice, ils constitueraient dans certaines juridictions plus de la moitié des dossiers civils. La double vente d'un même terrain y est un cas fréquent de spoliation foncière. Ces litiges opposent souvent des membres d'une même famille, des villageois à des élites urbaines, ou des communautés à des investisseurs.
\end{exemplebox}

\begin{popfigure}
\begin{center}
\begin{tikzpicture}[scale=0.92]
% Contour très simplifié du Cameroun
\draw[thick,popdark,fill=popcream] 
  (2.0,0.0) -- (3.6,0.2) -- (4.6,1.0) -- (4.4,2.2) -- (4.8,3.2) -- (4.2,4.2)
  -- (3.4,4.6) -- (3.6,5.6) -- (3.0,6.8) -- (3.4,8.0) -- (2.6,8.6)
  -- (1.8,7.6) -- (2.2,6.4) -- (1.6,5.2) -- (1.0,4.4) -- (0.4,3.4)
  -- (0.8,2.2) -- (0.2,1.2) -- (1.0,0.4) -- cycle;
% Nord (flèche)
\draw[->,thick,popdark] (5.4,7.6) -- (5.4,8.6) node[above]{\textbf{N}};
% Zones de tension / conflits fonciers
\fill[poporange!55,draw=poporange,thick] (2.9,7.4) circle (0.75);
\node[align=center,font=\scriptsize] at (2.9,7.4) {Extrême-Nord\\(insécurité)};
\fill[popink!35,draw=popink,thick] (1.35,3.9) circle (0.7);
\node[align=center,font=\scriptsize] at (1.35,3.9) {Nord-Ouest\\Sud-Ouest};
\fill[popblue!35,draw=popblue,thick] (2.6,1.6) circle (0.55);
\node[align=center,font=\scriptsize] at (2.6,1.6) {Douala};
\fill[popblue!35,draw=popblue,thick] (1.7,2.6) circle (0.5);
\node[align=center,font=\scriptsize] at (1.7,2.6) {Yaoundé};
\fill[popgold!45,draw=popgold,thick] (3.8,3.4) circle (0.55);
\node[align=center,font=\scriptsize] at (3.8,3.4) {Est / Adamaoua\\(pression foncière)};
% Villes
\fill[popdark] (2.6,1.6) circle (0.06);
\fill[popdark] (1.7,2.6) circle (0.06);
\fill[popdark] (2.9,7.4) circle (0.06);
\node[font=\scriptsize,anchor=west] at (3.0,8.1) {Maroua};
\fill[popdark] (2.95,8.05) circle (0.06);
\node[font=\scriptsize,anchor=east] at (1.0,4.6) {Bamenda, Buea};
\end{tikzpicture}
\end{center}
\legende{Carte schématique du Cameroun : principales zones de conflits fonciers et d'atteintes au patrimoine. 1 : Extrême-Nord (insécurité liée aux groupes armés, vols de bétail). 2 : Nord-Ouest et Sud-Ouest (crise sécuritaire, destructions d'écoles et de biens). 3 : Douala et Yaoundé (vols, escroqueries, doubles ventes de terrains). 4 : Adamaoua et Est (pression foncière, conflits agriculteurs-éleveurs). Croquis simplifié, sans précision cartographique.}
\end{popfigure}

\courssub{La corruption et le détournement de deniers publics}

Il existe enfin une atteinte au patrimoine d'un genre particulier : celle qui vise le \emph{patrimoine national}, c'est-à-dire les biens et l'argent de tous les citoyens.

\begin{definition}[Corruption et détournement de deniers publics]
La \cle{corruption} est le fait de donner, de recevoir ou d'exiger un avantage (argent, cadeau, service) en échange d'un acte qui relève de la fonction d'une personne (délivrance d'un document, attribution d'un marché, note d'examen).
Le \cle{détournement de deniers publics} est le fait, pour un agent public, de s'approprier ou d'utiliser à des fins personnelles des fonds, des biens ou des matériels appartenant à l'État ou à une collectivité.
\end{definition}

Ces infractions appauvrissent la nation entière : l'argent détourné aurait pu construire des écoles, des hôpitaux, des routes. La corruption fausse aussi la concurrence et décourage l'effort : pourquoi travailler si c'est le « piston » ou l'enveloppe qui décide ?

\begin{savaistu}[La CONAC]
La \cle{CONAC} (Commission nationale anti-corruption), créée en 2006, est l'institution camerounaise chargée de la lutte contre la corruption. Elle publie chaque année un \emph{rapport sur l'état de la lutte contre la corruption au Cameroun}, recueille les dénonciations des citoyens et transmet les dossiers à la justice. D'autres institutions participent à ce combat : la Chambre des comptes de la Cour suprême, qui contrôle l'emploi des deniers publics, et le Tribunal criminel spécial, créé en 2011, qui juge les affaires de détournement de fonds publics.
\end{savaistu}

\begin{popfigure}
\begin{center}
\begin{tikzpicture}
\begin{axis}[
  ybar, bar width=0.55cm,
  width=0.95\linewidth, height=6.2cm,
  ymin=0, ymax=40,
  ylabel={Fréquence indicative (\%)},
  symbolic x coords={Vols, Escroqueries, Dégradations, Conflits fonciers, Recel},
  xtick=data,
  x tick label style={rotate=18, anchor=east, font=\small},
  nodes near coords, nodes near coords style={font=\small},
  ymajorgrids, grid style={popline!40},
  axis line style={popdark},
]
\addplot[fill=popblue!70, draw=popblue] coordinates {(Vols,35) (Escroqueries,25) (Dégradations,15) (Conflits fonciers,15) (Recel,10)};
\end{axis}
\end{tikzpicture}
\end{center}
\legende{Types d'atteintes au patrimoine les plus fréquentes en milieu urbain camerounais (données indicatives, en pourcentage des plaintes enregistrées ; total : 100 \%). Les vols et les escroqueries — notamment par mobile money — arrivent en tête.}
\end{popfigure}

\courssub{Les réponses de la société : plainte, enquête, restitution, réparation}

Face aux atteintes au patrimoine, la victime n'est pas démunie. La loi lui offre une démarche claire.

\begin{methode}[Que faire en cas d'atteinte au patrimoine ?]
\begin{enumerate}
\item \textbf{Conserver les preuves} : reçus, messages, captures d'écran, témoins, photos des dégâts.
\item \textbf{Porter plainte} au commissariat de police ou à la brigade de gendarmerie ; pour les conflits fonciers, saisir aussi le sous-préfet ou la conservation foncière.
\item \textbf{L'enquête} est menée par les forces de l'ordre sous le contrôle du procureur de la République.
\item \textbf{Le procès} : le tribunal juge l'auteur et peut prononcer des peines de prison et des amendes.
\item \textbf{La réparation} : la victime peut se constituer \emph{partie civile} pour obtenir la \cle{restitution} du bien (s'il existe encore) ou des \cle{dommages et intérêts} (réparation en argent du préjudice subi).
\end{enumerate}
\end{methode}

\begin{aretenir}[L'essentiel]
Les atteintes au patrimoine visent les biens meubles et immeubles des personnes physiques et morales. On distingue le \cle{vol} (prendre), l'\cle{escroquerie} (faire donner par la ruse), l'\cle{abus de confiance} (garder indûment), la \cle{dégradation} (détruire), la \cle{spoliation foncière} (usurper une terre) et le \cle{détournement de deniers publics} (voler la communauté nationale). Acheter un bien volé est un délit (recel). La victime doit porter plainte ; l'auteur risque la prison, l'amende et l'obligation de restituer ou de réparer. La CONAC et les tribunaux luttent contre la corruption, atteinte au patrimoine de tous.
\end{aretenir}

\begin{exercice}[À vous de jouer]
\begin{enumerate}
\item Un jeune achète un smartphone à très bas prix dans la rue, sans facture, alors que le vendeur lui dit qu'il « vient de loin ». Quel délit ce jeune risque-t-il de commettre ? Justifiez.
\item Votre oncle reçoit un message lui annonçant qu'il a gagné 500 000 FCFA à un concours auquel il n'a jamais participé, à condition d'envoyer 10 000 FCFA de frais. Identifiez l'infraction et proposez trois réflexes de protection.
\item Lors d'une manifestation, des jeunes incendient le marché communal. Qui est victime de cette atteinte au patrimoine ? Expliquez pourquoi on dit que le bien public appartient à tous.
\end{enumerate}
\end{exercice}

\begin{corrige}[Exercice « À vous de jouer »]
\begin{enumerate}
\item Il risque le délit de \cle{recel} : acheter un bien que l'on sait (ou que l'on devrait soupçonner) volé est puni par la loi, même sans avoir participé au vol. L'absence de facture et le prix anormalement bas sont des indices de l'origine frauduleuse.
\item Il s'agit d'une tentative d'\cle{escroquerie} en ligne (faux concours). Réflexes : ne jamais envoyer d'argent pour recevoir un gain ; ne jamais communiquer ses codes secrets ; vérifier l'existence réelle du concours et signaler le numéro à l'opérateur ou à la police.
\item La victime est la \emph{communauté tout entière} : le marché communal est un bien public financé par les impôts de tous. Le reconstruire coûtera encore de l'argent public, et les commerçants perdent leur outil de travail. Dégrader un bien public, c'est donc s'appauvrir soi-même.
\end{enumerate}
\end{corrige}

\coursec{L'extrémisme violent}

Nous venons d'étudier les atteintes au patrimoine : vols, destructions, pillages. Mais il existe une forme de violence encore plus grave, qui combine à la fois les atteintes à l'intégrité physique, à l'intégrité morale et au patrimoine : c'est la violence commise au nom d'une \cle{idéologie}. Lorsqu'un groupe décide d'imposer ses convictions par les armes, les attentats ou la terreur, on parle d'\cle{extrémisme violent}. Cette forme de violence touche directement le Cameroun, dans l'Extrême-Nord comme dans les régions du Nord-Ouest et du Sud-Ouest. Il est donc essentiel de la comprendre pour savoir la prévenir et la combattre.

\courssub{Définition de l'extrémisme violent}

Commençons par une situation concrète. Dans un village, un jeune homme affirme : « Ceux qui ne pensent pas comme nous sont des ennemis. Il faut les chasser ou les punir. » Ce discours ne relève plus d'une simple opinion : il justifie la violence pour imposer une conviction. C'est exactement le mécanisme de l'extrémisme violent.

\begin{definition}[Extrémisme violent]
L'\cle{extrémisme violent} est une \textbf{idéologie radicale qui légitime l'usage de la violence} (attentats, enlèvements, destructions, meurtres) pour imposer des convictions \textbf{religieuses, politiques ou identitaires} à l'ensemble de la société. Il rejette le dialogue, la loi et le pluralisme, et considère que la fin (imposer ses idées) justifie tous les moyens, y compris les plus cruels.
\end{definition}

\begin{definition}[Radicalisation]
La \cle{radicalisation} est le \textbf{processus progressif} par lequel un individu adopte des idées de plus en plus extrêmes, se coupe de son entourage et finit par accepter l'idée que la violence est légitime. C'est un \textbf{cheminement}, pas un acte : on peut être radicalisé sans être (encore) passé à l'acte violent.
\end{definition}

\begin{aretenir}[Processus et passage à l'acte]
Il faut bien distinguer deux moments :
\begin{itemize}
\item la \textbf{radicalisation} : processus psychologique et social d'adhésion à des idées extrêmes ;
\item le \textbf{passage à l'acte violent} : le moment où l'individu commet réellement une violence (attentat, agression, destruction).
\end{itemize}
Tous les radicalisés ne passent pas à l'acte, mais tout passage à l'acte est précédé d'une radicalisation. Agir \textbf{tôt}, pendant le processus de radicalisation, est donc la clé de la prévention.
\end{aretenir}

\courssub{Les facteurs de vulnérabilité : comment bascule-t-on ?}

Pourquoi un jeune, souvent âgé de 15 à 25 ans, adhère-t-il à un discours violent ? Les spécialistes ont identifié plusieurs \cle{facteurs de vulnérabilité} qui se combinent :

\begin{itemize}
\item la \textbf{pauvreté et le chômage} : un jeune sans emploi ni perspective est une cible facile pour les recruteurs qui promettent argent et statut ;
\item l'\textbf{exclusion sociale} : sentiment d'injustice, d'humiliation, d'abandon par la société ou la famille ;
\item la \textbf{manipulation} : des recruteurs expérimentés exploitent les croyances religieuses ou les frustrations politiques, en déformant les textes et l'histoire ;
\item les \textbf{réseaux sociaux} : vidéos de propagande, groupes fermés, discours de haine diffusés massivement et sans contradicteur ;
\item la \textbf{rupture familiale ou scolaire} : l'isolement rend plus vulnérable à l'embrigadement.
\end{itemize}

\begin{attention}[Ne pas confondre religion et extrémisme]
Il est \textbf{faux et dangereux} d'assimiler une religion ou une communauté entière à l'extrémisme violent. Les extrémistes sont une \textbf{toute petite minorité} qui \emph{manipule} et \emph{déforme} les croyances pour justifier leurs crimes. Les premières victimes de Boko Haram, par exemple, sont des musulmans de l'Extrême-Nord qui refusent cette idéologie. Stigmatiser une communauté, c'est faire exactement ce que les extrémistes attendent : diviser la société.
\end{attention}

\begin{methode}[Repérer un discours radical en ligne]
Face à un message, une vidéo ou un groupe sur les réseaux sociaux, pose-toi quatre questions :
\begin{enumerate}
\item \textbf{Appel à la haine} : le message désigne-t-il un groupe de personnes comme des « ennemis » à combattre ou à éliminer ?
\item \textbf{Rejet du dialogue} : le message refuse-t-il toute discussion, tout compromis, toute opinion différente ?
\item \textbf{Glorification de la violence} : la violence y est-elle présentée comme héroïque, légitime ou récompensée ?
\item \textbf{Isolement du groupe} : le message demande-t-il de rompre avec la famille, les amis, l'école, et de ne faire confiance qu'au groupe ?
\end{enumerate}
Si plusieurs réponses sont « oui », tu es probablement face à un discours radical : ne le partage pas, signale-le, et parles-en à un adulte de confiance.
\end{methode}

\courssub{L'extrémisme violent au Cameroun : deux réalités}

\textbf{1. Boko Haram dans l'Extrême-Nord.} Depuis \textbf{2014}, la secte terroriste Boko Haram, née au Nigeria voisin, étend ses attaques à la région de l'Extrême-Nord du Cameroun (départements du Logone-et-Chari, du Mayo-Sava, du Mayo-Tsanaga). Ses exactions sont multiples : \textbf{attentats-suicides} dans les marchés et les mosquées, \textbf{enlèvements} de civils (femmes, enfants, chefs traditionnels), \textbf{attaques de villages} avec pillages, incendies et massacres, pose de mines. Des localités comme Fotokol, Amchidé ou Kerawa ont été durement touchées.

\textbf{2. La crise dans le Nord-Ouest et le Sud-Ouest.} Depuis 2016-2017, les régions du Nord-Ouest et du Sud-Ouest connaissent une crise née de revendications corporatistes (avocats, enseignants) qui ont dégénéré. Une partie des contestataires s'est \textbf{radicalisée} et a pris les armes pour réclamer la sécession : on parle de \textbf{violences séparatistes}. Les groupes armés commettent des enlèvements contre rançon, des assassinats, des incendies d'écoles et de marchés, et imposent des « villes mortes » (\emph{ghost towns}) qui paralysent la vie économique.

\begin{popfigure}
\begin{center}
\begin{tikzpicture}[scale=1.0, font=\small]
% Contour très simplifié du Cameroun
\draw[thick, popdark, fill=popcream] (2.0,0) -- (4.2,0.3) -- (4.6,1.6) -- (4.3,3.0) -- (4.6,4.2) -- (3.8,5.2) -- (2.8,5.6) -- (2.4,6.6) -- (2.0,7.6) -- (1.4,8.4) -- (0.8,8.2) -- (1.0,7.0) -- (1.2,5.8) -- (0.9,4.6) -- (0.4,3.4) -- (0.2,2.2) -- (0.8,1.0) -- cycle;
% Zone Boko Haram (Extrême-Nord)
\draw[fill=poporange!55, draw=none] (1.0,7.0) -- (0.8,8.2) -- (1.4,8.4) -- (2.0,7.6) -- (2.4,6.6) -- (1.6,6.4) -- cycle;
% Zone crise anglophone (Nord-Ouest / Sud-Ouest)
\draw[fill=poppink!70, draw=none] (0.2,2.2) -- (0.8,1.0) -- (2.0,0) -- (2.4,1.2) -- (1.8,2.4) -- (1.2,3.0) -- (0.4,3.4) -- cycle;
% Villes
\fill[popblue] (2.6,1.4) circle (0.07) node[right, popdark]{Yaoundé};
\fill[popblue] (1.2,0.7) circle (0.07) node[right, popdark]{Douala};
\fill[poporange] (1.5,7.6) circle (0.07) node[right, popdark]{Maroua};
\fill[poppink] (0.9,2.4) circle (0.07) node[right=1pt, popdark]{Bamenda};
\fill[poppink] (1.1,1.3) circle (0.07) node[left=1pt, popdark]{Buéa};
% Flèche du nord
\draw[-{Stealth}, thick, popdark] (5.2,7.4) -- (5.2,8.2) node[above]{N};
\end{tikzpicture}
\end{center}
\legende{Carte schématique du Cameroun. En orange : l'Extrême-Nord, zone affectée par les attaques de Boko Haram depuis 2014. En rose : les régions du Nord-Ouest et du Sud-Ouest, touchées par la crise séparatiste. Croquis simplifié, sans prétention de précision cartographique.}
\end{popfigure}

\courssub{Les conséquences de l'extrémisme violent}

L'extrémisme violent détruit des vies, mais aussi toute la vie sociale d'une région. Ses conséquences sont lourdes :

\begin{itemize}
\item des \textbf{déplacés internes} : des centaines de milliers de personnes fuient leurs villages pour se réfugier dans des villes ou des camps, dans des conditions très précaires ;
\item des \textbf{réfugiés} : des Camerounais passent les frontières, et le Cameroun accueille aussi des réfugiés nigérians fuyant Boko Haram (camp de Minawao, près de Mokolo) ;
\item des \textbf{écoles fermées} : dans le Nord-Ouest et le Sud-Ouest, des milliers d'écoles ont été fermées ou incendiées, privant des centaines de milliers d'enfants d'éducation ;
\item une \textbf{économie locale détruite} : marchés brûlés, champs abandonnés, « villes mortes », commerces fermés, insécurité sur les routes ;
\item des \textbf{traumatismes psychologiques} : orphelins, enfants enrôlés, femmes victimes de violences, populations vivant dans la peur permanente.
\end{itemize}

\begin{popfigure}
\begin{center}
\begin{tikzpicture}
\begin{axis}[
width=0.85\linewidth, height=6cm,
xlabel={Année}, ylabel={Déplacés internes (en milliers)},
xmin=2013.5, xmax=2024.5, ymin=0, ymax=1100,
xtick={2014,2016,2018,2020,2022,2024},
ytick={0,200,400,600,800,1000},
grid=major, grid style={popline!40},
legend style={at={(0.03,0.97)}, anchor=north west, font=\small},
axis line style={popdark},
]
\addplot[popcurve, mark=*, mark options={fill=poporange}] coordinates {
(2014,40) (2015,95) (2016,200) (2017,240) (2018,270) (2019,280) (2020,320) (2021,580) (2022,1000) (2023,1050) (2024,1000)
};
\legend{Ordre de grandeur (OCHA)}
\end{axis}
\end{tikzpicture}
\end{center}
\legende{Évolution indicative du nombre de déplacés internes liés aux conflits au Cameroun (2014-2024), en milliers de personnes. Les valeurs sont des ordres de grandeur inspirés des rapports OCHA : la forte hausse à partir de 2021 reflète l'aggravation de la crise du Nord-Ouest et du Sud-Ouest, qui s'ajoute à celle de l'Extrême-Nord.}
\end{popfigure}

\courssub{Les réponses de l'État}

Face à ces violences, l'État camerounais a mis en place plusieurs réponses complémentaires :

\begin{itemize}
\item l'\textbf{action militaire et sécuritaire} : opérations des forces de défense dans l'Extrême-Nord (dont la force multinationale mixte avec les pays voisins), sécurisation des villages, des axes routiers et des frontières ;
\item le \textbf{Comité national de désarmement, démobilisation et réintégration (DDR)} : créé en 2018, il accueille les ex-combattants qui déposent les armes, dans des centres (Bamenda, Buéa, Mora) où ils reçoivent une formation et une aide à la réinsertion sociale et professionnelle ;
\item la \textbf{reconstruction} : réouverture des écoles, reconstruction des infrastructures détruites, plans de relance économique pour les régions touchées, aide humanitaire aux déplacés.
\end{itemize}

\begin{savaistu}[Les comités de vigilance de l'Extrême-Nord]
Dans l'Extrême-Nord, les populations elles-mêmes se sont organisées en \textbf{comités de vigilance locaux}. Ces groupes de villageois coopèrent avec les forces de défense et de sécurité : ils surveillent les mouvements suspects, signalent les intrus, protègent les puits, les marchés et les écoles. Grâce à cette collaboration entre citoyens et forces de l'ordre, de nombreuses attaques ont été déjouées. C'est la preuve que la \textbf{sécurité est l'affaire de tous}, pas seulement de l'armée.
\end{savaistu}

\courssub{Les réponses citoyennes et préventives}

L'État ne peut pas tout. Chaque citoyen, y compris chaque élève, a un rôle à jouer dans la prévention de l'extrémisme violent :

\begin{itemize}
\item la \textbf{vigilance communautaire} : signaler aux autorités les comportements suspects, les recruteurs, les discours dangereux, sans stigmatiser personne ;
\item le \textbf{dialogue interreligieux et interculturel} : chrétiens, musulmans et adeptes d'autres croyances se rencontrent, se connaissent et refusent ensemble les discours de division ;
\item l'\textbf{éducation à la paix} : à l'école, dans les clubs et les associations, apprendre à régler les conflits par le dialogue, le respect et la négociation ;
\item la \textbf{dénonciation des discours de haine} : ne pas partager les messages haineux sur les réseaux sociaux, les signaler, et y répondre par des messages de tolérance et de fraternité.
\end{itemize}

\begin{exoresolu}[Analyser une situation]
Énoncé : Dans un quartier de Maroua, un jeune de 19 ans, au chômage depuis deux ans, passe ses journées sur son téléphone dans un groupe de discussion fermé. Il répète à ses amis : « Ceux qui ne sont pas avec nous sont des traitres. Un jour, ils paieront. » Il a quitté l'école et ne parle presque plus à sa famille. Un de ses amis, inquiet, vient te demander conseil. Analyse la situation et propose une conduite à tenir.
\tcblower
\textbf{Solution détaillée.}
\textbf{1. Identification des signaux.} Ce jeune cumule plusieurs facteurs de vulnérabilité et signes de radicalisation : chômage et exclusion sociale, isolement (rupture avec l'école et la famille), fréquentation d'un groupe fermé en ligne, discours qui désigne des « ennemis » et évoque une vengeance violente.
\textbf{2. Diagnostic.} Il s'agit probablement d'un début de \cle{radicalisation} (processus), sans passage à l'acte pour l'instant. C'est précisément le moment où la prévention est encore possible.
\textbf{3. Conduite à tenir.} L'ami ne doit ni ignorer la situation, ni accuser publiquement le jeune, ni diffuser son cas sur les réseaux sociaux. Il doit : (a) maintenir le lien et le dialogue avec lui, sans valider ses propos ; (b) en parler à un adulte de confiance (parent, enseignant, chef de quartier, autorité religieuse modérée) ; (c) si le risque semble grave, signaler discrètement la situation aux autorités compétentes.
\textbf{4. Conclusion.} Face à la radicalisation, la bonne réponse citoyenne combine \textbf{vigilance} (ne pas laisser faire) et \textbf{humanité} (ne pas rejeter la personne, mais l'aider à sortir du processus). Dénoncer un discours de haine, ce n'est pas trahir : c'est protéger la personne et la communauté.
\end{exoresolu}

\begin{exercice}[À toi de jouer]
1. Définis l'extrémisme violent et la radicalisation, en montrant la différence entre les deux notions.
2. Cite quatre facteurs de vulnérabilité qui favorisent la radicalisation des jeunes.
3. Explique pourquoi il est injuste et contre-productif d'assimiler une religion entière à l'extrémisme violent.
4. Décris deux réponses de l'État camerounais et deux réponses citoyennes face à l'extrémisme violent.
\end{exercice}

\begin{corrige}[Exercice]
1. L'extrémisme violent est une idéologie radicale qui légitime l'usage de la violence pour imposer des convictions religieuses, politiques ou identitaires. La radicalisation est le processus progressif d'adhésion à ces idées extrêmes ; elle précède le passage à l'acte, qui n'est pas systématique.
2. La pauvreté et le chômage, l'exclusion sociale, la manipulation par des recruteurs, la propagande sur les réseaux sociaux (ainsi que les ruptures familiales et scolaires).
3. Parce que les extrémistes sont une minorité qui déforme les croyances pour justifier ses crimes, et que les membres de cette religion sont souvent leurs premières victimes. Stigmatiser une communauté divise la société, ce qui sert le projet des extrémistes.
4. Réponses de l'État : action militaire et sécurisation ; Comité DDR pour les ex-combattants ; reconstruction et aide humanitaire. Réponses citoyennes : vigilance communautaire (comités de vigilance), dialogue interreligieux, éducation à la paix, dénonciation des discours de haine.
\end{corrige}

\begin{aretenir}[L'essentiel]
L'extrémisme violent est une idéologie qui justifie la violence pour imposer ses convictions. Il résulte d'un processus de \textbf{radicalisation}, favorisé par la pauvreté, l'exclusion, la manipulation et les réseaux sociaux. Au Cameroun, il frappe l'Extrême-Nord (Boko Haram) et les régions du Nord-Ouest et du Sud-Ouest (crise séparatiste), avec des conséquences dramatiques : déplacés, écoles fermées, économie détruite. Le combattre exige à la fois l'action de l'État (armée, DDR, reconstruction) et l'engagement de chaque citoyen : vigilance, dialogue, éducation à la paix et refus des discours de haine.
\end{aretenir}

\coursec{Prévenir et combattre les violences : le rôle du citoyen}

Dans la section précédente, nous avons étudié l'extrémisme violent, forme radicale de violence qui détruit les vies et déchire le tissu social. Face à toutes ces violences — atteintes à l'intégrité physique, morale, au patrimoine, extrémisme violent — une question se pose naturellement : que peut-on faire ? La réponse est claire : la lutte contre les violences n'est pas l'affaire des seuls policiers ou des juges. Elle commence par chacun de nous, dans la famille, à l'atelier, au lycée, dans le quartier. Cette dernière section montre comment le citoyen, entouré de nombreux acteurs, peut \cle{prévenir} les violences et y réagir efficacement.

\courssub{La prévention primaire : éduquer aux valeurs et au dialogue}

\begin{definition}[Prévention primaire]
La \cle{prévention primaire} désigne l'ensemble des actions menées \emph{en amont}, c'est-à-dire \textbf{avant} que la violence n'éclate, pour en éliminer les causes. Elle repose sur l'éducation aux valeurs, le dialogue et la gestion pacifique des conflits.
\end{definition}

L'idée est simple : il est plus facile et moins coûteux d'empêcher un incendie que de l'éteindre. De même, il vaut mieux former des jeunes au respect mutuel que de réparer des vies brisées. La prévention primaire repose sur trois piliers.

\begin{itemize}
\item \textbf{L'éducation aux valeurs} : le \cle{respect} de la dignité de chaque personne, la \cle{tolérance} (accepter la différence d'opinion, de religion, d'ethnie), la \cle{non-violence} (refuser de régler ses problèmes par la force physique ou verbale).
\item \textbf{Le dialogue} : parler, écouter, négocier. Le dialogue suppose que chacun accepte d'entendre le point de vue de l'autre, même quand on n'est pas d'accord.
\item \textbf{La gestion pacifique des conflits} : un conflit n'est pas forcément une violence. C'est un désaccord qui peut être réglé par la \cle{médiation} (intervention d'un tiers neutre), la \cle{négociation} ou l'\cle{arbitrage}.
\end{itemize}

\begin{exemplebox}[Un conflit réglé sans violence]
Dans un atelier de mécanique à Douala, deux apprentis se disputent l'usage d'une machine. Au lieu d'en venir aux mains, le chef d'atelier les réunit : chacun expose ses besoins, puis ils conviennent d'un planning d'utilisation. Le conflit est réglé par la médiation, sans violence, et la relation de travail est préservée.
\end{exemplebox}

\courssub{Les acteurs de la prévention : de la famille aux organisations internationales}

La prévention des violences fonctionne comme des cercles concentriques : chaque acteur, du plus proche au plus lointain, a un rôle complémentaire.

\begin{popfigure}
\begin{center}
\begin{tikzpicture}[scale=0.9]
\fill[popblue!15] (0,0) circle (4.6);
\fill[popgreen!20] (0,0) circle (3.8);
\fill[popgold!25] (0,0) circle (3.0);
\fill[poporange!30] (0,0) circle (2.2);
\fill[poppink!35] (0,0) circle (1.4);
\fill[popblue!50] (0,0) circle (0.7);
\draw[popdark,thick] (0,0) circle (4.6);
\draw[popdark,thick] (0,0) circle (3.8);
\draw[popdark,thick] (0,0) circle (3.0);
\draw[popdark,thick] (0,0) circle (2.2);
\draw[popdark,thick] (0,0) circle (1.4);
\draw[popdark,thick] (0,0) circle (0.7);
\node[font=\small\bfseries,text=white] at (0,0) {Individu};
\node[font=\small\bfseries] at (0,1.05) {Famille};
\node[font=\small\bfseries] at (0,1.8) {\'Ecole};
\node[font=\small\bfseries] at (0,2.6) {Communaut\'e};
\node[font=\small\bfseries] at (0,3.4) {\'Etat};
\node[font=\small\bfseries] at (0,4.2) {Organisations internationales};
\end{tikzpicture}
\end{center}
\legende{Les cercles concentriques de la prévention des violences : chaque niveau protège et éduque le niveau inférieur. 1 : l'individu (valeurs personnelles) ; 2 : la famille (première éducation) ; 3 : l'école (citoyenneté) ; 4 : la communauté (chefs traditionnels et religieux, médias) ; 5 : l'État (lois, justice, sécurité) ; 6 : les organisations internationales (ONU, ONG).}
\end{popfigure}

\courssub{Le rôle de la famille, de l'école, des chefs et des médias}

\textbf{La famille} est la première école de la vie en société. C'est elle qui transmet les premières valeurs : le respect des parents, des aînés, des autres enfants. Une famille où règnent les coups et les insultes fabrique souvent des adultes violents ; une famille où l'on dialogue forme des citoyens pacifiques.

\textbf{L'école} prolonge cette éducation. Par les cours d'ECM, les clubs de citoyenneté, les règlements intérieurs, elle apprend la vie en communauté, le respect des différences et le règlement pacifique des conflits. Le lycée technique, où filles et garçons de toutes origines apprennent ensemble un métier, est un lieu privilégié de la culture de la paix.

\textbf{Les chefs traditionnels et religieux} jouissent d'une grande autorité morale dans nos quartiers et villages. Ils règlent les conflits de voisinage, prêchent la paix, rappellent les valeurs de solidarité et de pardon. Leur parole, écoutée par tous, est un puissant instrument de prévention.

\textbf{Les médias} (radio, télévision, réseaux sociaux) façonnent les opinions. Ils peuvent propager la haine et les fausses informations, mais ils peuvent aussi diffuser des messages de paix, dénoncer les violences et éduquer le public. Un journaliste responsable vérifie ses informations ; un internaute responsable ne partage pas les appels à la haine.

\courssub{Le rôle de l'État, de la société civile et des organisations internationales}

\begin{definition}[Société civile]
La \cle{société civile} désigne l'ensemble des organisations non gouvernementales créées librement par les citoyens : associations, syndicats, ONG, mouvements de jeunes et de femmes.
\end{definition}

\textbf{L'État} agit à quatre niveaux :
\begin{itemize}
\item la \textbf{législation} : le Parlement vote des lois qui punissent les violences (le Code pénal camerounais réprime les coups et blessures, les viols, les vols, les destructions de biens) ;
\item la \textbf{justice} : les tribunaux jugent les auteurs de violences et réparent le préjudice des victimes ;
\item les \textbf{forces de sécurité} (police, gendarmerie) : elles protègent les personnes et les biens, interpellent les agresseurs ;
\item les \textbf{politiques sociales} : la lutte contre la pauvreté et le chômage des jeunes supprime un terreau de la violence et de l'extrémisme violent.
\end{itemize}

\textbf{Les organisations de la société civile} sensibilisent les populations, forment des médiateurs communautaires, accueillent et accompagnent les victimes (centres d'écoute, assistance juridique et psychologique).

\textbf{Les organisations internationales}, comme l'\cle{ONU} (avec l'UNICEF pour les enfants, l'UNESCO pour l'éducation à la paix) et les ONG de défense des droits humains (Amnesty International, Human Rights Watch), soutiennent les États, dénoncent les violences dans le monde et protègent les victimes des conflits.

\begin{exemplebox}[Les clubs de citoyenneté dans les lycées camerounais]
Dans de nombreux lycées techniques du Cameroun, des \cle{clubs de citoyenneté} et des comités de bien-être réunissent chaque semaine des dizaines d'élèves. Ils y organisent des débats sur la paix, des campagnes contre les violences en milieu scolaire, des journées de solidarité. Un club de 30 élèves formés peut sensibiliser à son tour plus de 500 camarades : c'est la prévention par les pairs, très efficace car le message passe mieux entre jeunes.
\end{exemplebox}

\begin{savaistu}[Le 2 octobre, Journée internationale de la non-violence]
L'Organisation des Nations unies célèbre chaque année le \textbf{2 octobre} la Journée internationale de la non-violence. Cette date n'est pas un hasard : c'est l'anniversaire de la naissance du \cle{Mahatma Gandhi} (1869-1948), leader de l'indépendance de l'Inde, qui a montré au monde qu'on peut vaincre un empire par la résistance pacifique, sans verser le sang. Son exemple a inspiré de nombreux mouvements de paix, y compris en Afrique.
\end{savaistu}

\courssub{Les conduites à tenir face à une violence}

Être témoin d'une violence — une bagarre au lycée, un vol dans le quartier, des menaces sur les réseaux sociaux — nous place devant un choix : détourner les yeux ou agir. Le citoyen responsable adopte cinq conduites.

\begin{methode}[Réagir face à une violence : les cinq conduites du citoyen]
\begin{enumerate}
\item \textbf{Protéger} : mettre la victime (et soi-même) à l'abri du danger immédiat, sans se mettre en péril. Ne jamais intervenir seul face à un groupe armé.
\item \textbf{Alerter} : appeler les secours et les autorités compétentes (police : 117 ; gendarmerie : 113 ; sapeurs-pompiers : 118), prévenir un adulte, un enseignant, un chef.
\item \textbf{Témoigner} : rapporter fidèlement ce que l'on a vu, sans exagérer ni minimiser. Le témoignage aide la justice à établir la vérité.
\item \textbf{Dénoncer} : signaler les faits aux autorités ou aux structures compétentes. Dénoncer n'est pas un acte de vengeance : c'est protéger la communauté et empêcher que la violence ne recommence.
\item \textbf{Accompagner la victime} : l'écouter, la rassurer, l'orienter vers un médecin, un psychologue, un centre d'écoute ou une association d'aide aux victimes. La victime isolée se replie ; la victime accompagnée se reconstruit.
\end{enumerate}
\end{methode}

\begin{attention}[La vengeance et l'autojustice sont aussi des violences]
Se faire justice soi-même est interdit et puni par la loi. Les \cle{lapidations} et exécutions sommaires de voleurs présumés par des foules (ce qu'on appelle parfois « justice populaire ») sont des \textbf{crimes} : ceux qui y participent peuvent être poursuivis pour coups et blessures, voire pour meurtre. De même, la vengeance entretient un cycle sans fin de violences. Au Cameroun comme dans tout État de droit, \textbf{seule la justice de l'État} peut juger et punir. Le citoyen alerte et dénonce ; il ne punit pas lui-même.
\end{attention}

\begin{exoresolu}[Que faire face à une rumeur violente sur les réseaux sociaux ?]
Énoncé : Dans le groupe WhatsApp de ta classe, un message accuse un élève d'un vol et appelle à « lui régler son compte » demain à la sortie du lycée. Plusieurs camarades partagent le message. Quelle doit être ta conduite de citoyen ? Justifie ta réponse.
\tcblower
Solution détaillée :
\begin{enumerate}
\item \textbf{Ne pas relayer} le message : partager une accusation non vérifiée, c'est participer à une atteinte à l'intégrité morale (diffamation) et préparer une atteinte physique.
\item \textbf{Protéger} : prévenir discrètement l'élève accusé pour qu'il évite le piège et se mette en sécurité.
\item \textbf{Alerter} : informer immédiatement un enseignant, le conseiller d'orientation ou le proviseur, qui peuvent prévenir la police si nécessaire.
\item \textbf{Témoigner} : conserver une capture du message comme preuve et rapporter les faits avec exactitude.
\item \textbf{Accompagner} : soutenir l'élève accusé, victime de la rumeur, et refuser de l'isoler.
\end{enumerate}
Conclusion : en agissant ainsi, tu brises la chaîne de la violence, tu protèges un innocent présumé (toute personne est présumée innocente tant qu'elle n'est pas jugée) et tu laisses la justice, seule compétente, établir la vérité.
\end{exoresolu}

\courssub{Synthèse : construire une culture de paix au quotidien}

\begin{definition}[Culture de paix]
La \cle{culture de paix} est un ensemble de valeurs, d'attitudes et de comportements qui rejettent la violence et privilégient le dialogue, le respect et la solidarité dans la vie quotidienne.
\end{definition}

Construire une culture de paix ne demande pas de grands discours : cela commence par des gestes simples — saluer son voisin, respecter la file d'attente, refuser les insultes, aider un camarade en difficulté, vérifier une information avant de la partager. Chaque citoyen pacifique est une pierre de l'édifice de la paix nationale.

\begin{methode}[Rédiger une conclusion argumentée (épreuve d'ECM)]
\begin{enumerate}
\item \textbf{Rappeler le problème} : reformule en une phrase la question posée (ex. : « Comment lutter contre les violences dans notre société ? »).
\item \textbf{Prendre position} : annonce clairement ton point de vue (ex. : « La lutte contre les violences est l'affaire de tous, et non de l'État seul »).
\item \textbf{Justifier} : appuie ta position avec deux ou trois arguments et des exemples précis (famille, école, justice, ONG).
\item \textbf{Proposer des solutions} : termine par des propositions concrètes et réalistes (clubs de citoyenneté, médiation, éducation aux valeurs, dénonciation des violences).
\end{enumerate}
\end{methode}

\begin{aretenir}[L'essentiel de la section]
\begin{itemize}
\item La \cle{prévention primaire} agit avant la violence : éducation aux valeurs, dialogue, gestion pacifique des conflits.
\item Les acteurs de la prévention forment des cercles concentriques : individu, famille, école, communauté (chefs, médias), État, organisations internationales.
\item Face à une violence, le citoyen \textbf{protège, alerte, témoigne, dénonce, accompagne}.
\item La vengeance et l'autojustice sont des violences punies par la loi : seule la justice de l'État juge et punit.
\item La paix se construit chaque jour, par les valeurs et les gestes de chacun.
\end{itemize}
\end{aretenir}

\begin{exercice}[Pour vérifier tes acquis]
\begin{enumerate}
\item Définis la prévention primaire et donne deux exemples d'actions de prévention primaire dans ton lycée.
\item Présente le rôle de trois acteurs de la prévention des violences au Cameroun.
\item Ton quartier organise une campagne contre l'autojustice. Rédige un court message (5 à 8 lignes) pour convaincre les jeunes de ne jamais participer à une lapidation.
\end{enumerate}
\end{exercice}

\begin{corrige}[Corrigé de l'exercice]
\begin{enumerate}
\item La prévention primaire est l'ensemble des actions menées avant l'éclatement de la violence pour en supprimer les causes. Exemples : les cours d'ECM sur le respect et la tolérance ; les débats et médiations organisés par le club de citoyenneté ; les campagnes d'affichage contre les bagarres.
\item La famille transmet les premières valeurs de respect et de dialogue. L'État légifère, fait juger les auteurs de violences par les tribunaux et assure la sécurité par la police et la gendarmerie. Les ONG et la société civile sensibilisent les populations et accompagnent les victimes (écoute, assistance juridique).
\item Message attendu (exemple) : « Jeune frère, jeune sœur, ne touche jamais à un présumé voleur, même pris sur le fait. La lapidation est un meurtre : tu peux tuer un innocent et finir en prison. Au Cameroun, seule la justice juge et punit. Si tu vois un vol, alerte la police au 117 ou la gendarmerie au 113. Refuse la foule, protège la victime, témoigne. Être fort, ce n'est pas frapper : c'est respecter la loi et la vie. »
\end{enumerate}
\end{corrige}

\newpage

\coursec{Activité d'intégration}

\coursec{Leçon N04 : Les formes de violences}

\courssub{Objectifs de l'activité}

Cette activité d'intégration mobilise l'ensemble des savoirs de la leçon sur les formes de violences. À l'issue du travail, l'élève doit être capable de :

\begin{itemize}
\item qualifier juridiquement et moralement des faits de violence observés dans la vie courante (atteintes à l'intégrité physique, morale, au patrimoine, extrémisme violent) ;
\item appliquer les notions de l'unité à des situations concrètes de la vie courante : agression, cyberharcèlement, double vente d'un terrain, recrutement radical, destruction de biens publics ;
\item construire une carte schématique localisant les principales formes de violences au Cameroun ;
\item rédiger et justifier une démarche argumentée : un plan d'action citoyen de prévention des violences pour son lycée ou son quartier ;
\item développer l'esprit de collaboration, l'argumentation orale et le sens de l'écoute lors de la restitution et du débat collectif.
\end{itemize}

\courssub{Situation}

\textbf{Enquête citoyenne au lycée technique de Nkol-Eton (Yaoundé).}\par

Dans le cadre de la semaine de la citoyenneté, le club ECM du lycée technique de Nkol-Eton décide de mener une \cle{enquête citoyenne} sur les violences qui touchent les jeunes et les communautés au Cameroun. Le président du club a constitué un \cle{dossier documentaire} à partir de coupures de presse camerounaises, de photographies, de données du MINPROFF (Ministère de la Promotion de la Femme et de la Famille) et d'OCHA (Bureau des Nations unies pour la coordination des affaires humanitaires), ainsi que d'extraits du Code pénal camerounais. Voici les documents remis à chaque groupe.

\textbf{Document 1 — Cinq faits divers à qualifier.}\par

\begin{enumerate}
\item \emph{Fait A} : À la sortie de l'atelier de mécanique, un élève de Terminale frappe un camarade de Première à coups de poing et lui casse le bras, à la suite d'une dispute sur un téléphone portable. La victime est hospitalisée huit jours.
\item \emph{Fait B} : Une élève de Première est victime, pendant trois mois, de messages insultants et de montages photo humiliants publiés dans un groupe WhatsApp de la classe. Elle refuse désormais de venir en cours et présente des signes de dépression.
\item \emph{Fait C} : Dans un quartier périphérique de Yaoundé, un individu vend le même terrain de \SI{500}{m^2} à deux acheteurs différents, moyennant \num{1500000} FCFA chacun, puis disparaît. Les deux familles se déchirent devant le chef de quartier.
\item \emph{Fait D} : Un jeune déscolarisé de 19 ans est contacté sur les réseaux sociaux par un recruteur qui lui promet de l'argent et un « idéal » s'il rejoint un groupe armé qui attaque des villages dans la région de l'Extrême-Nord. Le jeune commence à diffuser des vidéos de propagande à ses camarades.
\item \emph{Fait E} : Lors d'une nuit de troubles, des inconnus incendient le bloc administratif et la salle informatique d'une école publique, détruisant 40 ordinateurs et les archives des élèves.
\end{enumerate}

\textbf{Document 2 — Quelques données chiffrées (sources : MINPROFF, OCHA, rapports nationaux).}\par

\begin{center}
\begin{tabularx}{\linewidth}{|l|X|}
\hline
\rowcolor{popblueL} \textbf{Indicateur} & \textbf{Donnée} \\
\hline
Femmes ayant subi des violences physiques au cours de leur vie (Cameroun) & environ 4 sur 10 (près de 40 \%) \\
\hline
Personnes déplacées internes (régions de l'Extrême-Nord, du Nord-Ouest et du Sud-Ouest) & plus d'un million de personnes, dont une majorité de femmes et d'enfants \\
\hline
Écoles non fonctionnelles ou attaquées dans les zones en crise & plusieurs centaines, privant des milliers d'enfants d'éducation \\
\hline
Plaintes liées aux conflits fonciers devant les juridictions & les affaires foncières représentent une part très importante (souvent majoritaire) du contentieux civil \\
\hline
\end{tabularx}
\end{center}

\textbf{Document 3 — Extraits simplifiés du Code pénal camerounais.}\par

\begin{itemize}
\item Les coups et blessures volontaires sont punis d'emprisonnement et d'amende, les peines étant aggravées lorsque la victime est un mineur ou lorsque l'atteinte entraîne une incapacité de travail.
\item Le harcèlement, y compris par voie électronique, est réprimé ; la loi de 2010 sur la cybersécurité et la cybercriminalité punit les atteintes aux personnes commises au moyen des technologies de l'information et de la communication.
\item L'escroquerie et l'abus de confiance (dont la double vente d'un terrain) sont des infractions punies par le Code pénal.
\item Le terrorisme, l'apologie du terrorisme, le financement et le recrutement au profit de groupes armés sont sévèrement réprimés par la loi de 2014 sur la répression des actes de terrorisme.
\item La destruction volontaire des biens d'autrui, notamment des biens publics (écoles, hôpitaux), constitue un délit.
\end{itemize}

\courssub{Consignes}

Travail en groupes de quatre à cinq élèves. Chaque groupe dispose de 55 minutes de préparation, suivies d'une restitution orale de 5 minutes par groupe et d'un débat collectif.

\begin{enumerate}
\item \textbf{Qualifier.} Pour chacun des cinq faits divers (A à E), indiquez : (a) la forme de violence commise (atteinte à l'intégrité physique, atteinte à l'intégrité morale, atteinte au patrimoine, extrémisme violent) ; (b) la victime et le préjudice subi ; (c) la qualification juridique possible à l'aide du document 3 ; (d) la réaction citoyenne appropriée (plainte, signalement, médiation, alerte).
\item \textbf{Cartographier.} Sur une carte schématique du Cameroun, localisez à l'aide de symboles et d'une légende les principales formes de violences étudiées : violences basées sur le genre (tout le territoire), insécurité liée aux groupes armés (Extrême-Nord), tensions et destructions d'écoles (Nord-Ouest et Sud-Ouest), conflits fonciers (grandes villes et périphéries), cyberharcèlement (zones urbaines connectées).
\item \textbf{Analyser les données.} À partir du document 2, dégagez deux constats chiffrés sur l'ampleur des violences au Cameroun et expliquez pourquoi les femmes, les enfants et les déplacés internes sont particulièrement vulnérables.
\item \textbf{Proposer.} Rédigez un \cle{plan d'action citoyen} en cinq mesures concrètes pour prévenir et combattre les violences dans votre lycée ou votre quartier. Chaque mesure doit être justifiée par une notion précise du cours (intégrité physique, intégrité morale, patrimoine, refus de l'extrémisme, rôle du citoyen).
\item \textbf{Restituer et débattre.} Présentez votre qualification, votre carte et votre plan d'action à la classe. Répondez aux questions des autres groupes et du professeur, puis participez au débat : « La violence est-elle une fatalité dans notre environnement ? »
\end{enumerate}

\courssub{Ressources à mobiliser}

\begin{aretenir}[Notions utiles de la leçon]
\begin{itemize}
\item \textbf{Atteinte à l'intégrité physique} : toute violence qui porte atteinte au corps d'une personne (coups, blessures, mutilations, violences sexuelles, meurtre). Elle est punie par le Code pénal.
\item \textbf{Atteinte à l'intégrité morale} : toute violence qui blesse la dignité, l'honneur ou l'équilibre psychologique (insultes, diffamation, harcèlement, cyberharcèlement, discrimination, humiliation).
\item \textbf{Atteinte au patrimoine} : toute violence qui porte atteinte aux biens d'une personne ou de la collectivité (vol, escroquerie, double vente de terrain, destruction, vandalisme).
\item \textbf{Extrémisme violent} : adhésion à une idéologie radicale qui justifie l'usage de la violence (terrorisme, recrutement, propagande) ; il se combat par l'éducation, la vigilance citoyenne et la répression légale.
\item \textbf{Rôle du citoyen} : dénoncer, signaler, protéger les victimes, privilégier le dialogue et la médiation, respecter la loi, éduquer à la paix.
\end{itemize}
\end{aretenir}

\begin{methode}[Méthode de qualification d'un fait de violence]
\begin{enumerate}
\item Identifier les faits précis : qui fait quoi, à qui, avec quelles conséquences ?
\item Déterminer ce qui est atteint : le corps (intégrité physique), la dignité (intégrité morale), les biens (patrimoine), ou la sécurité collective par une idéologie radicale (extrémisme violent).
\item Relier les faits à un texte de loi (Code pénal, loi sur la cybercriminalité, loi antiterroriste).
\item Conclure par la conduite citoyenne à tenir : porter plainte, signaler, alerter une autorité, protéger la victime.
\end{enumerate}
\end{methode}

\courssub{Démarche et corrigé commenté}

\begin{exoresolu}[Qualification des faits, cartographie et plan d'action]
Énoncé : qualifier juridiquement les cinq faits divers, construire la carte schématique des violences au Cameroun et rédiger un plan d'action citoyen en cinq mesures justifiées.
\tcblower
\textbf{Étape 1 — Qualification des cinq faits.}\par
\emph{Fait A (agression à la sortie de l'atelier)} : l'atteinte porte sur le corps de la victime (bras cassé, huit jours d'hospitalisation). C'est une \cle{atteinte à l'intégrité physique}. Juridiquement, il s'agit de coups et blessures volontaires, punis par le Code pénal, avec circonstance aggravante (victime élève, incapacité temporaire de travail). Réaction citoyenne : secourir la victime, alerter l'administration du lycée, porter plainte au commissariat ou à la gendarmerie.\par
\emph{Fait B (groupe WhatsApp)} : l'atteinte porte sur la dignité et l'équilibre psychologique de l'élève (humiliation, dépression, déscolarisation). C'est une \cle{atteinte à l'intégrité morale}, sous la forme précise du cyberharcèlement. Juridiquement, la loi de 2010 sur la cybersécurité et la cybercriminalité et le Code pénal répriment ces agissements. Réaction citoyenne : ne pas relayer, conserver les preuves (captures d'écran), signaler à l'administration et à la plateforme, soutenir la victime.\par
\emph{Fait C (double vente du terrain)} : l'atteinte porte sur les biens des deux familles (perte de \num{1500000} FCFA chacune). C'est une \cle{atteinte au patrimoine}, qualifiée juridiquement d'escroquerie et d'abus de confiance. Réaction citoyenne : saisir le chef de quartier pour la conciliation, puis porter plainte ; vérifier toujours les titres fonciers avant tout achat.\par
\emph{Fait D (recrutement radical)} : le jeune adhère à une idéologie qui justifie la violence et diffuse de la propagande. C'est de l'\cle{extrémisme violent}. Juridiquement, le recrutement, l'apologie et le financement du terrorisme sont réprimés par la loi de 2014. Réaction citoyenne : ne jamais partager les contenus, alerter discrètement les autorités (forces de sécurité, chef d'établissement, famille), dialoguer avec le jeune pour le détourner du groupe.\par
\emph{Fait E (incendie de l'école)} : l'atteinte porte sur les biens de la collectivité (40 ordinateurs, archives). C'est une \cle{atteinte au patrimoine public}, qualifiée de destruction volontaire de biens, délit aggravé car il s'agit d'un bien public. Réaction citoyenne : dénoncer les auteurs, participer à la protection des infrastructures scolaires, sensibiliser au respect du bien commun.\par
\textbf{Étape 2 — Carte schématique.} Le groupe dessine un contour simplifié du Cameroun, place les villes repères (Yaoundé, Douala, Maroua, Bamenda, Buea) et associe à chaque zone un symbole légendé : poing barré pour les violences physiques et basées sur le genre (tout le territoire), étoile pour l'extrémisme violent (Extrême-Nord), flamme pour les écoles attaquées (Nord-Ouest et Sud-Ouest), parcelle pour les conflits fonciers (périphéries de Yaoundé et Douala), écran pour le cyberharcèlement (villes connectées).\par
\textbf{Étape 3 — Analyse des données.} Constat 1 : près de 40 \% des femmes ont subi des violences physiques, ce qui montre que les violences basées sur le genre sont un problème massif et national. Constat 2 : plus d'un million de déplacés internes, majoritairement des femmes et des enfants, vivent dans la précarité ; coupés de leur protection habituelle (famille, école, revenus), ils sont exposés à toutes les formes de violences.\par
\textbf{Étape 4 — Plan d'action citoyen en cinq mesures.}\par
\begin{enumerate}
\item Créer un comité d'écoute et d'alerte contre le harcèlement et le cyberharcèlement au lycée (justification : protection de l'intégrité morale).
\item Organiser une journée de sensibilisation aux violences physiques avec un officier de police judiciaire (justification : protection de l'intégrité physique, prévention par l'information sur les peines encourues).
\item Afficher dans le quartier les démarches légales de vérification des titres fonciers avant tout achat de terrain (justification : protection du patrimoine, lutte contre l'escroquerie).
\item Mettre en place des ateliers d'éducation aux médias pour repérer la propagande et les discours de haine en ligne (justification : prévention de l'extrémisme violent par l'esprit critique).
\item Lancer une brigade citoyenne de protection des infrastructures scolaires et un club de médiation des conflits (justification : protection du patrimoine public, rôle actif du citoyen dans la culture de la paix).
\end{enumerate}
\textbf{Étape 5 — Restitution.} Chaque groupe présente sa grille de qualification, sa carte et son plan ; le débat collectif fait ressortir que la violence n'est pas une fatalité : la loi la punit et l'action citoyenne la prévient.
\end{exoresolu}

\begin{popfigure}
\begin{center}
\begin{tikzpicture}[scale=1.1]
% Contour très simplifié du Cameroun
\draw[thick,popdark,fill=popcream] (0,0) -- (1.6,-0.4) -- (3.2,-0.2) -- (4.2,0.6) -- (4.0,1.8) -- (3.4,2.6) -- (3.6,3.8) -- (3.0,5.0) -- (2.2,6.2) -- (1.4,7.0) -- (0.8,6.0) -- (0.2,4.6) -- (-0.6,3.6) -- (-0.8,2.2) -- (-0.4,1.0) -- cycle;
% Villes repères
\fill[popblue] (1.2,1.6) circle (0.07); \node[right,font=\scriptsize] at (1.3,1.6) {Yaoundé};
\fill[popblue] (0.0,1.2) circle (0.07); \node[left,font=\scriptsize] at (-0.1,1.2) {Douala};
\fill[popblue] (1.5,6.3) circle (0.07); \node[right,font=\scriptsize] at (1.6,6.3) {Maroua};
\fill[popblue] (-0.2,2.6) circle (0.07); \node[left,font=\scriptsize] at (-0.3,2.6) {Bamenda};
\fill[popblue] (-0.4,1.8) circle (0.07); \node[left,font=\scriptsize] at (-0.5,1.8) {Buea};
% Symboles légendés directement sur la carte
% Extrémisme violent (Extrême-Nord) : étoile rouge
\fill[poppurple] (1.5,5.6) circle (0.1);
\node[font=\scriptsize\bfseries,text=poppurple] at (1.5,5.2) {Extrémisme violent};
% Écoles attaquées (Nord-Ouest / Sud-Ouest) : flamme orange
\draw[poporange,thick] (-0.2,3.0) -- (-0.1,2.6) -- (0.0,3.0) -- (-0.05,2.8) -- cycle;
\node[font=\scriptsize\bfseries,text=poporange] at (-0.2,3.4) {Écoles attaquées};
% Conflits fonciers (périphéries Yaoundé/Douala) : parcelle verte
\draw[popgreen,thick] (2.6,1.0) rectangle (2.9,1.3);
\draw[popgreen,thick] (2.6,1.0) -- (2.75,1.15) -- (2.9,1.0);
\node[font=\scriptsize\bfseries,text=popgreen] at (2.75,0.7) {Conflits fonciers};
% Cyberharcèlement (villes connectées) : écran rose
\draw[popink,thick,rounded corners] (2.8,2.0) rectangle (3.2,2.4);
\draw[popink,thick] (3.0,2.4) -- (3.0,2.6);
\node[font=\scriptsize\bfseries,text=popink] at (3.0,2.7) {Cyberharcèlement};
% VBG : tout le territoire (poing barré) - symbole au sud
\draw[popdark,thick] (2.2,0.2) -- (2.4,0.0) -- (2.6,0.2) -- (2.4,0.1) -- cycle;
\draw[popdark,thick] (2.3,0.1) -- (2.5,0.1);
\node[font=\scriptsize\bfseries,text=popdark] at (2.2,-0.3) {VBG : tout le territoire};
% Flèche Nord
\draw[->,thick,popdark] (4.6,6.0) -- (4.6,6.8); \node[font=\scriptsize] at (4.6,7.0) {N};
% Légende intégrée en bas à gauche
\begin{scope}[shift={(-1.2,-0.8)}]
\draw[popdark,fill=popcream] (-0.1,-0.1) rectangle (3.5,1.8);
\node[font=\scriptsize\bfseries] at (1.7,1.5) {Légende};
\fill[poppurple] (0.2,1.1) circle (0.07); \node[right,font=\scriptsize] at (0.35,1.1) {Extrémisme violent (Extrême-Nord)};
\draw[poporange,thick] (0.2,0.8) -- (0.1,0.6) -- (0.3,0.8) -- (0.15,0.7) -- cycle; \node[right,font=\scriptsize] at (0.35,0.8) {Écoles attaquées (NO/SO)};
\draw[popgreen,thick] (0.2,0.5) rectangle (0.35,0.65); \draw[popgreen,thick] (0.2,0.5) -- (0.275,0.575) -- (0.35,0.5); \node[right,font=\scriptsize] at (0.35,0.55) {Conflits fonciers (villes)};
\draw[popink,thick,rounded corners] (0.2,0.2) rectangle (0.4,0.4); \draw[popink,thick] (0.3,0.4) -- (0.3,0.5); \node[right,font=\scriptsize] at (0.35,0.3) {Cyberharcèlement (urbain)};
\draw[popdark,thick] (0.2,-0.1) -- (0.4,-0.3) -- (0.6,-0.1) -- (0.4,-0.2) -- cycle; \draw[popdark,thick] (0.3,-0.2) -- (0.5,-0.2); \node[right,font=\scriptsize] at (0.35,-0.15) {VBG (tout le territoire)};
\end{scope}
\end{tikzpicture}
\end{center}
\legende{Carte schématique (croquis simplifié) de localisation des principales formes de violences au Cameroun : extrémisme violent (Extrême-Nord, région de Maroua), écoles attaquées (Nord-Ouest et Sud-Ouest, régions de Bamenda et Buea), conflits fonciers (périphéries de Yaoundé et Douala), cyberharcèlement (grandes villes connectées), violences basées sur le genre — VBG (ensemble du territoire national). Croquis indicatif, sans précision cartographique.}
\end{popfigure}

\courssub{Bilan de l'activité}

Cette enquête citoyenne a permis de mettre en évidence plusieurs acquis essentiels. D'abord, les élèves ont appris à \cle{qualifier} précisément une violence : un même mot, « violence », recouvre des réalités différentes (atteinte au corps, à la dignité, aux biens, ou adhésion à une idéologie radicale), qui appellent des réponses juridiques et citoyennes différentes. Ensuite, la cartographie a montré que les violences ne sont pas abstraites : elles ont des lieux, des victimes et des chiffres, au Cameroun comme dans l'environnement proche du lycée. L'analyse des données du MINPROFF et d'OCHA a sensibilisé à la vulnérabilité particulière des femmes, des enfants et des déplacés internes. Enfin, la rédaction du plan d'action a transformé les notions du cours en \cle{compétences d'action} : prévenir, signaler, protéger, médiatiser. Le débat final a confirmé la conclusion centrale de la leçon : la violence n'est pas une fatalité ; entre la loi qui la punit et le citoyen qui la prévient, chaque jeune peut devenir un acteur de la paix dans son lycée, son quartier et son pays.

\newpage

\coursec{Méthodes et exercices résolus}

\courssub{Fiches méthodes et exercices résolus}

\begin{aretenir}[Synthèse des formes de violences au programme]
Pour le Probatoire / Baccalauréat technique, il faut maîtriser les quatre catégories d'atteintes et leur vocabulaire précis :
\begin{itemize}
\item \textbf{Atteinte à l'intégrité physique} : coups et blessures, séquestration, mutilation, homicide. Le corps est visé.
\item \textbf{Atteinte à l'intégrité morale} : injures, diffamation, calomnie, harcèlement (moral, scolaire, cyberharcèlement), menaces, discrimination, humiliations. La dignité, l'honneur, la réputation ou la tranquillité psychologique sont visés.
\item \textbf{Atteinte au patrimoine} : vol, escroquerie, abus de confiance, dégradation volontaire (vandalisme), incendie volontaire, racket. Les biens matériels ou l'argent sont visés.
\item \textbf{Extrémisme violent} : recours à la terreur (attentats, enlèvements, destructions massives) au nom d'une idéologie (religieuse, politique, identitaire) refusant le dialogue. Processus : radicalisation $\rightarrow$ recrutement $\rightarrow$ passage à l'acte. Exemple camerounais : Boko Haram (Extrême-Nord).
\end{itemize}
\textbf{Règle d'or :} une situation concrète combine souvent plusieurs formes (ex. racket = vol + menaces = atteinte au patrimoine \emph{et} morale). Il faut toutes les identifier et les nommer avec le vocabulaire du cours.
\end{aretenir}

\begin{methode}[Appliquer les notions de l'unité à une situation concrète de la vie courante]
Face à une situation (atelier, entreprise, quartier, réseaux sociaux, établissement scolaire), suivez cette démarche rigoureuse :
\begin{enumerate}
\item \textbf{Lire attentivement} et souligner les faits bruts : qui agit ? sur qui ? quelle action ? quelles conséquences immédiates ?
\item \textbf{Identifier la ou les formes de violence} en comparant chaque fait aux définitions du cours : atteinte physique (corps/santé), atteinte morale (dignité/honneur/psychisme), atteinte au patrimoine (biens/argent), extrémisme violent (idéologie + terreur).
\item \textbf{Nommer précisément la notion} avec le vocabulaire technique du cours (\cle{agression}, \cle{harcèlement}, \cle{diffamation}, \cle{vandalisme}, \cle{radicalisation}, \cle{racket}...) et rejeter les termes vagues (« méchanceté », « problème », « bagarre »).
\item \textbf{Justifier l'identification} : pour chaque fait, citer l'élément du texte qui correspond à la définition (ex. « il l'a frappé » $\rightarrow$ atteinte physique car atteinte au corps).
\item \textbf{Conclure} en synthétisant les violences subies et, si la question le demande, en indiquant la réaction citoyenne appropriée (conserver les preuves, alerter un adulte de confiance, saisir l'institution compétente : inspection du travail, chef d'établissement, police, gendarmerie, tribunal).
\end{enumerate}
\textbf{Réflexe Bac :} une même situation contient \textbf{souvent plusieurs formes de violences} ; il faut toutes les repérer, les distinguer et les classer sans en oublier.
\end{methode}

\begin{exoresolu}[Identifier et classer les formes de violences dans une situation professionnelle]
Énoncé. Dans un atelier de mécanique de Douala, l'apprenti K. est régulièrement insulté et ridiculisé devant les clients par son maître d'apprentissage, qui l'a également frappé à deux reprises et retient depuis trois mois le petit pécule que l'entreprise verse aux apprentis. Un jour, furieux, le maître détruit la boîte à outils que K. avait achetée avec ses économies personnelles.
\begin{enumerate}
\item Relevez dans cette situation les différentes formes de violences subies par K.
\item Classez chaque fait relevé dans la catégorie correspondante : atteinte à l'intégrité physique, atteinte à l'intégrité morale, atteinte au patrimoine.
\end{enumerate}
\tcblower
Solution détaillée.
\begin{enumerate}
\item \textbf{Relevé des faits (lecture active).} On isole quatre faits distincts dans le récit :
\begin{itemize}
\item (a) Insultes et ridiculisation publique devant la clientèle (répétition : « régulièrement »).
\item (b) Coups portés à deux reprises (agression physique directe).
\item (c) Rétention du pécule versé par l'entreprise pendant trois mois (privation de revenu).
\item (d) Destruction volontaire de la boîte à outils personnelle de K. (dégradation de bien propre).
\end{itemize}
\item \textbf{Classement argumenté à l'aide des définitions du cours.}
\begin{itemize}
\item \textbf{Fait (b) : Atteinte à l'intégrité physique.} Les coups portent une atteinte directe au corps et à la santé de l'apprenti. Juridiquement, il s'agit de \cle{coups et blessures volontaires} (Code pénal).
\item \textbf{Fait (a) : Atteinte à l'intégrité morale.} Les insultes et l'humiliation publique visent la dignité, l'honneur et l'estime de soi. La répétition caractérise un \cle{harcèlement moral} en milieu professionnel.
\item \textbf{Fait (c) : Atteinte au patrimoine.} Le maître prive K. d'une somme d'argent (le pécule) qui lui revient de droit. C'est un \cle{vol} ou un \cle{abus de confiance} (détournement de fonds confiés).
\item \textbf{Fait (d) : Atteinte au patrimoine.} La destruction de la boîte à outils (bien meuble appartenant à la victime) est une \cle{dégradation volontaire} (vandalisme), évaluable en argent (coût de remplacement).
\end{itemize}
\end{enumerate}
\textbf{Conclusion synthétique.} K. est victime de \textbf{trois formes de violences distinctes} : une atteinte à l'intégrité physique (coups), une atteinte à l'intégrité morale (injures, humiliations répétées = harcèlement) et une double atteinte au patrimoine (rétention d'argent + destruction d'outillage).
\textbf{Démarche citoyenne :} K. doit conserver des preuves (certificat médical pour les coups, témoins pour les insultes, attestation de versement du pécule), alerter l'inspection du travail ou son syndicat, et porter plainte auprès du procureur de la République ou de la police judiciaire.
\end{exoresolu}

\begin{methode}[Rédiger et justifier une réponse argumentée en ECM (épreuve écrite)]
Pour traiter une question de synthèse, un conseil à une victime ou une prise de position (10 à 15 lignes attendues au Bac) :
\begin{enumerate}
\item \textbf{Introduction :} Reformuler le sujet en une phrase d'accroche qui annonce votre thèse ou la réponse attendue.
\item \textbf{Définition :} Définir en une phrase la notion clé mobilisée (ex. : « Le cyberharcèlement est une atteinte à l'intégrité morale commise via des supports numériques... »).
\item \textbf{Développement (2 arguments minimum) :} Chaque argument = \textbf{Idée force} + \textbf{Justification} (connecteur logique : \emph{car, en effet, puisque}) + \textbf{Exemple concret} (tiré de l'énoncé ou de la vie courante camerounaise).
\item \textbf{Appui juridique / institutionnel :} Citer un texte de référence : Constitution (Préambule : dignité humaine), Code pénal camerounais (art. 277 s. coups, 318 s. vol, 305 s. diffamation), Loi n°2014/028 du 23 déc. 2014 (répression du terrorisme), Convention relative aux droits de l'enfant.
\item \textbf{Conclusion :} Répondre explicitement à la question posée et ouvrir sur l'attitude citoyenne (refus de la vengeance, recours à la loi, solidarité).
\end{enumerate}
\textbf{Exigence de rédaction :} Toute affirmation doit être suivie de sa justification. Pas d'affirmation « en l'air ».
\end{methode}

\begin{exoresolu}[Rédiger une réponse justifiée : conseiller une victime de cyberharcèlement]
Énoncé. Votre camarade N., élève en Première technique, est victime depuis plusieurs semaines de moqueries, d'insultes et de menaces sur un groupe de messagerie de la classe ; ses photos ont été détournées et diffusées sans son accord. Il envisage de répondre par la violence physique contre ses harceleurs. Rédigez un texte de dix à quinze lignes dans lequel vous identifiez la violence subie et vous le convainquez d'adopter une démarche citoyenne plutôt que la vengeance.
\tcblower
Solution détaillée (texte modèle attendu au Bac).
« N. est victime d'une \cle{atteinte à l'intégrité morale} sous forme de \cle{cyberharcèlement} : les moqueries répétées, les insultes, les menaces et le détournement de ses photos blessent gravement sa dignité, son honneur et sa tranquillité psychologique, sans porter atteinte à son corps. Cette violence est pleinement réelle et punissable, car la répétition des actes et l'intention de nuire la caractérisent, quel que soit le support utilisé.

Répondre par la violence physique serait une erreur lourde de conséquences : N. deviendrait à son tour auteur d'une \cle{atteinte à l'intégrité physique} (coups et blessures), passible de poursuites pénales (Code pénal), et il aggraverait le conflit au lieu de le résoudre. La vengeance transforme la victime en coupable et ne fait pas cesser le harcèlement.

La démarche citoyenne s'articule en trois temps : (1) \textbf{Conserver les preuves} (captures d'écran, copies des messages, signalement à la plateforme) pour étayer toute procédure ; (2) \textbf{Alerter immédiatement un adulte de confiance} (parents, censeur, conseiller d'orientation, professeur principal, infirmier scolaire) : l'établissement a l'obligation de protéger ses élèves et peut sanctionner les auteurs (règlement intérieur, conseil de discipline) ; (3) \textbf{Saisir la justice} si les menaces persistent : la diffamation, l'outrage et les menaces de mort sont des délits (Code pénal), la plainte auprès du procureur ou de la police judiciaire fait cesser l'impunité.

En refusant la spirale de la violence et en choisissant l'alerte, la preuve et le droit, N. protège sa dignité, fait respecter ses droits fondamentaux et contribue à construire un environnement scolaire plus sûr pour tous. »
\textbf{Analyse de la copie :} Structure respectée (Intro/Définition $\rightarrow$ 2 arguments développés/justifiés $\rightarrow$ Appui Code pénal $\rightarrow$ Conclusion citoyenne). Vocabulaire précis (cyberharcèlement, atteinte à l'intégrité morale, coups et blessures, diffamation, procureur). Ton adapté (conseil fraternel mais ferme).
\end{exoresolu}

\begin{methode}[Distinguer sans erreur les trois atteintes : physique, morale, patrimoine]
Pour classer un fait de violence dans le bon tableau ou la bonne colonne :
\begin{enumerate}
\item \textbf{Poser la question décisive :} \emph{« Qu'est-ce qui est atteint, visé, endommagé ? »}
\begin{itemize}
\item Le \textbf{corps, la santé, l'intégrité corporelle} $\rightarrow$ \cle{Atteinte à l'intégrité physique}.
\item La \textbf{dignité, l'honneur, la réputation, la tranquillité psychologique, l'image} $\rightarrow$ \cle{Atteinte à l'intégrité morale}.
\item Les \textbf{biens matériels, l'argent, les objets, les récoltes, le logement, les outils de travail} $\rightarrow$ \cle{Atteinte au patrimoine}.
\end{itemize}
\item \textbf{Vérifier la nature du dommage :} Blessure/douleur (physique) ; Souffrance psychologique/humiliation/peur/atteinte à la réputation (morale) ; Perte matérielle évaluable en argent/destruction (patrimoine).
\item \textbf{Attribuer le vocabulaire précis (lexique du Bac) :}
\begin{center}
\begin{tabular}{|l|l|}\hline
\textbf{Physique} & Coups et blessures, séquestration, mutilation, homicide, torture \\ \hline
\textbf{Morale} & Injures, diffamation, calomnie, menaces, harcèlement (moral, scolaire, cyber), discrimination, outrages \\ \hline
\textbf{Patrimoine} & Vol, escroquerie, abus de confiance, vandalisme (dégradation), incendie volontaire, racket, recel \\ \hline
\end{tabular}
\end{center}
\item \textbf{Gérer les cas mixtes (fréquents au Bac) :} Un même acte ou une même situation peut produire plusieurs atteintes. Ex. : Racket = Vol (patrimoine) + Menaces (morale) + éventuellement Coups (physique). \textbf{Obligation :} toutes les identifier et les lier aux faits précis de l'énoncé.
\end{enumerate}
\end{methode}

\begin{exoresolu}[Classer des faits de violence dans un tableau à double entrée]
Énoncé. Recopiez et complétez le tableau suivant en plaçant chaque fait numéroté dans la ou les colonnes qui conviennent :
\begin{enumerate}
\item Des inconnus incendient le hangar de stockage d'un commerçant au marché central.
\item Une rumeur mensongère accuse faussement une élève de vol dans tout le quartier.
\item Un agent de sécurité (vigile) gifle un client qui contestait un contrôle.
\item Un groupe de « grands » bloque un élève chaque matin à l'entrée du lycée et lui prend son argent de transport en le menaçant de représailles.
\item Un employeur humilie quotidiennement son personnel en le traitant publiquement d'« incapables » et de « bons à rien ».
\end{enumerate}
\begin{center}
\begin{tabularx}{\linewidth}{|l|X|}\hline
\rowcolor{popblueL} \textbf{Forme d'atteinte} & \textbf{Fait(s) correspondant(s) (numéro + qualification précise)} \\ \hline
Atteinte à l'intégrité physique & \\ \hline
Atteinte à l'intégrité morale & \\ \hline
Atteinte au patrimoine & \\ \hline
\end{tabularx}
\end{center}
\tcblower
Solution détaillée et commentée.
\begin{itemize}
\item \textbf{Fait (1) :} Incendie du hangar $\rightarrow$ destruction d'un bien immobilier professionnel. \textbf{Atteinte au patrimoine} (incendie volontaire, art. 238 Code pénal).
\item \textbf{Fait (2) :} Rumeur fausse imputant un vol $\rightarrow$ atteinte à l'honneur et à la réputation. \textbf{Atteinte à l'intégrité morale} (\cle{diffamation} si public, \cle{calomnie} si dénonciation à l'autorité).
\item \textbf{Fait (3) :} Gifle $\rightarrow$ contact physique non consenti, atteinte au corps. \textbf{Atteinte à l'intégrité physique} (\cle{coups et blessures volontaires}, même sans ITT).
\item \textbf{Fait (4) :} \textbf{Cas mixte type Bac.} (a) Prise de l'argent de transport $\rightarrow$ \textbf{Atteinte au patrimoine} (\cle{vol} avec violence ou \cle{racket}). (b) Menaces répétées quotidiennes créant la peur $\rightarrow$ \textbf{Atteinte à l'intégrité morale} (\cle{harcèlement}, \cle{menaces} réitérées). (c) Blocage physique $\rightarrow$ possible \textbf{Atteinte à l'intégrité physique} (\cle{séquestration} brève ou violences).
\item \textbf{Fait (5) :} Humiliations publiques quotidiennes $\rightarrow$ atteinte à la dignité des travailleurs. \textbf{Atteinte à l'intégrité morale} (\cle{harcèlement moral} au travail).
\end{itemize}
\textbf{Tableau complété (modèle de réponse attendue) :}
\begin{center}
\begin{tabularx}{\linewidth}{|l|X|}\hline
\rowcolor{popblueL} \textbf{Forme d'atteinte} & \textbf{Fait(s) correspondant(s)} \\ \hline
Atteinte à l'intégrité physique & (3) Gifle du vigile sur le client ; (4) Blocage physique de l'élève (séquestration/violences) \\ \hline
Atteinte à l'intégrité morale & (2) Rumeur diffamatoire sur l'élève ; (5) Humiliations quotidiennes du personnel (harcèlement moral) ; (4) Menaces de représailles quotidiennes (harcèlement, menaces) \\ \hline
Atteinte au patrimoine & (1) Incendie du hangar (destruction) ; (4) Vol de l'argent de transport (racket) \\ \hline
\end{tabularx}
\end{center}
\textbf{Point de vigilance :} Le fait (4) est le « piège » classique : il rapporte des points dans \textbf{trois colonnes}. Ne pas le signaler dans toutes ses dimensions fait perdre la moitié des points de l'exercice.
\end{exoresolu}

\begin{methode}[Comprendre, expliquer et situer l'extrémisme violent (contexte camerounais)]
Pour traiter toute question sur l'extrémisme violent (définition, mécanisme, prévention, exemple) :
\begin{enumerate}
\item \textbf{Définir rigoureusement :} L'\cle{extrémisme violent} est l'adoption et la mise en œuvre d'une idéologie (religieuse, politique, identitaire) qui \textbf{rejette le dialogue, la différence, l'État de droit et la loi}, et qui \textbf{légitime l'usage de la terreur} (attentats, enlèvements, assassinats ciblés, destructions massives) pour imposer sa vision.
\item \textbf{Décrire le processus d'embrigadement (radicalisation) :} Ce n'est pas un acte unique mais un parcours :
\begin{enumerate}
\item \textbf{Fragilités préexistantes :} pauvreté, chômage, déscolarisation, sentiment d'injustice, d'exclusion ou d'humiliation, quête de sens/reconnaissance.
\item \textbf{Rencontre avec un recruteur / exposition à la propagande :} souvent via les réseaux sociaux (vidéos, sermons déviants, promesses de paradis/argent/pouvoir) ou par endoctrinement direct.
\item \textbf{Rupture progressive :} isolement familial/amical, adoption d'un discours binaire (nous/eux, croyants/mécréants), légitimation de la violence.
\item \textbf{Passage à l'acte :} adhésion au groupe armé, entraînement, participation aux exactions.
\end{enumerate}
\item \textbf{Citer l'exemple camerounais de référence :} Les exactions du groupe \textbf{Boko Haram} (et de l'État islamique en Afrique de l'Ouest - ISWAP) dans la région de l'\textbf{Extrême-Nord} : attaques de villages, marchés, écoles, postes de sécurité ; enlèvements massifs (ex. lycéens de Dabanga, 2014) ; utilisation d'enfants-soldats et de kamikazes (souvent des filles) ; destruction de l'économie locale (agriculture, commerce transfrontalier).
\item \textbf{Identifier les réponses de l'État et de la société :}
\begin{itemize}
\item \textbf{Judiciaire :} Loi n°2014/028 du 23 décembre 2014 portant répression des actes de terrorisme (peines lourdes, peine de mort pour les auteurs d'attentats mortels).
\item \textbf{Militaire :} Opérations de l'armée (BIR, secteur opérationnel) et Force multinationale mixte (FMM) avec Nigeria, Tchad, Niger.
\item \textbf{Prévention / Déradicalisation :} Centres de désengagement (ex. centre de Mérina), programmes de réinsertion socio-économique (formation professionnelle, agriculture), sensibilisation par les leaders religieux (imams, pasteurs) et traditionnels, éducation à la citoyenneté (ECM) à l'école.
\item \textbf{Vigilance citoyenne :} Signalement des mouvements suspects, refus de la propagande haineuse, solidarité communautaire.
\end{itemize}
\end{enumerate}
\end{methode}

\begin{exoresolu}[Analyser un cas d'extrémisme violent et proposer des réponses préventives]
Énoncé. Dans un village de l'Extrême-Nord, plusieurs jeunes déscolarisés, sans emploi et sans perspective, sont approchés par un individu charismatique. Il leur promet de l'argent, une arme, un « grand destin » et la « défense de la vraie religion » s'ils rejoignent un groupe armé qui « combat les infidèles ». Deux d'entre eux acceptent et participent quelques semaines plus tard à une attaque meurtrière contre un marché hebdomadaire.
\begin{enumerate}
\item Définissez l'extrémisme violent et montrez, en vous appuyant sur le texte, que cette situation en relève.
\item Expliquez par quels mécanismes précis les deux jeunes ont été entraînés vers la violence.
\item Proposez deux mesures de prévention concrètes et adaptées à cette communauté villageoise.
\end{enumerate}
\tcblower
Solution détaillée.
\begin{enumerate}
\item \textbf{Définition et qualification.} L'extrémisme violent est le recours à la violence terroriste (attentats, attaques de civils, enlèvements) pour imposer une idéologie totalitaire (religieuse, politique, identitaire) qui refuse le pluralisme, la démocratie et l'État de droit. La situation en relève pleinement : le groupe justifie ses actes par une idéologie religieuse déviée (« défense de la vraie religion », combat contre « les infidèles ») et commet une \textbf{attaque directe contre des civils} sur un marché (cible molle, terreur), ce qui est la signature tactique de Boko Haram dans l'Extrême-Nord.
\item \textbf{Mécanismes d'embrigadement (radicalisation) des deux jeunes :}
\begin{enumerate}
\item \textbf{Terrain de vulnérabilité :} Déscolarisation + chômage + pauvreté $\rightarrow$ sentiment d'inutilité sociale, d'injustice, absence d'horizon, fragilité psychologique.
\item \textbf{Manipulation / Endoctrinement :} Le recruteur exploite ces failles : promesse d'argent (appât matériel), promesse de statut social (« grand destin », héros), réinterprétation religieuse (violence = devoir sacré, martyr = paradis). Rupture avec la famille/la communauté traditionnelle.
\item \textbf{Engagement progressif :} Adhésion au groupe $\rightarrow$ endoctrinement intensif $\rightarrow$ entraînement $\rightarrow$ passage à l'acte (attaque du marché) pour « prouver » son engagement et sceller l'appartenance.
\end{enumerate}
\item \textbf{Mesures de prévention adaptées (niveau local) :}
\begin{enumerate}
\item \textbf{Insertion socio-économique des jeunes (cause racine) :} Création d'un centre de formation professionnelle décentralisé (métiers agricoles, artisanat, élevage, transformation de produits locaux) + accompagnement à l'installation (microcrédit, foncier). \emph{Justification :} Occuper les jeunes, leur donner un revenu et une dignité retire le « vivier » de recrutement.
\item \textbf{Vigilance communautaire et contre-discours religieux (cause idéologique) :} Mise en place d'un comité de veille villageois (chef traditionnel, imam modéré, notables, femmes, jeunes) chargé de : (a) repérer et signaler tout recruteur suspect aux forces de défense ; (b) organiser des causeries éducatives à la mosquée/église/marché sur le « vrai islam » (paix, travail, respect de la vie) pour démonter l'endoctrinement. \emph{Justification :} La prévention passe par la reprise en main du religieux par les leaders légitimes et la solidarité villageoise (loi 2014 : obligation de signalement).
\end{enumerate}
\end{enumerate}
\textbf{Conclusion.} L'extrémisme violent se combat par la \textbf{force de la loi} (répression) et par la \textbf{force du développement} (prévention). Offrir un avenir aux jeunes et protéger les esprits contre la haine sont les deux piliers de la citoyenneté active dans l'Extrême-Nord comme partout au Cameroun.
\end{exoresolu}

\begin{methode}[Construire un plan d'action citoyen de prévention et de lutte contre les violences]
Pour proposer des mesures concrètes, efficaces et juridiquement fondées (ex. : club citoyen, conseil de la vie lycéenne, association de quartier) :
\begin{enumerate}
\item \textbf{Nommer le problème} avec le vocabulaire exact (ex. : « Le racket est une atteinte au patrimoine doublée d'une atteinte morale par les menaces »).
\item \textbf{Structurer l'action sur 4 niveaux complémentaires} (approche systémique) :
\begin{itemize}
\item \textbf{Niveau 1 - Individuel / Victime :} Refuser de se taire, conserver les preuves, fuir le danger, ne pas riposter par la violence.
\item \textbf{Niveau 2 - Scolaire / Familial / Proximité :} Cellule d'écoute/mediation (pairs formés + adulte référent), alerte immédiate au chef d'établissement / parents / inspection du travail, sanction éducative (réparation, excuse publique).
\item \textbf{Niveau 3 - Communautaire / Institutionnel :} Sensibilisation régulière (causeries, affiches, théâtre-forum), formation des acteurs (élèves, enseignants, chefs de quartier), amélioration de l'environnement (éclairage, présence adulte aux abords).
\item \textbf{Niveau 4 - Judiciaire / Répressif :} Plainte systématique (police, gendarmerie, procureur), application du Code pénal et de la loi antiterroriste, protection des témoins/victimes.
\end{itemize}
\item \textbf{Justifier chaque mesure} par son effet attendu : « Cette mesure permet de... car... ».
\item \textbf{Conclure par une phrase forte} liant non-violence, État de droit et responsabilité citoyenne.
\end{enumerate}
\textbf{Interdit absolu :} Proposer la vengeance, la justice privée, le lynchage, le silence (« il ne faut pas créer d'histoires »). La citoyenneté, c'est \textbf{la loi, pas le poing}.
\end{methode}

\begin{exoresolu}[Élaborer un plan d'action citoyen dans un établissement technique]
Énoncé. Suite à une recrudescence de rackets (vols avec menaces aux abords du lycée) et de bagarres (coups et blessures) dans un lycée technique de Yaoundé, le Club des Élèves Citoyens, dont vous êtes le président, est chargé par le proviseur de proposer un plan d'action concret en quatre mesures. Rédigez ce plan en justifiant chaque mesure, puis concluez.
\tcblower
Solution détaillée (texte modèle structuré pour le Bac).
« Le racket constitue une \cle{atteinte au patrimoine} (vol) aggravée par des \cle{menaces} (atteinte à l'intégrité morale), et les bagarres sont des \cle{atteintes à l'intégrité physique} (coups et blessures volontaires). Pour y mettre fin durablement, le Club des Élèves Citoyens propose le plan d'action suivant, articulé sur quatre niveaux :

\begin{enumerate}
\item \textbf{Niveau Prévention / Sensibilisation (Individuel \& Collectif) :} Organiser, chaque premier mardi du mois, une \emph{« Heure Citoyenne »} obligatoire en classe, animée par les membres du club avec l'appui du CPE ou d'un policier invité, sur les thèmes : « Connaître ses droits », « Les peines encourues pour racket et violences (Code pénal) », « Gérer ses conflits sans violence ». \textbf{Justification :} L'ignorance de la loi et la banalisation de la violence favorisent le passage à l'acte ; informer, c'est dissuader.
\item \textbf{Niveau Écoute / Médiation / Alerte (Scolaire / Proximité) :} Mettre en place une \emph{Cellule « Écoute \& Médiation »} permanente (local identifié, permanences quotidiennes 12h-14h), composée de 6 élèves médiateurs formés (par une ONG partenaire ou le CPE) et d'un enseignant référent. Toute victime ou témoin peut s'y rendre anonymement. \textbf{Justification :} La peur de parler (représailles, honte) est le premier allié des racketteurs ; un espace de confiance brise la loi du silence et permet de régler les conflits naissants avant l'irréparable.
\item \textbf{Niveau Signalement / Sanction Éducative (Institutionnel) :} Établir un \emph{circuit d'alerte clair et connu de tous} : Victime/Témoin $\rightarrow$ Cellule / Professeur principal / Censeur $\rightarrow$ Proviseur + Parents $\rightarrow$ Conseil de discipline / Police (si infraction pénale). Afficher ce circuit dans chaque salle de classe et au portail. \textbf{Justification :} La rapidité et la certitude de la réponse institutionnelle (sanction éducative + réparation + éventuelle plainte) restaurent l'autorité de l'école et protègent les victimes.
\item \textbf{Niveau Sécurisation / Partenariat Judiciaire (Communautaire / Répressif) :} Signer une \emph{convention de partenariat} avec le Commissariat de police du secteur et la Brigade de gendarmerie pour : (a) patrouilles mixtes (police + vie scolaire) aux heures d'entrée/sortie aux abords du lycée ; (b) intervention rapide en cas de signalement ; (c) co-animation d'une séance de sensibilisation sur « La justice des mineurs et les peines réelles ». \textbf{Justification :} La présence visible de la force publique dissuade les auteurs externes (souvent des déscolarisés) et rappelle que le racket est un délit (art. 318 Code pénal : 5 à 10 ans de prison), pas une « bêtise de jeunes ».
\end{enumerate}

\textbf{Conclusion.} En choisissant résolument la \textbf{sensibilisation}, la \textbf{médiation}, l'\textbf{alerte institutionnelle} et le \textbf{partenariat judiciaire} plutôt que la loi du talion ou le silence complice, les élèves de ce lycée technique exercent pleinement leur citoyenneté : ils construisent un établissement où la \textbf{dignité} de chacun est protégée par la \textbf{loi}, où le \textbf{dialogue} prime sur la force, et où la \textbf{formation professionnelle} peut se dérouler dans la sérénité. C'est cela, l'École de la République. »
\textbf{Évaluation de la copie :} Problème nommé + 4 mesures structurées/niveau + Justification cause/effet + Vocabulaire précis (racket, atteinte patrimoine/morale/physique, Code pénal, conseil de discipline, convention partenariat) + Conclusion citoyenne. Note maximale.
\end{exoresolu}

\begin{attention}[Les 5 erreurs fatales qui font perdre des points au Probatoire / Bac Technique ECM]
\begin{enumerate}
\item \textbf{Confondre les trois atteintes fondamentales.} Classer une insulte en « atteinte physique » ou un vol en « atteinte morale ». \textbf{Parade :} Avant d'écrire, posez la question : « Qu'est-ce qui est touché ? Le corps $\rightarrow$ Physique. La dignité/honneur $\rightarrow$ Morale. Les biens/argent $\rightarrow$ Patrimoine. »
\item \textbf{Utiliser un langage courant au lieu du vocabulaire scientifique.} Écrire : « Il l'a traité de tous les noms », « Il lui a pris ses affaires », « Il l'a tapé ». \textbf{Exigence Bac :} « Injures publiques / Diffamation », « Vol / Racket / Abus de confiance », « Coups et blessures volontaires / Agression physique ». Le correcteur note le \textbf{lexique de l'unité}.
\item \textbf{Ignorer les situations mixtes (cas combinés).} Ne relever qu'une violence là où il y en a deux ou trois (ex. : racket = vol + menaces ; violences conjugales = coups + humiliations + destruction téléphone). \textbf{Parade :} Relire l'énoncé 3 fois, surligner \textbf{tous} les verbes d'action, classer \textbf{chaque} verbe.
\item \textbf{Préconiser la vengeance, l'auto-défense disproportionnée ou la justice privée.} Réponses types à proscrire : « Il faut qu'il se défende », « Qu'il appelle ses grands frères », « Qu'on le tabasse à son tour ». \textbf{Seule réponse acceptable :} Alerte (adulte, institution) + Preuves + Droit (plainte, police, justice) + Médiation. La citoyenneté = \textbf{monopole de la violence légitime à l'État}.
\item \textbf{Affirmer sans justifier (le « parce que » manquant).} Écrire : « C'est du harcèlement. » Sans ajouter : « \emph{car} les actes sont répétés, intentionnels, visent à dégrader les conditions de vie/travail de la victime. » \textbf{Règle :} Aucune affirmation nue. Structure : \textbf{Notion} + \textbf{Car / En effet / Puisque} + \textbf{Fait précis du texte / Exemple réel / Article de loi}.
\end{enumerate}
\end{attention}

\newpage

\coursec{Exercices}

\courssub{Je vérifie mes connaissances}

\begin{exercice}[QCM : reconnaître les formes de violences]
Pour chaque question, une seule réponse est exacte. Recopie la lettre de la bonne réponse.
\begin{enumerate}
\item L'atteinte à l'intégrité morale comprend :
\begin{enumerate}
\item le vol de bétail ;
\item l'injure et la diffamation ;
\item la destruction d'un champ de maïs.
\end{enumerate}
\item La calomnie se distingue de la diffamation parce qu'elle :
\begin{enumerate}
\item est toujours écrite ;
\item accuse une personne d'un fait en le sachant faux ;
\item ne peut pas être punie par la loi.
\end{enumerate}
\item Le recel consiste à :
\begin{enumerate}
\item détenir ou vendre un bien dont on sait qu'il provient d'un vol ;
\item tromper une personne pour obtenir de l'argent ;
\item dégrader le bien d'autrui.
\end{enumerate}
\item L'extrémisme violent est :
\begin{enumerate}
\item une simple opinion politique radicale ;
\item le passage à la violence au nom d'une idéologie radicale ;
\item une forme de grève ouvrière.
\end{enumerate}
\end{enumerate}
\end{exercice}

\begin{exercice}[Vrai ou faux, avec justification]
Indique si chaque affirmation est vraie ou fausse, puis justifie ta réponse en une ou deux phrases.
\begin{enumerate}
\item Frapper un camarade à l'atelier est une atteinte à l'intégrité morale.
\item La double vente d'un même terrain est une atteinte au patrimoine.
\item La radicalisation désigne le processus par lequel une personne adopte progressivement des idées extrémistes.
\item Seule la police peut prévenir les violences dans un quartier.
\end{enumerate}
\end{exercice}

\begin{exercice}[Définitions courtes]
Définis en deux phrases maximum chacun des termes suivants : \cle{injure}, \cle{escroquerie}, \cle{viol}, \cle{extrémisme violent}.
\end{exercice}

\begin{exercice}[Associer la situation à l'atteinte]
Recopie le tableau et complète la deuxième colonne en indiquant s'il s'agit d'une atteinte à l'intégrité \cle{physique}, \cle{morale} ou au \cle{patrimoine}.
\begin{center}
\begin{tabularx}{\linewidth}{|X|l|}
\hline
\rowcolor{popblueL} \textbf{Situation observée} & \textbf{Forme d'atteinte} \\
\hline
Un élève est blessé à coups de bâton derrière l'atelier & \\
\hline
Une rumeur mensongère circule sur WhatsApp contre une vendeuse & \\
\hline
Un hangar est incendié volontairement pendant la nuit & \\
\hline
Des menaces de mort sont envoyées par SMS à un chef de quartier & \\
\hline
\end{tabularx}
\end{center}
\end{exercice}

\courssub{Je m'entraîne}

\begin{exercice}[Qualifier juridiquement des faits]
Dans chaque cas, nomme précisément l'infraction commise (injure, diffamation, calomnie, vol, escroquerie, recel, coups et blessures) et justifie ton choix.
\begin{enumerate}
\item À Yaoundé, un vendeur de téléphones remet un appareil volé à son ami en lui disant : « Tu sais d'où ça vient, tais-toi ». L'ami le revend.
\item Un client mécontent écrit sur Facebook qu'un mécanicien de Douala « vole les pièces de ses clients », alors qu'il sait que c'est faux.
\item Un inconnu arrache le sac d'une dame au marché et s'enfuit.
\item Lors d'une dispute, un voisin traite un autre de « bon à rien » devant tout le quartier.
\end{enumerate}
\end{exercice}

\begin{exercice}[Distinguer radicalisation et passage à l'acte]
Un jeune de l'Extrême-Nord commence par fréquenter un groupe qui prône la haine des autres communautés, puis refuse tout dialogue, et finit par participer à une attaque contre un village.
\begin{enumerate}
\item Repère dans ce récit les étapes de la \cle{radicalisation}.
\item Indique le moment du \cle{passage à l'acte} violent.
\item Propose deux actions que la famille, l'école ou les autorités auraient pu mener avant le passage à l'acte.
\end{enumerate}
\end{exercice}

\begin{exercice}[Débat argumenté rédigé]
Sujet : « Les réseaux sociaux sont-ils devenus le premier espace de violence au Cameroun ? »
Rédige un texte argumenté de 15 à 20 lignes dans lequel tu :
\begin{enumerate}
\item présentes deux arguments en faveur de cette idée (exemples : cyberharcèlement, rumeurs, insultes en ligne) ;
\item présentes deux arguments contre (exemples : violences physiques dans les quartiers, conflits fonciers) ;
\item rédiges et justifies une conclusion personnelle claire.
\end{enumerate}
\end{exercice}

\begin{exercice}[Identifier victimes, auteurs et voies de recours]
À Bafoussam, M. Kamdem est roué de coups par deux individus qui lui prennent sa moto et son téléphone. Transporté à l'hôpital, il dépose plainte.
\begin{enumerate}
\item Identifie la ou les victimes et les auteurs des faits.
\item Qualifie les deux infractions commises (atteinte physique et atteinte au patrimoine).
\item Cite deux voies de recours ouvertes à M. Kamdem (plainte au commissariat ou à la gendarmerie, action en justice, demande de réparation).
\end{enumerate}
\end{exercice}

\courssub{Je me prépare au Bac}

\begin{exercice}[Étude de cas : agression et double vente à Douala]
À Douala, Mme Essomba achète un terrain à M. Bello, moyennant 2 500 000 FCFA, avec un reçu signé devant témoins. Deux mois plus tard, elle découvre que M. Bello a revendu le même terrain à M. Ngoa, qui a commencé à construire. En allant protester, Mme Essomba est insultée publiquement par M. Ngoa (« voleuse, menteuse ») devant les voisins, puis bousculée, ce qui lui cause une entorse au poignet constatée par un certificat médical.
\begin{enumerate}
\item Qualifie juridiquement chacun des faits commis : la double vente du terrain, les propos publics de M. Ngoa, la bousculade.
\item Pour chaque fait, précise s'il s'agit d'une atteinte à l'intégrité physique, morale ou au patrimoine.
\item Identifie les victimes et les auteurs de chaque infraction.
\item Indique les voies de recours dont dispose Mme Essomba pour obtenir réparation (plainte, action civile, rôle du certificat médical).
\item Dégage de ce cas une leçon de vie citoyenne sur le règlement pacifique des conflits fonciers.
\end{enumerate}
\end{exercice}

\begin{exercice}[Analyse de documents : l'extrémisme violent dans l'Extrême-Nord]
Document 1 (témoignage) : « Quand les hommes armés sont arrivés dans notre village de l'Extrême-Nord, ils ont brûlé les maisons, enlevé des jeunes et tué ceux qui résistaient. Nous avons fui à pied jusqu'à un site de déplacés. »

Document 2 (données) : la région de l'Extrême-Nord compte plusieurs centaines de milliers de personnes déplacées internes à cause des attaques attribuées au groupe Boko Haram ; de nombreuses écoles ont été fermées et des marchés abandonnés.
\begin{enumerate}
\item À partir du document 1, relève deux atteintes à l'intégrité physique et deux atteintes au patrimoine.
\item Définis l'\cle{extrémisme violent} et montre que les faits décrits en relèvent.
\item À partir des deux documents, dégage deux causes et trois conséquences (sociales, économiques, scolaires) de cette violence.
\item Propose trois actions concrètes de prévention et de lutte contre l'extrémisme violent (rôle de l'État, des chefs traditionnels et religieux, des jeunes).
\end{enumerate}
\end{exercice}

\begin{exercice}[Dissertation guidée type Probatoire]
Sujet : « Les violences sous toutes leurs formes menacent la cohésion sociale au Cameroun. » Montrez-le à partir d'exemples précis et proposez des solutions citoyennes.

Consignes : rédige un devoir complet comportant une introduction, un développement en deux parties et une conclusion.
\begin{enumerate}
\item Introduction : accroche, définition des termes clés du sujet (\cle{violence}, \cle{cohésion sociale}), annonce du plan.
\item Première partie : montre, avec au moins trois exemples précis (atteintes physiques, atteintes morales, atteintes au patrimoine, extrémisme violent), comment les violences fragilisent la vie en société au Cameroun.
\item Deuxième partie : propose au moins trois solutions citoyennes (éducation, dénonciation, dialogue, respect de la loi, rôle des médias et des réseaux sociaux).
\item Conclusion : bilan et ouverture sur le vivre-ensemble.
\end{enumerate}
\end{exercice}

\newpage

\coursec{Corrigés des exercices}

\begin{corrige}[Exercice 1 — QCM : reconnaître les formes de violences]
\begin{enumerate}
\item \textbf{Réponse b} : l'injure et la diffamation. Ce sont des atteintes à l'honneur et à la dignité, donc à l'intégrité \cle{morale}. Le vol de bétail (a) et la destruction d'un champ (c) sont des atteintes au \cle{patrimoine}.
\item \textbf{Réponse b} : la calomnie accuse une personne d'un fait en le sachant faux. Elle ajoute donc à la diffamation la conscience du mensonge. Elle peut être orale ou écrite et elle est punie par la loi.
\item \textbf{Réponse a} : le recel consiste à détenir, cacher ou vendre un bien dont on sait qu'il provient d'une infraction (vol, escroquerie). La tromperie pour obtenir de l'argent (b) est l'escroquerie ; la dégradation (c) est une destruction de bien.
\item \textbf{Réponse b} : l'extrémisme violent est le passage à la violence au nom d'une idéologie radicale. Une opinion radicale (a) relève de la liberté d'opinion tant qu'elle ne débouche pas sur la violence ; une grève (c) est un droit encadré par la loi.
\end{enumerate}
\end{corrige}

\begin{corrige}[Exercice 2 — Vrai ou faux, avec justification]
\begin{enumerate}
\item \textbf{Faux.} Frapper un camarade est une atteinte à l'intégrité \cle{physique} (coups et blessures), car le corps de la victime est atteint. L'atteinte morale vise l'honneur, la dignité ou la réputation.
\item \textbf{Vrai.} La double vente d'un même terrain est une atteinte au \cle{patrimoine} : elle porte sur un bien (le terrain) et prive l'un des acheteurs de son droit de propriété ; elle peut constituer une escroquerie.
\item \textbf{Vrai.} La \cle{radicalisation} est bien le processus progressif par lequel une personne adopte des idées extrémistes, se referme au dialogue et peut finir par justifier la violence.
\item \textbf{Faux.} La prévention des violences est l'affaire de tous : famille, école, chefs traditionnels et religieux, médias, associations de jeunes et chaque citoyen (dialogue, dénonciation, éducation). La police intervient surtout pour réprimer et protéger.
\end{enumerate}
\end{corrige}

\begin{corrige}[Exercice 3 — Définitions courtes]
\begin{itemize}
\item \textbf{\cle{Injure}} : parole, écrit ou geste outrageant adressé à une personne pour blesser son honneur ou sa dignité. C'est une atteinte à l'intégrité morale punie par la loi.
\item \textbf{\cle{Escroquerie}} : fait de tromper une personne (mensonges, faux documents, abus de confiance) pour obtenir d'elle de l'argent, un bien ou un avantage. C'est une atteinte au patrimoine.
\item \textbf{\cle{Viol}} : rapport sexuel imposé à une personne par la force, la menace ou la surprise. C'est l'une des plus graves atteintes à l'intégrité physique et morale, lourdement punie par la loi.
\item \textbf{\cle{Extrémisme violent}} : passage à la violence (attentats, enlèvements, destructions) au nom d'une idéologie radicale, politique ou pseudo-religieuse, pour imposer ses idées par la terreur.
\end{itemize}
\end{corrige}

\begin{corrige}[Exercice 4 — Associer la situation à l'atteinte]
\begin{center}
\begin{tabularx}{\linewidth}{|X|l|}
\hline
\rowcolor{popblueL} \textbf{Situation observée} & \textbf{Forme d'atteinte} \\
\hline
Un élève est blessé à coups de bâton derrière l'atelier & Intégrité \cle{physique} (coups et blessures) \\
\hline
Une rumeur mensongère circule sur WhatsApp contre une vendeuse & Intégrité \cle{morale} (diffamation) \\
\hline
Un hangar est incendié volontairement pendant la nuit & \cle{Patrimoine} (destruction volontaire de bien) \\
\hline
Des menaces de mort sont envoyées par SMS à un chef de quartier & Intégrité \cle{morale} (menaces), qui fait craindre une atteinte physique \\
\hline
\end{tabularx}
\end{center}
\textbf{Point de vigilance} : les menaces de mort ne blessent pas encore le corps ; elles atteignent la tranquillité et la dignité de la personne, d'où leur classement parmi les atteintes morales, tout en restant punies par la loi.
\end{corrige}

\begin{corrige}[Exercice 5 — Qualifier juridiquement des faits]
\begin{enumerate}
\item \textbf{Recel} (et vol en amont). L'ami détient puis revend un téléphone en sachant qu'il est volé (« Tu sais d'où ça vient ») : c'est la définition exacte du recel. Le voleur initial commet, lui, un vol.
\item \textbf{Diffamation}, voire \textbf{calomnie}. Le client porte publiquement (sur Facebook) atteinte à l'honneur du mécanicien en l'accusant d'un fait précis ; comme il sait que c'est faux, il y a calomnie, qui est une diffamation aggravée par la conscience du mensonge.
\item \textbf{Vol} (arraché de sac). L'inconnu s'approprie le bien d'autrui sans son consentement, avec violence surprise : c'est un vol, une atteinte au patrimoine.
\item \textbf{Injure publique}. Traiter quelqu'un de « bon à rien » devant tout le quartier est une parole outrageante qui blesse l'honneur : injure, atteinte à l'intégrité morale.
\end{enumerate}
\textbf{Point de vigilance} : bien distinguer injure (outrage sans fait précis), diffamation (allégation d'un fait précis portant atteinte à l'honneur) et calomnie (diffamation dont l'auteur sait le fait faux).
\end{corrige}

\begin{corrige}[Exercice 6 — Distinguer radicalisation et passage à l'acte]
\begin{enumerate}
\item \textbf{Étapes de la radicalisation} : (a) fréquentation d'un groupe qui prône la haine des autres communautés (embrigadement, endoctrinement) ; (b) refus de tout dialogue (fermeture de l'esprit, rupture avec l'entourage et les valeurs de tolérance).
\item \textbf{Passage à l'acte violent} : la participation à l'attaque contre le village. C'est le moment où les idées radicales se transforment en violence concrète : on parle alors d'\cle{extrémisme violent}.
\item \textbf{Actions préventives possibles} (au choix, deux parmi) :
\begin{itemize}
\item la famille et les chefs religieux ouvrent le dialogue, rappellent les valeurs de tolérance et de respect de la vie ;
\item l'école et les éducateurs détectent les signes de repli (isolement, discours haineux) et alertent les autorités compétentes ;
\item les autorités (administration, forces de l'ordre) surveillent et démantèlent les réseaux d'endoctrinement, proposent des programmes de réinsertion des jeunes ;
\item les associations de jeunes proposent des activités (sport, formation professionnelle) qui éloignent de l'embrigadement.
\end{itemize}
\end{enumerate}
\end{corrige}

\begin{corrige}[Exercice 7 — Débat argumenté rédigé]
\textbf{Corrigé type (à adapter avec tes propres exemples) :}

\emph{Introduction.} Avec la généralisation des téléphones portables et de l'internet au Cameroun, les réseaux sociaux (WhatsApp, Facebook, TikTok) occupent une place centrale dans la vie des jeunes. Certains affirment qu'ils sont devenus le premier espace de violence dans le pays. Cette idée est-elle fondée ?

\emph{Arguments en faveur.} D'abord, les réseaux sociaux sont le théâtre de nombreuses atteintes à l'intégrité morale : injures, cyberharcèlement entre élèves, diffusion de rumeurs mensongères qui détruisent la réputation d'un commerçant ou d'un élève en quelques heures. Ensuite, la rapidité et l'anonymat d'internet amplifient le mal : une fausse information partagée dans des groupes WhatsApp peut provoquer des conflits réels dans un quartier, et l'auteur se croit intouchable derrière son écran.

\emph{Arguments contre.} Cependant, les violences les plus graves restent physiques et hors ligne : agressions, coups et blessures dans les quartiers, violences faites aux femmes, attaques de l'extrémisme violent dans l'Extrême-Nord. De plus, les conflits fonciers (doubles ventes de terrains, spoliations) et les vols touchent directement le patrimoine des familles, sans aucun lien avec internet. Les réseaux sociaux ne sont donc qu'un espace de violence parmi d'autres.

\emph{Conclusion personnelle.} Les réseaux sociaux ne sont pas le premier espace de violence, mais ils sont devenus un amplificateur redoutable des atteintes morales. La solution n'est pas de les rejeter, car ils servent aussi à informer et à unir, mais d'éduquer les jeunes à un usage responsable : vérifier l'information avant de la partager, refuser l'injure, dénoncer les contenus haineux. Le citoyen responsable se comporte en ligne comme dans la rue : avec respect et dans le cadre de la loi.

\textbf{Barème indicatif} : introduction et problématique (2 pts) ; deux arguments « pour » étayés (4 pts) ; deux arguments « contre » étayés (4 pts) ; conclusion personnelle justifiée (2 pts) ; qualité de la langue et de l'organisation (2 pts). Total : 14 pts, ramené selon le barème de l'évaluation.

\textbf{Points de vigilance} : chaque argument doit être accompagné d'un exemple concret ; la conclusion doit trancher clairement et être justifiée, non simplement répéter le sujet.
\end{corrige}

\begin{corrige}[Exercice 8 — Identifier victimes, auteurs et voies de recours]
\begin{enumerate}
\item \textbf{Victime} : M. Kamdem, qui subit les coups et la perte de ses biens. \textbf{Auteurs} : les deux individus qui l'ont agressé et dépouillé.
\item \textbf{Infractions} : (a) les coups portés constituent des \cle{coups et blessures}, atteinte à l'intégrité \cle{physique} ; (b) l'enlèvement de la moto et du téléphone constitue un \cle{vol} avec violence, atteinte au \cle{patrimoine}.
\item \textbf{Voies de recours} (deux parmi) :
\begin{itemize}
\item déposer plainte au commissariat de police ou à la brigade de gendarmerie (ce qu'il a fait) ;
\item saisir la justice : action pénale pour faire punir les auteurs et action civile (constitution de partie civile) pour obtenir des dommages et intérêts ;
\item demander la restitution de la moto et du téléphone et la réparation du préjudice (frais d'hôpital, blessures), le certificat médical servant de preuve.
\end{itemize}
\end{enumerate}
\textbf{Point de vigilance} : ne pas confondre victime et témoin ; citer le certificat médical comme preuve de l'atteinte physique est un plus apprécié au Bac.
\end{corrige}

\begin{corrige}[Exercice 9 — Étude de cas : agression et double vente à Douala]
\begin{enumerate}
\item \textbf{Qualification des faits} :
\begin{itemize}
\item la double vente du même terrain est une \cle{escroquerie} (tromperie sur un bien déjà vendu) et une atteinte au droit de propriété ;
\item les propos publics de M. Ngoa (« voleuse, menteuse ») sont des \cle{injures publiques} et une \cle{diffamation}, car ils accusent Mme Essomba d'un fait (vol, mensonge) devant témoins ;
\item la bousculade ayant causé une entorse constitue des \cle{coups et blessures} (violences volontaires).
\end{itemize}
\item \textbf{Classification} : double vente $\rightarrow$ atteinte au \cle{patrimoine} ; injures et diffamation $\rightarrow$ atteinte à l'intégrité \cle{morale} ; bousculade et entorse $\rightarrow$ atteinte à l'intégrité \cle{physique}.
\item \textbf{Victimes et auteurs} :
\begin{itemize}
\item escroquerie : victimes Mme Essomba (et M. Ngoa, second acheteur trompé) ; auteur M. Bello ;
\item injures et diffamation : victime Mme Essomba ; auteur M. Ngoa ;
\item coups et blessures : victime Mme Essomba ; auteur M. Ngoa.
\end{itemize}
\item \textbf{Voies de recours} : Mme Essomba peut (a) déposer plainte au commissariat ou à la gendarmerie contre M. Bello (escroquerie) et contre M. Ngoa (injures, coups et blessures) ; (b) saisir le juge pour faire annuler la seconde vente et confirmer son droit sur le terrain (reçu signé devant témoins comme preuve de paiement) ; (c) se constituer partie civile pour obtenir des dommages et intérêts ; (d) utiliser le \cle{certificat médical} comme preuve de l'entorse pour l'action en coups et blessures.
\item \textbf{Leçon de vie citoyenne} : les conflits fonciers doivent se régler par le dialogue, la médiation des chefs et la voie judiciaire, jamais par l'insulte ou la violence. Avant tout achat, le citoyen prudent vérifie la propriété du terrain (titre foncier, enquête auprès des autorités) et exige des actes écrits : c'est le respect de la loi qui garantit la paix sociale.
\end{enumerate}
\textbf{Barème indicatif} : qualification des trois faits (6 pts) ; classification (3 pts) ; victimes et auteurs (3 pts) ; voies de recours (4 pts) ; leçon citoyenne (2 pts) ; expression (2 pts). Total : 20 pts.

\textbf{Points de vigilance} : exiger la justification de chaque qualification (éléments du texte cités) ; ne pas oublier que M. Ngoa est à la fois victime (de l'escroquerie) et auteur (injures, violences) ; citer le certificat médical et le reçu comme preuves.
\end{corrige}

\begin{corrige}[Exercice 10 — Analyse de documents : l'extrémisme violent dans l'Extrême-Nord]
\begin{enumerate}
\item \textbf{Atteintes relevées dans le document 1} :
\begin{itemize}
\item atteintes à l'intégrité physique : les \cle{enlèvements} de jeunes et les \cle{meurtres} de ceux qui résistaient ;
\item atteintes au patrimoine : l'\cle{incendie des maisons} et la fuite qui contraint les habitants à abandonner leurs biens, champs et marchés.
\end{itemize}
\item \textbf{Définition et application} : l'\cle{extrémisme violent} est le passage à la violence (tueries, destructions, terreur) au nom d'une idéologie radicale pour imposer ses vues. Les faits décrits en relèvent : un groupe armé (Boko Haram) attaque des villages, tue, enlève et détruit pour imposer sa doctrine par la terreur, en dehors de toute loi et de tout dialogue.
\item \textbf{Causes et conséquences} :
\begin{itemize}
\item causes (deux) : la pauvreté et le chômage des jeunes qui facilitent le recrutement ; l'endoctrinement idéologique et pseudo-religieux (radicalisation) ;
\item conséquences (trois) : sociale — des centaines de milliers de \cle{déplacés internes} et des familles brisées ; économique — marchés abandonnés, champs et biens détruits, appauvrissement de la région ; scolaire — fermeture de nombreuses écoles, enfants privés d'éducation.
\end{itemize}
\item \textbf{Trois actions de prévention et de lutte} :
\begin{itemize}
\item rôle de l'État : sécuriser les populations (forces de défense), rouvrir les écoles, développer la région (emplois, routes) et poursuivre les auteurs devant la justice ;
\item rôle des chefs traditionnels et religieux : prêcher la tolérance et la paix, détecter et dénoncer l'embrigadement, accueillir et réinsérer les jeunes repentis ;
\item rôle des jeunes : refuser les discours de haine, participer aux comités de vigilance, s'engager dans les activités éducatives, sportives et professionnelles.
\end{itemize}
\end{enumerate}
\textbf{Barème indicatif} : relevé des atteintes (4 pts) ; définition et application (4 pts) ; causes et conséquences (5 pts) ; actions de lutte (5 pts) ; expression (2 pts). Total : 20 pts.

\textbf{Points de vigilance} : appuyer chaque réponse sur les documents (citer les éléments du texte) ; distinguer clairement causes (avant) et conséquences (après) ; les actions proposées doivent être concrètes et attribuées à un acteur précis.
\end{corrige}

\begin{corrige}[Exercice 11 — Dissertation guidée type Probatoire]
\textbf{Corrigé type (plan détaillé rédigé) :}

\emph{Introduction.} Accroche : chaque jour, les faits divers rapportent agressions, rumeurs, vols ou attaques dans notre pays. Définitions : la \cle{violence} est tout acte qui porte atteinte à la personne (corps, honneur) ou à ses biens ; la \cle{cohésion sociale} est le lien de confiance et de solidarité qui permet aux membres d'une société de vivre ensemble en paix. Problématique : en quoi les violences fragilisent-elles le vivre-ensemble au Cameroun, et que peut faire le citoyen ? Annonce du plan : nous montrerons d'abord que les violences, sous toutes leurs formes, détruisent la vie en société, puis nous proposerons des solutions citoyennes.

\emph{Première partie : les violences fragilisent la vie en société.}
\begin{itemize}
\item Les atteintes \cle{physiques} (coups et blessures, agressions, viols) créent la peur et brisent la confiance entre voisins, élèves ou collègues d'atelier.
\item Les atteintes \cle{morales} (injures, diffamation, rumeurs sur les réseaux sociaux) détruisent l'honneur, enveniment les relations et provoquent des vengeances.
\item Les atteintes au \cle{patrimoine} (vols, escroqueries, doubles ventes de terrains, destructions) appauvrissent les familles et nourrissent les conflits, parfois violents, entre communautés.
\item L'\cle{extrémisme violent} (attaques dans l'Extrême-Nord) tue, déplace des populations entières, ferme les écoles et divise les communautés : c'est la négation même du vivre-ensemble.
\end{itemize}

\emph{Deuxième partie : des solutions citoyennes.}
\begin{itemize}
\item L'\cle{éducation} : famille et école transmettent le respect d'autrui, la tolérance et la culture de la paix dès le plus jeune âge.
\item La \cle{dénonciation} et le respect de la loi : signaler les violences aux autorités, refuser de se faire justice soi-même, saisir les tribunaux.
\item Le \cle{dialogue} et la médiation : régler les conflits (fonciers, de voisinage) par la discussion et l'arbitrage des chefs plutôt que par la force.
\item Le rôle des \cle{médias et des réseaux sociaux} : vérifier l'information, refuser les contenus haineux, promouvoir les messages de paix.
\end{itemize}

\emph{Conclusion.} Les violences, qu'elles touchent le corps, l'honneur ou les biens, minent la confiance sans laquelle aucune société ne peut vivre. Mais chaque citoyen, par son comportement quotidien, peut être un artisan de paix : c'est à ce prix que le Cameroun restera un pays où l'on vit ensemble, dans la diversité et le respect mutuel.

\textbf{Barème indicatif} : introduction complète (accroche, définitions, plan) (4 pts) ; première partie avec trois exemples précis et analysés (6 pts) ; deuxième partie avec trois solutions argumentées (6 pts) ; conclusion avec ouverture (2 pts) ; langue et présentation (2 pts). Total : 20 pts.

\textbf{Points de vigilance} : les exemples doivent être précis et camerounais (pas de généralités vagues) ; chaque solution doit être justifiée (pourquoi elle réduit la violence) ; respecter la structure annoncée dans l'introduction ; soigner l'orthographe et les connecteurs logiques (d'abord, ensuite, enfin, cependant).
\end{corrige}

\newpage

\coursec{Fiche bilan}

\coursec{Les formes de violences}

\begin{popfigure}
\begin{tikzpicture}[mindmap, grow cyclic, every node/.style={concept, text=white, font=\small},
  root concept/.append style={concept color=popink, font=\bfseries},
  level 1/.append style={sibling angle=60, level distance=3.4cm},
  level 1 concept/.append style={font=\footnotesize\bfseries},
  level 2/.append style={level distance=2.2cm},
  level 2 concept/.append style={concept color=popmuted, font=\scriptsize}]
\node[root concept]{Les formes de violences}
  child[concept color=popblue]{ node{Généralités}
    child{ node{Définition} }
    child{ node{Formes et victimes} }
  }
  child[concept color=poporange]{ node{Atteintes à l'intégrité physique}
    child{ node{Coups et blessures} }
    child{ node{Homicide, séquestration} }
  }
  child[concept color=poppurple]{ node{Atteintes à l'intégrité morale}
    child{ node{Diffamation, injure} }
    child{ node{Harcèlement, discrimination} }
  }
  child[concept color=popgreen]{ node{Atteintes au patrimoine}
    child{ node{Vol, escroquerie} }
    child{ node{Dégradation, abus de confiance} }
  }
  child[concept color=popgold]{ node{Extrémisme violent}
    child{ node{Idéologie radicale} }
    child{ node{Terrorisme, recrutement} }
  }
  child[concept color=poppink]{ node{Prévenir et combattre}
    child{ node{Dénoncer, alerter} }
    child{ node{Éducation, dialogue} }
  };
\end{tikzpicture}
\legende{Carte mentale de la leçon : les six grandes parties autour de la notion centrale de violence.}
\end{popfigure}

\begin{aretenir}[L'essentiel de la leçon]
\textbf{Définitions clés à connaître}
\begin{itemize}
  \item \cle{Violence} : tout acte ou comportement qui porte atteinte à l'intégrité physique, morale ou au patrimoine d'une personne.
  \item \cle{Atteinte à l'intégrité physique} : acte qui cause un dommage au corps d'autrui (coups et blessures, homicide, séquestration, violences corporelles).
  \item \cle{Atteinte à l'intégrité morale} : acte qui blesse l'honneur, la dignité ou la psychologie d'autrui (injure, diffamation, harcèlement, discrimination, humiliation).
  \item \cle{Atteinte au patrimoine} : acte qui porte préjudice aux biens d'autrui (vol, escroquerie, abus de confiance, dégradation, extorsion).
  \item \cle{Extrémisme violent} : adhésion à une idéologie radicale qui justifie l'usage de la violence, pouvant conduire au terrorisme.
\end{itemize}

\textbf{Repères à retenir par cœur}
\begin{center}
\begin{tabularx}{\linewidth}{|l|X|X|}\hline
\rowcolor{popblueL} \textbf{Forme de violence} & \textbf{Exemples} & \textbf{Conséquences} \\ \hline
Physique & coups, blessures, homicide & blessures, traumatismes, mort \\ \hline
Morale & injures, diffamation, harcèlement & souffrance psychique, isolement \\ \hline
Patrimoine & vol, escroquerie, dégradation & appauvrissement, insécurité \\ \hline
Extrémisme violent & radicalisation, terrorisme & insécurité, pertes en vies humaines \\ \hline
\end{tabularx}
\end{center}

\textbf{Méthodes express}
\begin{itemize}
  \item Face à une situation concrète : identifier la victime, qualifier l'atteinte (physique, morale ou patrimoine), citer la sanction encourue, proposer une solution citoyenne.
  \item Face à l'extrémisme violent : repérer les signes de radicalisation (isolement, discours haineux), alerter les autorités, promouvoir le dialogue et l'éducation.
  \item Rédaction au Probatoire/Bac : définir les mots-clés du sujet, illustrer par un exemple de la vie courante (atelier, quartier, réseaux sociaux), conclure par un engagement citoyen justifié.
\end{itemize}
\end{aretenir}

\begin{aretenir}[Checklist avant le Bac]
\begin{itemize}
  \item $\square$ Je sais définir la violence et citer ses trois grandes formes d'atteintes.
  \item $\square$ Je sais distinguer atteinte physique, atteinte morale et atteinte au patrimoine.
  \item $\square$ Je connais au moins deux exemples concrets pour chaque forme de violence.
  \item $\square$ Je sais définir l'extrémisme violent et expliquer ses causes.
  \item $\square$ Je connais les conséquences de l'extrémisme violent sur la société camerounaise.
  \item $\square$ Je sais citer les moyens de prévention des violences (éducation, dialogue, dénonciation).
  \item $\square$ Je connais le rôle du citoyen et des institutions dans la lutte contre les violences.
  \item $\square$ Je sais analyser une étude de cas : victime, atteinte, sanction, solution.
  \item $\square$ Je sais rédiger une conclusion engagée et justifiée en quelques phrases.
  \item $\square$ Je respecte les valeurs de paix, de tolérance et de vivre-ensemble dans mes propos.
\end{itemize}
\end{aretenir}

\findecours

\end{document}
